\documentclass[11pt]{article}
\usepackage[T1]{fontenc}
\usepackage[utf8]{inputenc}
\usepackage{amsmath,amsthm,mathtools,aliascnt}
\usepackage{newtxtext,newtxmath}
\usepackage[a4paper,margin=27mm,headheight=15pt]{geometry}
\usepackage{microtype}
\usepackage{enumitem,booktabs,longtable,array}
\usepackage{xcolor}
\usepackage{fancyhdr}
\usepackage{tikz}
\usetikzlibrary{arrows.meta,positioning}
\usepackage[most]{tcolorbox}
\usepackage{hyperref}
\usepackage[nameinlink,noabbrev]{cleveref}
\definecolor{ink}{HTML}{17324D}
\definecolor{soft}{HTML}{F2F5F8}
\hypersetup{colorlinks=true,linkcolor=ink,citecolor=ink,urlcolor=ink,
 pdftitle={Discrete Initial Groups and Omnific Normalization},
 pdfauthor={AI-assisted research draft prepared for the Surreal project},
 pdfsubject={Sign-tree surgery, a proposed solution of an Ehrlich--Kaplan question, convex factors and profinite obstructions}}
\pagestyle{fancy}
\fancyhf{}
\fancyhead[L]{\small\textsc{Discrete initial groups}}
\fancyhead[R]{\small\textsc{Omnific normalization}}
\fancyfoot[C]{\thepage}
\renewcommand{\headrulewidth}{0.3pt}
\setlength{\parskip}{3pt}
\setlength{\emergencystretch}{2em}
\setcounter{tocdepth}{2}
\setlist[enumerate]{leftmargin=*,itemsep=3pt,topsep=4pt}
\newtheorem{theorem}{Theorem}[section]
% Shared counter through aliascnt, so that cleveref names each type
% correctly (with a plain [theorem] counter every reference read "theorem").
\newcommand{\isgthm}[3]{%
 \newaliascnt{#1}{theorem}%
 \newtheorem{#1}[#1]{#2}%
 \aliascntresetthe{#1}%
 \crefname{#1}{#3}{#3s}%
 \Crefname{#1}{#2}{#2s}}
\isgthm{lemma}{Lemma}{lemma}
\isgthm{proposition}{Proposition}{proposition}
\newaliascnt{corollary}{theorem}
\newtheorem{corollary}[corollary]{Corollary}
\aliascntresetthe{corollary}
\crefname{corollary}{corollary}{corollaries}
\Crefname{corollary}{Corollary}{Corollaries}
\isgthm{fact}{Imported fact}{imported fact}
\theoremstyle{definition}
\isgthm{definition}{Definition}{definition}
\isgthm{example}{Example}{example}
\theoremstyle{remark}
\isgthm{remark}{Remark}{remark}
\isgthm{question}{Question}{question}
\newtheorem{maintheorem}{Theorem}
\renewcommand{\themaintheorem}{\Alph{maintheorem}}
\crefname{maintheorem}{Theorem}{Theorems}
\newcommand{\No}{\ensuremath{\mathbf{No}}}
\newcommand{\Oz}{\ensuremath{\mathbf{Oz}}}
\newcommand{\On}{\ensuremath{\mathbf{On}}}
\newcommand{\R}{\mathbb{R}}
\newcommand{\C}{\mathbb{C}}
\newcommand{\Z}{\mathbb{Z}}
\newcommand{\N}{\mathbb{N}}
\newcommand{\D}{\mathbb{D}}
\newcommand{\supp}{\operatorname{supp}}
\newcommand{\Pred}{\operatorname{Pred}}
\newcommand{\Rs}{\operatorname{R}_{s}}
\newcommand{\Ls}{\operatorname{L}_{s}}
\newcommand{\len}{\ell}
\newcommand{\cat}{\mathbin{\frown}}
\newcommand{\plus}{\langle+\rangle}
\newcommand{\minus}{\langle-\rangle}
\newcommand{\emp}{\varnothing}
\newcommand{\pr}{\mathrel{\prec}}
\newcommand{\pre}{\mathrel{\preccurlyeq}}
\newcommand{\Hahn}[2]{\mathcal{H}_{#1}(#2)}
\newcommand{\Rect}{\mathcal{N}_{\alpha,n}}
\newcommand{\tht}{\Theta_{\alpha}}
\newcommand{\psiA}{\Psi_{\alpha}}
\newcommand{\ct}{\operatorname{ct}}
\newcommand{\lc}{\operatorname{lc}}
\newcommand{\hdeg}{\operatorname{deg}_{H}}
\newcommand{\lex}{\mathbin{\overrightarrow{\times}}}
% Notation of the second source (convex factors); see isg:cf:sec:conventions.
\newcommand{\Q}{\mathbb{Q}}
\newcommand{\Zhat}{\widehat{\Z}}
\newcommand{\IR}{\mathsf{Init}}
\newcommand{\otp}{\operatorname{otp}}
\newcommand{\res}{\operatorname{res}}
\newcommand{\Ext}{\operatorname{Ext}}
\newcommand{\dv}{\mathrm{div}}
\newcommand{\cmp}[1]{\flat_{#1}}
\newcommand{\source}[1]{\textup{\cite{#1}}}
\newcommand{\replabel}[1]{\nolinkurl{#1}}
\newcommand{\fn}[1]{\nolinkurl{#1}}
\newtcolorbox{statusbox}{colback=soft,colframe=ink,boxrule=0.5pt,
 arc=1mm,left=3mm,right=3mm,top=2mm,bottom=2mm}

\begin{document}
\hypersetup{pageanchor=false}
\begin{titlepage}
\vspace*{2mm}
{\color{ink}\Large\textsc{Surreal numbers / ordered algebra}}\par
\vspace{4mm}
{\color{ink}\huge\bfseries Discrete Initial Groups\\[3mm]
and Omnific Normalization\par}
\vspace{4mm}
{\Large Sign-tree surgery, a sharp image classification,\\
an Ehrlich--Kaplan question, and convex factors\par}
\vspace{5mm}
{\large AI-assisted research draft prepared for the Surreal project\par}
\vspace{2mm}
23 September 2026 \textperiodcentered{} two-source report of the
collection
\vspace{4mm}
\begin{abstract}
We give an explicit construction proposed as an affirmative solution of
Ehrlich and Kaplan's question whether every discretely ordered initial
subgroup of the surreal numbers is order-isomorphic to an initial subgroup
of the omnific integers (Question~2 of the arXiv preprint, Question~9.1 of
the journal version). The construction is not ordinary multiplication
by the reciprocal of the least positive element. Instead, a sign-sequence
operation moves the minimum exponent to zero, while preserving all
right-ancestor relations away from that minimum; a separate rescaling
changes only the bottom coefficient group to the ordinary integers.
The Ehrlich--Kaplan characterization of initial Hahn subgroups then
establishes initiality of the image.

The argument yields more than existence. For a prescribed least positive
element $2^{-n}\omega^{-\alpha}$, we characterize the images exactly by
containment of an explicit dyadic-spine subgroup $K_\alpha\subseteq\Oz$.
We identify the smallest and largest initial groups with that least
positive element, give sharp cardinality consequences, construct an
initial realization of the quotient by the least positive cyclic subgroup,
and extend the additive normalization to rectangular surcomplex modules.
All normal-form transports preserve leading truncations and set-indexed
strong summability. Counterexamples explain why scalar normalization,
exponent translation, and unrestricted inversion do not suffice.

A second manuscript realizes every convex subquotient of an initially
realizable group by sign-tree compression, and shows that a Presburger
group is initially realizable exactly when its residues in
$\widehat{\Z}$ are ordinary integers.
\end{abstract}
\vspace{2mm}
\begin{statusbox}
\textbf{Status and scope.} This is an unrefereed research manuscript with
full proposed proofs, not an independently certified resolution or a
claim of established priority. The principal external mathematical input
is explicitly identified. Finite regression checks accompany the text;
they do not certify the transfinite argument. No new Lean formalization
is claimed. A targeted literature and repository review found the cited
question and its motivating example, but did not establish exhaustive
novelty. At placement in the collection every step was checked by hand
and no gap was found (\cref{isg:sec:provenance}); that check is not a
refereeing, and the literature search was targeted, not exhaustive.
The second manuscript (\cref{isg:cf:sec:intro}) has the same status.
\end{statusbox}
\vfill
{\small\textbf{Keywords:} surreal numbers; omnific integers; initial
subgroups; sign sequences; Hahn series; ordered abelian groups.\par
\textbf{Mathematical focus:} additive and simplicity-theoretic structure,
not multiplicative or exponential equivalence.}
\end{titlepage}
\hypersetup{pageanchor=true}

\begingroup
\small
\setlength{\parskip}{0pt}
\tableofcontents
\endgroup
\newpage

\section{The problem and the main results}\label{isg:sec:intro}

\subsection{A specific published question}
An additive subgroup of $\No$ is \emph{initial} when it contains every
proper sign-sequence prefix of each of its members. This is a much stronger
requirement than being an ordered subgroup. An initial subgroup is
\emph{discrete} when it is nonzero and has a least positive element.
Throughout, ``isomorphism'' of ordered groups preserves addition and
order; it need not preserve the ambient simplicity relation.

Ehrlich and Kaplan ask:
\begin{quote}
``Is every discrete initial subgroup of $\No$ isomorphic to an initial
subgroup of $\Oz$?''
\end{quote}
This is Question~2 in Section~9 of the explicitly versioned preprint
\cite{EK}, printed page~18. The journal article appeared in 2018
\cite{EKjournal}. We use the preprint's numbering throughout to avoid
mixing versions. The question also appears in Kaplan's thesis
\cite{Kaplan}.

In the journal version the same question is Question~9.1, with the same
wording and the same motivating example. The results of \cite{EK} used
below are numbered differently there; the correspondence, checked when
this report was placed in the collection against the author-hosted text
of the journal version \cite{EKjournal}, is:
\begin{center}\small
\begin{tabular}{@{}lll@{}}
\toprule
arXiv v1 \cite{EK} & journal version \cite{EKjournal} & use here\\
\midrule
Lemma~3 & Lemma~5.1 & initial subgroups of $\R$, \cref{isg:fact:real}\\
Theorem~1 & Theorem~5.1 & the initiality criterion, \cref{isg:fact:EK}\\
Proposition~9 & Proposition~5.2 & shape of the least positive element\\
Question~2 & Question~9.1 & the question addressed here\\
Question~3 & Question~9.2 & the initial-monoid problem, not addressed\\
Theorem~6 & Theorem~8.1 & optimality with class models (source~02)\\
Question~1 & Question~8.1 & set-model optimality, not answered\\
\bottomrule
\end{tabular}
\end{center}
Below, the arXiv numbers are used, except where the wording of the
journal version matters (after \cref{isg:fact:EK}).

The distinction between a group and a ring matters. A discrete initial
subgroup can contain $1/2$ or a nonzero infinitesimal, whereas an omnific
integer has neither a nonintegral constant coefficient nor a negative
normal-form exponent. The task is therefore to change the initial
\emph{realization} of the ordered group, not to prove that its original
realization already lies in $\Oz$.

The paper addresses precisely this existence question and several
structural refinements. It does not claim a classification of all ordered
abelian groups admitting initial embeddings, nor a solution of the
separate initial-monoid problem (Question~3 of \cite{EK}, Question~9.2
of the journal version).

\subsection{Normalization and its exact range}
Write $\N=\{0,1,2,\ldots\}$ and $\D=\Z[1/2]$. If $\alpha$ is an ordinal, let
\[
 p_\alpha=-\alpha=\minus^{\alpha},
 \qquad q_\beta=\plus\cat\minus^{\beta}\quad(\beta\in\On).
\]
Here $\cat$ means concatenation of sign sequences, not surreal addition.
For example, $q_0=1$, $q_1=1/2$, and $q_m=2^{-m}$ for finite $m$.
Define the finite-support additive subgroup
\begin{equation}\label{isg:eq:Kintro}
 K_\alpha=\Z\oplus\bigoplus_{\beta<\alpha}\D\,\omega^{q_\beta}
 \ \subseteq\Oz.
\end{equation}
The direct sum means that each element has only finitely many nonzero
summands. In particular, $K_0=\Z$.

\begin{maintheorem}[Omnific normalization]\label{isg:thm:A}
Let $A\subseteq\No$ be a nonzero discrete initial additive subgroup.
There are unique $\alpha\in\On$ and $n\in\N$ such that its least positive
element is
\[
 \varepsilon=2^{-n}\omega^{-\alpha}.
\]
An explicitly defined map $\Rect$ restricts to an order-preserving group
isomorphism
\[
 \Rect:A\longrightarrow H,
 \qquad H\subseteq\Oz\text{ initial},\qquad \Rect(\varepsilon)=1.
\]
The map extends to a real-linear order-isomorphism between the full
normal-form spaces supported in $[-\alpha,\infty)$ and $[0,\infty)$.
This extension preserves support order types, all leading truncations,
and set-indexed strong summability in both directions.
\end{maintheorem}

\begin{maintheorem}[Sharp image classification]\label{isg:thm:B}
Fix $\alpha\in\On$ and $n\in\N$. Under the same explicit map $\Rect$,
the initial groups with least positive element
$2^{-n}\omega^{-\alpha}$ correspond bijectively, with inclusions
preserved and reflected, to the initial groups $H\subseteq\Oz$ satisfying
\[
 K_\alpha\subseteq H.
\]
The inverse is explicit. Thus the dyadic-spine condition is both necessary
and sufficient, not merely an artifact of a sufficient construction.
\end{maintheorem}

\begin{maintheorem}[Removing the bottom cyclic layer]\label{isg:thm:C}
For every group $A$ in \cref{isg:thm:A}, the subgroup $\Z\varepsilon$ is
convex, and the ordered quotient $A/\Z\varepsilon$ has an explicit
initial realization in $\No$. The corresponding short exact sequence
splits as an additive group, with a canonical splitting relative to the
given Conway normal forms. Equivalently,
\[
 A\cong B\lex\Z
\]
for an explicitly constructed initial subgroup $B\subseteq\No$, where
the first coordinate is dominant in the lexicographic order.
\end{maintheorem}

The main construction is in \cref{isg:sec:surgery,isg:sec:transport}; the proofs
of the three theorems are completed in
\cref{isg:sec:normalization,isg:sec:classification,isg:sec:quotient}, respectively.
Theorem~B contains more information than Theorem~A: it identifies exactly
which normalized realizations can be moved to any prescribed least
positive scale by this construction.

\subsection{Why a change of exponents is necessary}
Multiplication by $\varepsilon^{-1}$ certainly gives an ordered-group
isomorphism. Initiality, however, refers to sign prefixes rather than
numerical intervals, and multiplication does not generally preserve it.
Even $A=\frac12\Z+\Z\omega$ is an initial group for which
$2A=\Z+2\Z\omega$ is not initial: $2\omega$ belongs to the image but
its simpler predecessor $\omega$ does not.

Our operation treats two different issues separately. First, it relocates
the minimum of the \emph{exponent tree} to zero. Second, it rescales only
the real coefficient at that minimum. No claim of multiplicativity is
made, and none is needed for the published ordered-group question.

\subsection{Conventions shared with the collection}\label{isg:sec:conventions}
This report is built from two manuscripts: source~01, the normalization
manuscript (\cref{isg:sec:intro,isg:sec:background,isg:sec:bottom,isg:sec:surgery,isg:sec:transport,isg:sec:normalization,isg:sec:classification,isg:sec:quotient,isg:sec:complex,isg:sec:examples}
and Appendix~\ref{isg:app:compact}), and source~02, the convex-factor manuscript
(\cref{isg:cf:sec:intro}, \cref{isg:cf:sec:cuts,isg:cf:sec:compression,isg:cf:sec:factors,isg:cf:sec:presburger,isg:cf:sec:factorial,isg:cf:sec:modeltheory,isg:cf:sec:complex,isg:cf:sec:boundaries}
and Appendix~\ref{isg:cf:app:core}). Their provenance, the placement reviews and
the relation to other reports are in \cref{isg:sec:provenance}. Statements
added in the merge of the two sources, and proved here, are marked
\textup{[merge]}. The following conventions are fixed here once; the
notation of source~02 is fixed in \cref{isg:cf:sec:conventions}.
\begin{enumerate}
\item \emph{Initial} always means closed under proper sign-sequence
      prefixes, as in the definition of the definable-surreals report
      (\replabel{dsn:sec:foundations}). It never means an initial segment
      of the numerical order, and it has nothing to do with the
      \emph{initial forms} (leading forms of polynomial solutions) of the
      Diophantine report (\replabel{odg:prop:initial}).
\item $\Oz=\Z\oplus\Pi$, with $\Pi$ the class of purely infinite normal
      forms, is the omnific ring of the collection
      (\replabel{odg:def:rings}, \replabel{osq:eq:Ozdef},
      \replabel{found:eq:omnific}); \eqref{isg:eq:Oz} below is the same
      description. For $x\in\Oz$, $\ct(x)$ is its coefficient at
      exponent~$0$, as there. Only the additive and order structure of
      $\Oz$ is used; its ring structure plays no role.
\item The symbols $\Theta_\alpha$, $\Psi_\alpha$, $q_\beta$ and
      $\mathcal N_{\alpha,n}$ are local to this report. They are unrelated
      to the theta series and to the nome or dilation parameter $q$ of
      other reports.
\end{enumerate}

\subsection{Convex factors and profinite obstructions}\label{isg:cf:sec:intro}
\Cref{isg:thm:C} removes the least positive cyclic subgroup. The second
source of this report, \emph{Convex Factors and Profinite Obstructions for
Initial Surreal Groups}, removes an arbitrary convex subgroup. Call an
ordered abelian group \emph{initially realizable} if it is order-isomorphic
to an additive subgroup of $\No$ closed under all proper sign-sequence
prefixes, and write $G\in\IR$ for this property; the zero group is
included. A \emph{convex factor} of $G$ is $C'/C$, where $C\subseteq C'$
are convex subgroups of $G$ and the quotient has its usual ordered group
structure.

\begin{maintheorem}[Convex-factor closure]\label{isg:cf:main:factor}
If $G\in\IR$ and $C\subseteq C'$ are convex subgroups of $G$, then
$C'/C\in\IR$. For an actual initial realization $G\subseteq\No$, the
construction is explicit: restrict normal forms to the corresponding
convex exponent interval and compress its sign tree.

Moreover, every convex subgroup sequence
\[
 0\longrightarrow C\longrightarrow G\longrightarrow G/C
 \longrightarrow0
\]
splits as a sequence of abelian groups. In the chosen normal-form
realization the section is order-preserving and the decomposition is
lexicographic. The coordinate projections and compressed factor
realizations satisfy precise composition laws under nested restrictions.
\end{maintheorem}

The theorem concerns \emph{initial realizability}, not literal initiality
of the original subgroup $C$ or of a selected complement. It also does
not assert a splitting invariant under every abstract automorphism.
The explicit splitting is attached to a chosen initial realization.
With $C'=G$ it answers \cref{isg:q:convex}: every convex subgroup of an
initial group has an explicit initial realization of the quotient by a
normal-form transport, and for $C=\Z\varepsilon$ that transport is the
construction of \cref{isg:sec:quotient}
(\cref{isg:cf:prop:cyclic}). It also answers \cref{isg:q:limit}: no
extra hypothesis is needed at a limit stage (\cref{isg:cf:cor:limit}).

A \emph{$\Z$-group} is a discretely ordered abelian group $G$ with least
positive element $e$ such that $G/\Z e$ is divisible. These are exactly
the ordered-group models of additive Presburger arithmetic
\cite{Jerabek}. Their coherent standard residues define
\[
 \res_G:G\longrightarrow\Zhat=\varprojlim_{n\ge1}\Z/n\Z,
 \qquad \res_G(e)=1.
\]
Here $\Z$ on the right is the ordinary, diagonally embedded copy of the
integers, not a nonstandard integer substructure.

\begin{maintheorem}[Presburger initiality criterion]\label{isg:cf:main:presburger}
For a set-sized $\Z$-group $G$, the following are equivalent:
\[
\begin{gathered}
 G\in\IR;\qquad
 G\text{ has an initial realization in }\Oz;\qquad
 \Z e\text{ is a direct summand of }G;\\
 \res_G(G)=\Z;\qquad
 G\cong G^{\dv}\lex\Z,\quad G^{\dv}=\bigcap_{n\ge1}nG.
\end{gathered}
\]
The group $G^{\dv}$ is divisible. In the split case its complement to
$\Z e$ is unique.
\end{maintheorem}

This is a global compatibility condition, not a condition on any one
finite modulus. Every $\Z$-group has the same individual quotients
$G/nG\cong\Z/n\Z$. Initiality distinguishes whether the residues of a
\emph{single element}, across all standard moduli, come from one ordinary
integer.

\begin{maintheorem}[Explicit countermodels]\label{isg:cf:main:models}
There are $2^{\aleph_0}$ pairwise nonisomorphic countable rigid
$\Z$-groups which are not initially realizable, although they are all
elementarily equivalent to $\Z$. There are also countable densely
ordered groups, elementarily equivalent to an explicit initial subgroup
of $\No$, which are not initially realizable.

No recursively saturated $\Z$-group, and no nonprincipal ultrapower over
$\N$ of a $\Z$-group, is initially realizable. The ordered additive group
completion of any nonstandard model of Peano arithmetic is likewise not
initially realizable.
\end{maintheorem}

The proofs are in \cref{isg:cf:sec:factors} (\cref{isg:cf:main:factor}),
\cref{isg:cf:sec:presburger} (\cref{isg:cf:main:presburger}) and
\cref{isg:cf:sec:modeltheory} (\cref{isg:cf:main:models}). In source~02
these three theorems are Theorems~A, B and~C; they are lettered D, E
and~F here because source~01's Theorems~A--C keep their letters.

\paragraph{What is new in source~02, and what is not.} The primary external
structural input is again the Ehrlich--Kaplan characterization of initial
Hahn subgroups (\cref{isg:fact:EK}). Their treatment of convex subgroups
asks which subgroups in the \emph{given} realization are themselves
initial. The construction of source~02 changes the realization of a
factor and preserves its ordered additive structure. Source~01 already
gives the normalization of \cref{isg:thm:A}, the integer-bottom suspension
and the relative splitting criterion for discrete groups
(\cref{isg:thm:suspension,isg:cor:relative}); its \cref{isg:sec:quotient}
explicitly leaves arbitrary convex quotients outside its conclusions.
Source~02 does not claim those results again as new. It replaces the
special positive-root deletion by a compression operation valid for
arbitrary convex exponent intervals; in particular it handles cuts with no
endpoint and convex factors not accessible by a finite iteration of
least-positive quotients.

The profinite residue map and its divisible kernel belong to the
classical structure theory of $\Z$-groups; see especially
\cite[Sections~2 and~4]{Jerabek}. Likewise, the extension parameter
$\Ext(\Q,\Z)\cong\Zhat/\Z$ is classical \cite[Example~8.27]{Strickland}.
Source~02's arithmetic contribution is their explicit connection with
initial surreal realizability, together with the displayed constructions
and obstructions. The existence of rigid Presburger models is not claimed
as a new phenomenon. The results are not announced as a resolution of the
set-model version of Ehrlich--Kaplan's optimality question (Question~1 of
\cite{EK}, Question~8.1 of the journal version); a precise comparison is
in \cref{isg:cf:sec:optimality}. A targeted review found no matching
formulation of the main convex-factor theorem, but this is not an
exhaustive priority search.

\section{Foundations and the imported initiality criterion}\label{isg:sec:background}

\subsection{Sets, classes, signs, and two orders}
We work in NBG with global choice, in the conventional class-sized
formulation of surreal numbers. Every individual surreal is a
set-length sign sequence. Every individual normal form has set support.
All strongly summable families considered below are indexed by sets.
Statements about families of class-sized groups are understood as
correspondence statements about class parameters, not as assertions that
a class of all classes exists.

For sign sequences $u,v$, write $u\pr v$ for a proper prefix and
$u\pre v$ for a prefix allowing equality. Their ordinary numerical order
is determined at the first disagreement, with
\[
 -\ <\ \text{undefined}\ <\ +.
\]
The empty sequence is the number $0$. Put
\[
 \Pred(x)=\{y:y\pr x\},\qquad
 \Rs(x)=\{y\pr x:x<y\},\qquad
 \Ls(x)=\{y\pr x:y<x\}.
\]
In sign language, a prefix obtained by stopping just before a minus sign
is a right ancestor; stopping just before a plus sign gives a left
ancestor. Each of these ancestor collections is a set.

A subclass $E\subseteq\No$ is initial precisely when
$\Pred(x)\subseteq E$ for every $x\in E$. This does not mean that $E$
is downward closed in the numerical order. For instance, $\No_{\ge0}$
is initial as a sign tree but is a numerical final segment.

If $u$ and $v$ have lengths $\lambda$ and $\mu$, then
$u\cat v$ has length $\lambda+_{\mathrm{ord}}\mu$. All additions of
\emph{lengths} in this paper are ordinal additions. Concatenation is
associative and preserves comparison between tails sharing the same
prefix. At a limit ordinal there is no last sign to remove: ``remove the
first sign'' or ``remove a specified concatenated block'' always means
restricting to the appropriate interval and taking its order type.
These conventions agree with the sign-sequence framework of
\cite{BagHoeven}.

\subsection{Hahn spaces without an exponent group}
For an ordered subclass $E\subseteq\No$, denote by $\Hahn{\R}{E}$
the additive space of formal expressions
\begin{equation}\label{isg:eq:nf}
 x=\sum_{\gamma\in S} r_\gamma\omega^\gamma,
 \qquad r_\gamma\in\R\setminus\{0\},
\end{equation}
where $S\subseteq E$ is a set reverse-well-ordered by $<$.
Conway normal forms identify these expressions with the corresponding
surreal numbers. The ordering is the sign of the coefficient at the
greatest exponent of a nonzero difference. For $x\ne0$, write
\[
 \hdeg(x)=\max\supp(x),\qquad
 \lc(x)=r_{\hdeg(x)}.
\]
We do not assert that $E$ is an additive group. Accordingly,
$\Hahn{\R}{E}$ is used here as an additive, real-linear Hahn space, not
necessarily as a ring. This distinction permits nonadditive exponent
reindexings.

Suppose $S=\{\gamma_\xi:\xi<\lambda\}$ is enumerated in strictly
decreasing order. A leading truncation of $x$ is
\[
 x_{<\eta}=\sum_{\xi<\eta}r_{\gamma_\xi}\omega^{\gamma_\xi}
 \quad(\eta\le\lambda).
\]
``Truncation closed'' will always refer to these leading truncations,
including those at limit indices, not to closure under arbitrary
subseries. The normal-form and monomial facts used here are part of the
standard theory summarized in \cite{EK}; we never choose a noncanonical
Hahn embedding of an arbitrary ordered group.

\begin{definition}[Strong summability]\label{isg:def:strong}
A set-indexed family $(x_i)_{i\in I}$ in a Hahn space is strongly
summable if $\bigcup_i\supp(x_i)$ is reverse-well-ordered and each
exponent occurs in only finitely many supports. Its sum is the normal
form obtained by the finite coefficient sum at each exponent, with zero
coefficients discarded.
\end{definition}

For a subgroup $G\subseteq\Hahn{\R}{E}$, set
\[
 G_\gamma=\{r\in\R:r\omega^\gamma\in G\}.
\]
It is \emph{cross-sectional on $E$} when $\omega^\gamma\in G$ for
every $\gamma\in E$.

\begin{fact}[Ehrlich--Kaplan, concrete form]\label{isg:fact:EK}
An initial subgroup of $\No$ is a truncation-closed, cross-sectional
Hahn subgroup on its initial exponent class
$\Gamma(G)=\{\gamma:\omega^\gamma\in G\}$. Its nonzero coefficient
groups are initial subgroups of $\R$, and
\begin{equation}\label{isg:eq:dyadic}
 \gamma,\eta\in\Gamma(G),\quad \eta\in\Rs(\gamma)
 \quad\Longrightarrow\quad \D\subseteq G_\eta.
\end{equation}
Conversely, these conditions imply that the canonical normal-form image
is initial in $\No$.
\end{fact}

This is the form of \cite[Theorem~1 and its proof]{EK} needed here.
The last sentence is important: the sufficiency proof establishes
initiality of the \emph{specified canonical image}, not just existence of
some other initial embedding.

The journal version states this concrete form outright. Its Theorem~5.1,
the counterpart of \cite[Theorem~1]{EK}, says that a subgroup of $\No$ is
initial if and only if it is isomorphic ``via the canonical isomorphism''
to a truncation-closed, cross-sectional Hahn subgroup whose exponent
class is an initial ordered subclass of $\No$ and whose coefficient
groups satisfy the two conditions above \cite{EKjournal}. The canonical
isomorphism is the normal-form identification of $\No$ with the Hahn
group, so the journal statement is \cref{isg:fact:EK} as used here.
The criterion is imported as a published statement. Neither the source
manuscript nor the placement review re-proves it
(\cref{isg:sec:provenance}).

\begin{fact}[Real coefficient groups]\label{isg:fact:real}
A nonzero initial additive subgroup of $\R$ either contains $\D$ or is
$2^{-m}\Z$ for an integer $m\ge0$ \cite[Lemma~3]{EK}
\textup{(}Lemma~5.1 of the journal version\textup{)}.
In particular, a discrete one has least positive element $2^{-m}$.
\end{fact}

No proof in the present paper replaces \cref{isg:fact:EK} by the weaker
assertion that initiality of the exponent set alone is sufficient. The
coefficient groups and the right-ancestor condition are essential.

\subsection{Omnific integers and coefficient isolation}
We use the normal-form description
\begin{equation}\label{isg:eq:Oz}
 \Oz=\left\{\sum_{\gamma>0}r_\gamma\omega^\gamma+m:
           m\in\Z\right\},
\end{equation}
with set reverse-well-ordered positive support and arbitrary real
coefficients. This is Conway's canonical integer part, as recalled in
\cite[Section~6.2]{Kaplan}, and the collection's $\Oz=\Z\oplus\Pi$
(\cref{isg:sec:conventions}). It is an initial additive subgroup: apply
\cref{isg:fact:EK} to $E=\No_{\ge0}$, coefficient group $\Z$ at zero and
$\R$ at every positive exponent. Zero is never a right ancestor of a
nonnegative number, so the exceptional coefficient group causes no
violation of \eqref{isg:eq:dyadic}. Every positive element of \Oz is at
least $1$; this discreteness is also \replabel{odg:prop:units} of the
Diophantine report.

\begin{lemma}[Isolating a normal-form term]\label{isg:lem:isolate}
Let $G$ be an additive subgroup closed under leading truncations. If
$x\in G$ has a nonzero coefficient $r_\gamma$ at exponent $\gamma$,
then $r_\gamma\omega^\gamma\in G$.
\end{lemma}
\begin{proof}
If $\gamma=\gamma_\eta$ in the decreasing enumeration of the support,
subtract the truncation at $\eta$ from the truncation at $\eta+1$.
Both belong to $G$. This also works when $\eta$ is a limit ordinal.
\end{proof}

This lemma permits removing a specified individual term. It does not
license an arbitrary infinite subseries of an element of $G$.

\section{Locating the least positive scale}\label{isg:sec:bottom}

\begin{lemma}[A minimum in an initial tree]\label{isg:lem:min-tree}
If a nonempty initial subclass $E\subseteq\No$ has a numerical minimum
$p$, then $p=-\alpha$ for a unique ordinal $\alpha$.
Moreover, the cone $C_\alpha=\{x\in\No:x\ge-\alpha\}$ is initial.
\end{lemma}
\begin{proof}
A plus sign in $p$ would give a proper prefix smaller than $p$.
Initiality would place that prefix in $E$, contradicting minimality.
Thus every sign in $p$ is minus.

For the second assertion, let $p=-\alpha$, $x\ge p$, and $v\pr x$.
If $v<p$, then at the first disagreement with the all-minus sequence $p$
one would have to find either a sign smaller than minus, which is
impossible, or a minus sign after $p$ has ended. The latter would force
$x$ itself to be below $p$. Hence $v\ge p$.
\end{proof}

\begin{proposition}[Bottom coefficient structure]\label{isg:prop:bottom}
Let $A$ be a nonzero discrete initial subgroup with least positive
member $\varepsilon$. There are unique $\alpha\in\On$, $n\in\N$
such that, writing $p=-\alpha$ and $c=2^{-n}$,
\begin{equation}\label{isg:eq:bottomdata}
 \varepsilon=c\omega^p,\qquad
 \min\Gamma(A)=p,\qquad A_p=c\Z.
\end{equation}
Every support of every member of $A$ lies in $C_\alpha$.
\end{proposition}
\begin{proof}
By \cref{isg:fact:EK}, $A$ is truncation closed. Write the leading term of
$\varepsilon$ as $r\omega^p$ with $r>0$. This term belongs to $A$.
If the remaining tail $t$ were nonzero, then $|t|\in A$ and
$0<|t|<\varepsilon$, since the leading exponent of $t$ is strictly
smaller than $p$. This contradicts the choice of $\varepsilon$.
Consequently $\varepsilon=r\omega^p$.

The exponent $p$ lies in $\Gamma(A)$ by the cross-sectional normal-form
description. If $\gamma<p$ belonged to $\Gamma(A)$, then
$0<\omega^\gamma<r\omega^p=\varepsilon$, another contradiction.
Thus $p$ is the minimum of $\Gamma(A)$. By \cref{isg:lem:min-tree},
$p=-\alpha$ for a unique ordinal $\alpha$.

The real coefficient group $A_p$ has least positive element $r$.
By \cref{isg:fact:real}, it is $2^{-n}\Z$ for a unique $n\in\N$, and
$r=2^{-n}$. All supports lie in $\Gamma(A)$, proving the last assertion.
\end{proof}

The shape of the least positive element is consistent with
\cite[Proposition~9 and its proof]{EK} (Proposition~5.2 of the journal
version); the argument above also records
the support and coefficient data needed later.

\begin{corollary}[Forced ancestors]\label{isg:cor:forced}
Under \eqref{isg:eq:bottomdata}, for every $\beta<\alpha$ one has
\[
 -\beta\in\Gamma(A),\qquad \D\omega^{-\beta}\subseteq A.
\]
\end{corollary}
\begin{proof}
The number $-\beta$ is a right ancestor of $-\alpha$.
Initiality gives the exponent, and \eqref{isg:eq:dyadic} gives all dyadic
coefficients there.
\end{proof}

The bottom coefficient group can be discrete, but all these ancestors
must carry dyadically divisible coefficient groups. This is the information
that survives normalization as the subgroup $K_\alpha$.

\section{Moving a minimum to zero in the sign tree}\label{isg:sec:surgery}

\subsection{Definition and inverse}
Fix $\alpha\in\On$, put $p=-\alpha$, and keep
$C_\alpha=[p,\infty)$. Every $x\in C_\alpha$ belongs to exactly one of
the following cases: $x=p$; $x=p\cat\plus\cat u$ for a unique tail
$u$; or $p$ is not a prefix of $x$.
Indeed, an extension through $p\cat\minus$ would lie below the cone.

\begin{definition}[Root relocation]\label{isg:def:theta}
Define $\tht:C_\alpha\to\No_{\ge0}$ by
\begin{equation}\label{isg:eq:theta}
 \tht(x)=
 \begin{cases}
  0, & x=p,\\
  \plus\cat p\cat u, & x=p\cat\plus\cat u,\\
  \plus\cat x, & p\not\pre x.
 \end{cases}
\end{equation}
The inverse candidate $\psiA:\No_{\ge0}\to C_\alpha$ is
\begin{equation}\label{isg:eq:psi}
 \psiA(0)=p,\qquad
 \psiA(\plus\cat v)=
 \begin{cases}
  p\cat\plus\cat u,&v=p\cat u,\\
  v,&p\not\pre v.
 \end{cases}
\end{equation}
\end{definition}

The symbols in these formulas are sign blocks. For $\alpha=0$, the block
$p$ is empty and both maps are the identity. For $\alpha>0$, one may
picture the operation as deleting the old minimum node, moving its
right subtree into the vacated position, and adding a new common root.
The formulas, not this picture, govern limit-length cases.

\begin{lemma}[The inverse is well-defined]\label{isg:lem:inverse}
The maps \eqref{isg:eq:theta} and \eqref{isg:eq:psi} are inverse bijections
between $C_\alpha$ and $\No_{\ge0}$.
\end{lemma}
\begin{proof}
Every nonempty image under $\tht$ begins with plus and is therefore
positive. Conversely, let $y=\plus\cat v>0$. If $p\pre v$, the
specified preimage extends $p$ through plus and is greater than $p$.
If $p\not\pre v$, then $v>p$: a sequence below an all-minus word must
extend that entire word through minus, contrary to the assumption.
Thus $\psiA$ takes values in the cone.

If $x=p$, the composite $\psiA\tht$ returns $p$. If
$x=p\cat\plus\cat u$, removing the first sign of
$\plus\cat p\cat u$ leaves $p\cat u$, and the first inverse case
returns $x$. In the remaining case, removing the first sign leaves
$x$, to which the second inverse case applies. The same two cases for
$v$ prove $\tht\psiA(y)=y$ for positive $y$, and zero is immediate.
Uniqueness of tails follows from the order types of the corresponding
ordinal intervals, so none of these steps requires a successor length.
\end{proof}

\begin{lemma}[Order preservation]\label{isg:lem:thetaorder}
The bijection $\tht$ is strictly increasing.
\end{lemma}
\begin{proof}
The minimum $p$ maps to the minimum $0$. On the block
\[
 U=\{p\cat\plus\cat u:u\in\No\},
\]
comparison is comparison of the tails $u$, before and after the map.
On the remaining class $V=C_\alpha\setminus(U\cup\{p\})$, adding
one common initial plus sign also preserves comparison.

It remains to compare the two blocks. Every member of $U$ is smaller
than every member $v$ of $V$: before the end of $p$, such a $v$ either
stops or first has plus where $p$ has minus. In both cases $v$ is greater
than every extension of $p$. After applying $\tht$, cancel the common
initial plus sign. One then compares an extension $p\cat u$ of $p$
with $v$, and the same first disagreement proves the same inequality.
When $\alpha=0$, the block $V$ is empty. This covers all cases.
\end{proof}

\subsection{The exact predecessor identity}
The following identity is the central combinatorial statement.
It makes the exception at the old minimum completely explicit.

\begin{theorem}[Predecessor and right-ancestor transport]\label{isg:thm:tree}
For every $x\in C_\alpha$ with $x>p$,
\begin{equation}\label{isg:eq:predidentity}
 \Pred(\tht(x))=\tht\bigl[\Pred(x)\cup\{p\}\bigr].
\end{equation}
Consequently,
\begin{align}
 \Rs(\tht(x))&=\tht[\Rs(x)],\label{isg:eq:rightidentity}\\
 \Ls(\tht(x))&=\tht[\Ls(x)\cup\{p\}].\label{isg:eq:leftidentity}
\end{align}
At $x=p$ the image has no proper ancestors, while
\begin{equation}\label{isg:eq:lost}
 \Rs(p)=\Pred(p)=\{-\beta:\beta<\alpha\},\qquad
 \tht[\Rs(p)]=\{q_\beta:\beta<\alpha\}.
\end{equation}
For $x,y>p$, simplicity is preserved and reflected:
\[
 x\pr y\quad\Longleftrightarrow\quad\tht(x)\pr\tht(y).
\]
\end{theorem}
\begin{proof}
First suppose $p\not\pre x$. No prefix $v\pre x$ can equal or extend
$p$; otherwise $p\pre x$. Thus $\tht(v)=\plus\cat v$ for all
proper prefixes $v$ of $x$. The proper prefixes of $\plus\cat x$ are
exactly the empty word, together with these $\plus\cat v$.
The empty word is $\tht(p)$, proving \eqref{isg:eq:predidentity} in this
case.

Now suppose $x=p\cat\plus\cat u$. The proper prefixes of $x$ are
\[
 \{-\beta:\beta\le\alpha\}
 \ \cup\ \{p\cat\plus\cat v:v\pr u\}.
\]
Under $\tht$ these become, respectively,
\[
 \{q_\beta:\beta<\alpha\}\cup\{0\},
 \qquad \{\plus\cat p\cat v:v\pr u\}.
\]
These are precisely the proper prefixes of $\plus\cat p\cat u$.
When $u$ is empty the second collection is empty; when $u$ has limit
length it includes exactly its proper ordinal restrictions. This proves
\eqref{isg:eq:predidentity} without a missing limit case.

Filter \eqref{isg:eq:predidentity} according to numerical comparison with
$\tht(x)$. By \cref{isg:lem:thetaorder}, comparison of images is comparison
of preimages. Since $p<x$, the extra element $\tht(p)=0$ is a left
ancestor, never a right ancestor. This gives
\eqref{isg:eq:rightidentity} and \eqref{isg:eq:leftidentity}.
Formula \eqref{isg:eq:lost} follows directly from the all-minus word $p$.
Finally, if $x,y>p$, membership of $\tht(x)$ in the predecessor set of
$\tht(y)$ is, by injectivity and \eqref{isg:eq:predidentity}, exactly
membership of $x$ in $\Pred(y)$.
\end{proof}

\begin{remark}[A stronger identity that would be false]\label{isg:rem:exception}
For $\alpha>0$ one must not assert \eqref{isg:eq:rightidentity} at $x=p$.
Already $p=-1$ has right ancestor $0$, whose image is $1$, while
$\tht(p)=0$ has no ancestors at all. The missing relations at precisely
this one point account for the extra condition in \cref{isg:thm:B}.
\end{remark}

\begin{corollary}[Initiality of the image]\label{isg:cor:initialimage}
If $\Gamma$ is initial and has minimum $p=-\alpha$, then
$\tht[\Gamma]$ is initial and has minimum $0$.
\end{corollary}
\begin{proof}
The minimum assertion follows from the order-isomorphism. The image of
$p$ has no predecessors. For $x>p$ in $\Gamma$, the right side of
\eqref{isg:eq:predidentity} lies in $\tht[\Gamma]$, because both
$\Pred(x)$ and $p$ lie in $\Gamma$.
\end{proof}

\begin{proposition}[The exact inverse obstruction]\label{isg:prop:inversetree}
Let $\Delta\subseteq\No_{\ge0}$ be nonempty and initial. Put
\[
 \Lambda_\alpha=\{0\}\cup\{q_\beta:\beta<\alpha\}.
\]
Then $\psiA[\Delta]$ is initial if and only if
$\Lambda_\alpha\subseteq\Delta$.
\end{proposition}
\begin{proof}
The empty word $0$ belongs to $\Delta$, so $p$ belongs to the inverse
image. If the inverse image is initial, it contains every $-\beta$ with
$\beta<\alpha$; applying $\tht$ gives $q_\beta\in\Delta$.

Conversely, suppose $\Lambda_\alpha\subseteq\Delta$. All predecessors
of $p$ are then in the inverse image. For any other inverse-image element
$x=\psiA(y)$, where $y>0$, \eqref{isg:eq:predidentity} gives
\[
 \Pred(x)\cup\{p\}=\psiA[\Pred(y)]\subseteq\psiA[\Delta].
\]
Hence every inverse-image element has all its predecessors in that image.
\end{proof}

\subsection{Exponent lengths and a finite picture}
\begin{proposition}[An exponent-birthday bound]\label{isg:prop:length}
For every $x\in C_\alpha$,
\[
 \len(\tht(x))\le 1+_{\mathrm{ord}}\len(x).
\]
In particular, an infinite exponent length never increases.
\end{proposition}
\begin{proof}
At $x=p$ the image length is zero. In the outside case the length is
exactly $1+\len(x)$. In the inside case, if $\len(u)=\mu$, the source
length is $\alpha+1+\mu$, and the image length is $1+\alpha+\mu$.
Monotonicity of ordinal addition gives
$1+\alpha+\mu\le1+\alpha+1+\mu$. Finally, $1+\lambda=\lambda$ for
any infinite ordinal $\lambda$: write $\lambda$ as a nonzero limit
ordinal followed by a finite tail.
\end{proof}

This is a bound for \emph{exponents}. It is not a theorem bounding the
birthday of the entire transformed surreal normal form.

\begin{center}
\begin{tikzpicture}[>=Stealth,every node/.style={font=\small},
                    x=1.30cm,y=1.0cm]
\node (a) at (0,0) {$0$};
\node (b) at (-1,-1) {$-1$};
\node (c) at (1,-1) {$1$};
\node (d) at (-0.5,-2) {$-\tfrac12$};
\node (e) at (0.6,-2) {$\tfrac12$};
\draw (a)--(b);\draw(a)--(c);\draw(b)--(d);\draw(c)--(e);
\node at (0,-2.65) {Source: a finite part of $C_1$};
\node (r) at (5,0) {$0$};
\node (s) at (5,-1) {$1$};
\node (t) at (6,-2) {$2$};
\node (u) at (4,-2) {$\tfrac12$};
\node (v) at (5.6,-3) {$\tfrac32$};
\draw(r)--(s);\draw(s)--(t);\draw(s)--(u);\draw(t)--(v);
\draw[->,thick] (1.8,-0.7)--node[above]{$\Theta_1$}(3.2,-0.7);
\node at (5,-3.55) {Image and its prefix relations};
\end{tikzpicture}
\end{center}
Here $-1\mapsto0$, $0\mapsto1$, $1\mapsto2$,
$-\tfrac12\mapsto\tfrac12$, and
$\tfrac12\mapsto\tfrac32$. The numerical labels are exponents, not
the monomials with those exponents.

\section{Strong transport of normal forms}\label{isg:sec:transport}

\begin{definition}[Bottom-normalized transport]\label{isg:def:normalization}
For $\alpha\in\On$, $n\in\N$, $p=-\alpha$, and $c=2^{-n}$, define
\begin{equation}\label{isg:eq:N}
 \Rect\left(\sum_{\gamma\ge p}r_\gamma\omega^\gamma\right)
   =c^{-1}r_p+
     \sum_{\gamma>p}r_\gamma\omega^{\tht(\gamma)}.
\end{equation}
A coefficient not in the support is interpreted as zero. The inverse is
\begin{equation}\label{isg:eq:Ninv}
 \Rect^{-1}\left(s_0+\sum_{\delta>0}s_\delta\omega^\delta\right)
   =c s_0\omega^p+
     \sum_{\delta>0}s_\delta\omega^{\psiA(\delta)}.
\end{equation}
\end{definition}

\begin{theorem}[Ambient transport]\label{isg:thm:transport}
The formulas above define inverse real-linear order-isomorphisms
\[
 \Rect:\Hahn{\R}{C_\alpha}
       \longrightarrow\Hahn{\R}{\No_{\ge0}}.
\]
They preserve support cardinality and decreasing support order type,
commute with every leading truncation, and preserve and reflect strong
summability of set-indexed families. For every such strongly summable
family,
\begin{equation}\label{isg:eq:strongtransport}
 \Rect\left(\sum_{i\in I}x_i\right)
       =\sum_{i\in I}\Rect(x_i).
\end{equation}
For a nonzero $x$,
\begin{equation}\label{isg:eq:leadingtransport}
 \hdeg(\Rect(x))=\tht(\hdeg(x)),
 \qquad
 \lc(\Rect(x))=
 \begin{cases}
  c^{-1}\lc(x),&\hdeg(x)=p,\\
  \lc(x),&\hdeg(x)>p.
 \end{cases}
\end{equation}
\end{theorem}
\begin{proof}
An order-isomorphism takes a reverse-well-ordered support to a
reverse-well-ordered support and preserves its order type. The support is
still a set: it is the image of a set under a definable class function.
Every coefficient is multiplied by a nonzero real number, so no support
point is lost and no new point is created. These observations show that
both displayed formulas are legitimate normal forms. Their inverse
identities follow coefficient by coefficient from
\cref{isg:lem:inverse}.

Both addition of normal forms and the proposed map are coefficientwise,
with the same fixed multiplier $c^{-1}$ at the bottom point. Hence the
map is real-linear. The greatest nonzero coefficient retains its sign,
because all multipliers are positive. Thus the map and its inverse
preserve order. This proves \eqref{isg:eq:leadingtransport} as well.

A leading truncation is specified by an initial segment of the decreasing
support enumeration. The same enumeration indices specify the
corresponding truncation of the transported support. Applying the map
therefore commutes with truncation, including at a limit index.

For a set-family $(x_i)$, one has the exact equality
\[
 \bigcup_i\supp(\Rect(x_i))
 =\tht\left[\bigcup_i\supp(x_i)\right].
\]
Reverse well-ordering of these unions is equivalent in both directions.
For each $\gamma$, the index set of family members contributing at
$\gamma$ is exactly the index set contributing at $\tht(\gamma)$ after
transport. The point-finiteness requirement is therefore also equivalent.
When it holds, coefficient extraction reduces \eqref{isg:eq:strongtransport}
to a finite real sum, through which the fixed coefficient multiplier
commutes. This proves the entire assertion without invoking convergence
in an order topology.
\end{proof}

\begin{remark}[No implicit completion of a subgroup]\label{isg:rem:nocompletion}
The ambient Hahn spaces in \cref{isg:thm:transport} are full spaces. A
particular initial subgroup $A$ need not be closed under all strong sums
or all infinite subseries. We define its image as the actual set or class
$\Rect[A]$, not as the space of all series with the same coefficient
groups. If a strongly summable family in $A$ has sum outside $A$,
\eqref{isg:eq:strongtransport} still holds in the ambient spaces, but is not
an assertion of an internal summation operation on $A$.
\end{remark}

\begin{remark}[What is and is not preserved]\label{isg:rem:preserved}
The map preserves the ordered chain of Archimedean classes, with their
exponents reindexed by $\tht$. It is not a homomorphism of additive
\emph{exponent} groups. Consequently \eqref{isg:eq:leadingtransport} should
not be described as preservation of a multiplicative valuation with the
same value-group structure.
\end{remark}

\section{Proof of omnific normalization}\label{isg:sec:normalization}

We now prove \cref{isg:thm:A}. Let $A$ be as there, and obtain
$\alpha,n,p,c$ from \cref{isg:prop:bottom}. Write
\[
 \Gamma=\Gamma(A),\qquad \Delta=\tht[\Gamma],
 \qquad H=\Rect[A].
\]
The ambient transport restricts to an ordered-group isomorphism onto
$H$. It remains to prove that this actual image is an initial subgroup
of $\Oz$.

\begin{lemma}[Transported coefficients]\label{isg:lem:coefficients}
The group $H$ is truncation closed and cross-sectional on $\Delta$,
and its coefficient groups are
\begin{equation}\label{isg:eq:coefftransport}
 H_0=\Z,\qquad H_{\tht(\gamma)}=A_\gamma\quad(\gamma>p).
\end{equation}
Its exponent class is exactly $\Delta$.
\end{lemma}
\begin{proof}
Truncation closure follows from \cref{isg:thm:transport}. For
$\gamma>p$, one has
$\Rect(\omega^\gamma)=\omega^{\tht(\gamma)}$. At the bottom,
$\Rect(c\omega^p)=1$. Thus all the leaders on $\Delta$ belong to $H$.
No image element can have a support point outside $\Delta$.

The formula at a positive image exponent leaves its coefficient
unchanged, and the inverse formula has the same property. This proves
both inclusions in the second identity of \eqref{isg:eq:coefftransport}.
At zero the coefficient group is $c^{-1}A_p=\Z$.
\end{proof}

\begin{proof}[Proof of \cref{isg:thm:A}]
By \cref{isg:cor:initialimage}, $\Delta$ is initial with minimum zero.
By \cref{isg:lem:coefficients}, the image is truncation closed and
cross-sectional, and all its coefficient groups are initial real
subgroups.

Consider an obligation in the dyadic right-ancestor condition for $H$:
let $\delta,\eta\in\Delta$ and $\eta\in\Rs(\delta)$.
The point $\delta$ cannot be zero. Write $\delta=\tht(\gamma)$
with $\gamma>p$. By \eqref{isg:eq:rightidentity}, there exists
$\xi\in\Rs(\gamma)$ with $\eta=\tht(\xi)$. Since $\xi>\gamma>p$,
its coefficient group was not rescaled. The initiality of $A$ gives
\[
 \D\subseteq A_\xi=H_\eta.
\]
Every target obligation is therefore satisfied. The concrete sufficiency
in \cref{isg:fact:EK} proves that $H$ is initial in $\No$.

All exponents of $H$ are nonnegative and its zero coefficient is an
integer. Hence $H\subseteq\Oz$ by \eqref{isg:eq:Oz}. Its minimum positive
element is $1$, either by order-isomorphism or directly from its normal
forms. Finally, $\Rect(\varepsilon)=1$, and all support and summability
claims are the restrictions of \cref{isg:thm:transport}. Uniqueness of
$\alpha,n$ was established in \cref{isg:prop:bottom}.
\end{proof}

\begin{corollary}[The cited existence question]\label{isg:cor:answer}
Every discrete initial subgroup of $\No$ is isomorphic, as an ordered
abelian group, to an initial subgroup of $\Oz$.
\end{corollary}

\begin{remark}[An initial image is not a simplicity isomorphism]\label{isg:rem:notsimplicity}
The conclusion is exactly an ordered-group isomorphism to an initial
image. For example, on $A=\frac12\Z$ the map $\mathcal N_{0,1}$
sends $1$ to $2$ and $1/2$ to $1$. Although $1\pr1/2$, one does not
have $2\pr1$. Thus the normalization is not generally an isomorphism
of the full simplicity-enriched structures. Any initial cyclic subgroup
of $\Oz$ is $\Z$, and the order-isomorphism from $\frac12\Z$ to
$\Z$ is forced. Thus this example rules out a generally
simplicity-preserving version. The published question does not demand
that stronger conclusion.
\end{remark}

\section{Exact images, extremal groups, and admissible scales}\label{isg:sec:classification}

\subsection{The compulsory dyadic spine}
\begin{lemma}[The spine group is initial]\label{isg:lem:K}
For every ordinal $\alpha$, the group $K_\alpha$ of
\eqref{isg:eq:Kintro} is initial in $\No$, has least positive element $1$,
and lies in $\Oz$. Its exponent class is $\Lambda_\alpha$.
\end{lemma}
\begin{proof}
The proper prefixes of $q_\beta$ are $0$ and the $q_\eta$ with
$\eta<\beta$. Thus $\Lambda_\alpha$ is initial. The finite-support
group on this class is cross-sectional and truncation closed. Its
coefficient group is $\Z$ at zero and $\D$ at every other point.
Any right ancestor lies among those positive spine points, where the
required dyadic containment holds. Apply \cref{isg:fact:EK}. The remaining
claims follow from its normal forms.
\end{proof}

\begin{proposition}[Necessary image data]\label{isg:prop:Knecessary}
If $A$ has least positive element $2^{-n}\omega^{-\alpha}$ and is
initial, then $K_\alpha\subseteq\Rect[A]$.
\end{proposition}
\begin{proof}
The image contains $1$. For each $\beta<\alpha$, \cref{isg:cor:forced}
gives $\D\omega^{-\beta}\subseteq A$. These terms are above the
bottom exponent, so their coefficients are unchanged, and
$\tht(-\beta)=q_\beta$. Thus the image contains every generating
summand of $K_\alpha$ and hence all their finite sums.
\end{proof}

\begin{proof}[Proof of \cref{isg:thm:B}]
Necessity has just been proved. Conversely, let $H\subseteq\Oz$ be
initial with $K_\alpha\subseteq H$, and put $A=\Rect^{-1}[H]$.
Since $H$ contains $1$, it has least positive element $1$ and $H_0=\Z$.
Let $\Delta=\Gamma(H)$. The inclusion of $K_\alpha$ implies
$\Lambda_\alpha\subseteq\Delta$, so
$\Gamma=\psiA[\Delta]$ is initial by \cref{isg:prop:inversetree}.
Its minimum is $p$.

By ambient transport, $A$ is an additive, truncation-closed Hahn subgroup.
Its bottom coefficient group is $c\Z$, which contains $1$ since
$c=2^{-n}$; all other coefficient groups are inherited from $H$.
It is therefore cross-sectional and has initial real coefficient groups.

We check the source right-ancestor condition. If $x\in\Gamma$ satisfies
$x>p$ and $y\in\Rs(x)$, then
$\tht(y)\in\Rs(\tht(x))$. The coefficient group at $y$ is unchanged
from that at $\tht(y)$, so it contains $\D$ by the initiality of $H$.
If instead $x=p$, its right ancestors are exactly the $-\beta$ for
$\beta<\alpha$. The hypothesis $K_\alpha\subseteq H$ says that all
coefficient groups at their images $q_\beta$ contain $\D$; again these
coefficients are unchanged. This checks the only obligations not already
covered by \eqref{isg:eq:rightidentity}.

The criterion \cref{isg:fact:EK} now proves $A$ initial. Its least positive
element is $\Rect^{-1}(1)=c\omega^p$, since the inverse is increasing.
The inverse identities show uniqueness, and a fixed bijection preserves
and reflects inclusion of its images. This proves the correspondence.
\end{proof}

\subsection{Smallest and largest initial realizations at a fixed scale}
Define
\begin{align}
 L_{\alpha,n}
   &=c\Z\omega^p\ \oplus\!
     \bigoplus_{\beta<\alpha}\D\omega^{-\beta},\label{isg:eq:L}\\
 M_{\alpha,n}
   &=\{x\in\Hahn{\R}{C_\alpha}:[\omega^p]x\in c\Z\}.
   \label{isg:eq:M}
\end{align}
Here $[\omega^p]x$ denotes coefficient extraction. The direct sum in
\eqref{isg:eq:L} again permits only finite supports; the class in
\eqref{isg:eq:M} permits every set reverse-well-ordered support in the cone.

\begin{theorem}[Extremal groups]\label{isg:thm:extremal}
The groups $L_{\alpha,n}$ and $M_{\alpha,n}$ are initial with least
positive element $c\omega^p$. Every initial group $A$ with that least
positive element satisfies
\[
 L_{\alpha,n}\subseteq A\subseteq M_{\alpha,n}.
\]
Moreover,
\begin{equation}\label{isg:eq:extremalimages}
 \Rect[L_{\alpha,n}]=K_\alpha,
 \qquad \Rect[M_{\alpha,n}]=\Oz,
 \qquad M_{\alpha,n}=c\omega^p\Oz.
\end{equation}
\end{theorem}
\begin{proof}
The first two identities in \eqref{isg:eq:extremalimages} follow directly
from the formulas: every positive image coefficient is unrestricted in
$M$, whereas precisely the finite dyadic spine terms occur for $L$.
By \cref{isg:thm:B}, the inverse images of $K_\alpha$ and $\Oz$ are
initial and have the required minimum. The lower bound on $A$ follows
from \cref{isg:cor:forced} and $A_p=c\Z$; the upper bound follows from
\cref{isg:prop:bottom}.

For the last identity, ordinary multiplication by $c\omega^p$ sends a
normal form on $[0,\infty)$ to one on $[p,\infty)$ by the arithmetic
translation $\delta\mapsto p+\delta$. Its bottom coefficient belongs
to $c\Z$, and the other coefficients range over all of $\R$ because
multiplication by $c\ne0$ is a bijection of $\R$. Conversely, translate
any support in the cone back by $-p$ and divide its coefficients by $c$.
The resulting series belongs to $\Oz$.
\end{proof}

The equality $M_{\alpha,n}=\varepsilon\Oz$ does not make the theorem
trivial. It concerns the \emph{largest} group, whose higher coefficients
are all real and whose support class is the full cone. Arithmetic
translation can destroy initiality for smaller initial subgroups, even
though it behaves correctly for this maximal one.

\begin{corollary}[Sharp size obstruction]\label{isg:cor:size}
Every set-sized initial group $A$ with least positive element
$2^{-n}\omega^{-\alpha}$ satisfies
\[
 |A|\ge\max(\aleph_0,|\alpha|).
\]
For every ordinal $\alpha$ and every $n\in\N$ this bound is attained
by $L_{\alpha,n}$. In particular, a countable initial group can have
such a least positive element only when $\alpha$ is countable.
\end{corollary}
\begin{proof}
The group contains the countably infinite cyclic subgroup generated by
its least positive member and contains distinct terms $\omega^{-\beta}$
for $\beta<\alpha$. This gives the lower bound. The finite-support
sum defining $L_{\alpha,n}$ has countable coefficient groups and an
index set of size $|\alpha|+1$, giving exactly the displayed cardinality.
Its initiality and minimum follow from \cref{isg:thm:extremal}.
\end{proof}

This is a restriction on a specified initial realization. The ordinal
$\alpha$ is not an invariant of the abstract ordered group under arbitrary
isomorphisms: normalization itself changes the least positive element
to $1$.

\subsection{Dyadic depth and the spectrum of inverse scales}
\begin{definition}[Dyadic-spine depth]\label{isg:def:depth}
Let $H\subseteq\Oz$ be a nonzero initial subgroup. If there exists an
ordinal $\beta$ such that $\D\omega^{q_\beta}\not\subseteq H$, let
$d(H)$ be the least such $\beta$. Otherwise write $d(H)=\On$ as a
formal unbounded-depth marker, not as a surreal or a set ordinal.
\end{definition}

\begin{corollary}[Admissible scales]\label{isg:cor:depth}
If $d(H)$ is an ordinal, then for every $\alpha\in\On$ and every
$n\in\N$,
\[
 \Rect^{-1}[H]\text{ is initial}
 \quad\Longleftrightarrow\quad \alpha\le d(H).
\]
If $d(H)=\On$, every ordinal $\alpha$ is allowed. Whenever the inverse
is initial, its least positive member is $2^{-n}\omega^{-\alpha}$.
The parameter $n$ imposes no further obstruction.
\end{corollary}
\begin{proof}
By \cref{isg:thm:B}, the condition is exactly that all
$\D\omega^{q_\beta}$ for $\beta<\alpha$ lie in $H$. This holds
precisely through the first missing index, inclusively, and fails at all
larger indices. In the unbounded case there is no missing index.
\end{proof}

For example, $d(K_\alpha)=\alpha$, since the exponent $q_\alpha$
is absent while all the earlier required dyadic terms are present.
Thus the endpoint in \cref{isg:cor:depth} can be a limit ordinal and must
not be replaced with a strict inequality. For a set-sized $H$, a first
missing index exists: the distinct monomials $\omega^{q_\beta}$
cannot all belong to a set. The subgroup
$\bigcup_{\alpha\in\On}K_\alpha$, interpreted as the class of all
finite dyadic-spine sums with an integer constant, is an example with
unbounded depth. It is initial by \cref{isg:fact:EK}, but is much smaller
than $\Oz$, for example it does not contain $\omega^2$.

\section{The quotient by the least positive cyclic subgroup}\label{isg:sec:quotient}

\subsection{Convexity and an actual splitting}
Fix $A$, $p$, $c$, and $\varepsilon=c\omega^p$ as before. Define
\[
 \pi_p:A\longrightarrow\Z,
 \qquad \pi_p(x)=c^{-1}[\omega^p]x,
 \qquad A^{>p}=\ker\pi_p.
\]
Although its notation emphasizes the higher exponents, $A^{>p}$ also
contains zero.

\begin{lemma}[Bottom-layer decomposition]\label{isg:lem:split}
The map $\pi_p$ is a surjective additive homomorphism,
$\pi_p(m\varepsilon)=m$, and
\[
 A=A^{>p}\oplus\Z\varepsilon
\]
as additive groups. The subgroup $\Z\varepsilon$ is convex. Every
nonzero positive element of $A^{>p}$ dominates all integer multiples of
$\varepsilon$.
\end{lemma}
\begin{proof}
Coefficient extraction is additive, and $A_p=c\Z$. This proves all
assertions about $\pi_p$. Given $x\in A$, the term
$\pi_p(x)\varepsilon$ belongs to $A$ and subtraction produces an element
of $A^{>p}$. The intersection of the two displayed summands is zero.

A nonzero element of $A^{>p}$ has greatest exponent strictly above $p$.
Its absolute value therefore exceeds $m\varepsilon$ for every positive
integer $m$. If $0\le x\le m\varepsilon$ in $A$ for some integer
$m\ge0$, its component in $A^{>p}$ must consequently be zero; thus
$x\in\Z\varepsilon$. Translation and sign reversal prove convexity.
\end{proof}

The summand $A^{>p}$ need not itself be initial as a concrete subclass
of $\No$. For example, with $A=\Z+\Z\omega$, it is $\Z\omega$,
which omits the proper prefix $1$ of $\omega$. A second exponent operation
is needed for an \emph{initial} quotient realization.

\subsection{Deleting a positive root}
Every positive surreal exponent has a unique form $\plus\cat v$.
Define
\[
 \rho:\No_{>0}\longrightarrow\No,\qquad
 \rho(\plus\cat v)=v.
\]
Unlike $\tht$, this map has no exceptional point inside its domain.
It preserves comparison and preserves and reflects prefix relations
between positive source points.

\begin{lemma}[Positive-root deletion]\label{isg:lem:rho}
If $\Delta\subseteq\No_{\ge0}$ is initial, then
$E=\rho[\Delta\setminus\{0\}]$ is initial, possibly empty. For
$\delta>0$,
\[
 \Rs(\rho(\delta))=\rho[\Rs(\delta)].
\]
\end{lemma}
\begin{proof}
If $v\pr\rho(\delta)$, then $\plus\cat v\pr\delta$ is a positive
prefix and lies in $\Delta$. Its image is $v$, establishing initiality.
Every proper prefix of $\delta=\plus\cat u$ is either $0$, which is
smaller than $\delta$, or $\plus\cat v$ with $v\pr u$.
Comparisons of the latter with $\delta$ are comparisons of $v$ with $u$.
This proves the right-ancestor identity.
\end{proof}

Let $H=\Rect[A]\subseteq\Oz$ and
\[
 H^+=\{h\in H:\ct(h)=0\}.
\]
The superscript $+$ here means ``strictly positive exponents,'' not
``positive elements.'' To avoid any order-cone interpretation, define
its transported group explicitly by
\begin{equation}\label{isg:eq:Bquotient}
 B=\left\{\sum_{\delta>0}s_\delta\omega^{\rho(\delta)}:
              \sum_{\delta>0}s_\delta\omega^\delta\in H^+\right\}.
\end{equation}

\begin{proposition}[An initial model of the quotient]\label{isg:prop:quotient}
The class $B$ in \eqref{isg:eq:Bquotient} is an initial additive subgroup of
$\No$. The normal-form map in that formula induces ordered-group
isomorphisms
\[
 A/\Z\varepsilon\cong H/\Z\cong H^+\cong B.
\]
\end{proposition}
\begin{proof}
The first isomorphism is induced by $\Rect$. Since constants of $H$
are integers and all integers lie in $H$, every coset in $H/\Z$ has a
unique representative with zero constant coefficient. The quotient
ordering agrees with the ordering of those representatives, because a
nonzero positive-exponent term dominates all integers.

The group $H^+$ is truncation closed and cross-sectional on
$\Delta\setminus\{0\}$. Its coefficient groups at those points are
the same as those of $H$. Reindexing by the order-isomorphism $\rho$
produces a truncation-closed, cross-sectional subgroup with the same
coefficient groups. Its exponent class is initial by \cref{isg:lem:rho}.
The right-ancestor identity in that lemma transfers every dyadic
obligation from $H$ to $B$. Hence $B$ is initial by \cref{isg:fact:EK}.
If there are no positive exponents, the group is $\{0\}$ and the same
conclusion is immediate.
\end{proof}

\begin{proof}[Proof of \cref{isg:thm:C}]
Combine \cref{isg:lem:split,isg:prop:quotient}. The explicit isomorphism sends
$x\in A$ to
\[
 \left(\rho_*\left(\Rect(x)-\pi_p(x)\right),\ \pi_p(x)\right),
\]
where $\rho_*$ denotes the normal-form reindexing of
\eqref{isg:eq:Bquotient}. The higher component determines the sign whenever
it is nonzero; otherwise the sign is the sign of the integer component.
This is precisely the indicated lexicographic order.
\end{proof}

\begin{corollary}[Finite iteration]\label{isg:cor:iteration}
Starting with an initial group, repeatedly quotient by its least positive
cyclic subgroup whenever the current quotient is nonzero and discrete.
At every finite stage reached in this way, the quotient has an explicit
initial realization in $\No$.
\end{corollary}
\begin{proof}
Apply \cref{isg:thm:C} at the first stage and then apply it to the initial
realization just constructed at each subsequent stage. This is ordinary
finite induction and makes no assertion at a limit stage.
\end{proof}

Limit stages are treated in \cref{isg:cf:cor:limit}, which rests on
source~02: at every stage, limit stages included, the quotient has an
explicit initial realization, it is the one above at finite stages, and
the realizations at different stages are related by compression.

\subsection{The converse construction: adjoining an integer bottom layer}
There is also an explicit converse to the lexicographic decomposition.
For an initial subgroup $B\subseteq\No$, including $B=\{0\}$, put
\[
 \sigma(\gamma)=\plus\cat\gamma,
 \qquad
 \sigma_*\left(\sum_\gamma r_\gamma\omega^\gamma\right)
      =\sum_\gamma r_\gamma\omega^{\sigma(\gamma)},
 \qquad J(B)=\Z\oplus\sigma_*[B].
\]
The map $\sigma_*$ changes exponents but not coefficients. Its image has
no constant term, and each nonzero image element dominates every ordinary
integer in absolute value.

\begin{theorem}[Integer-bottom suspension]\label{isg:thm:suspension}
The construction $B\mapsto J(B)$ gives an inclusion-preserving and
inclusion-reflecting correspondence between initial subgroups of $\No$
and nonzero initial subgroups of $\Oz$. Its inverse sends $H$ to the
initial group obtained by deleting the leading plus sign from every
positive exponent of $\ker(\ct|_H)$.
Moreover, $J(B)\cong B\lex\Z$ as ordered abelian groups.
\end{theorem}
\begin{proof}
For nonzero $B$, let $E=\Gamma(B)$. The exponent class of $J(B)$ is
$\{0\}\cup\sigma[E]$. It is initial: the prefixes of
$\plus\cat\gamma$ are zero and $\plus\cat v$ for $v\pr\gamma$.
The group is truncation closed and cross-sectional. Its coefficient
group at zero is $\Z$, while the coefficient group at
$\sigma(\gamma)$ is $B_\gamma$. Every right ancestor of
$\sigma(\gamma)$ has the form $\sigma(v)$ with $v\in\Rs(\gamma)$,
so its coefficient group contains $\D$. The criterion
\cref{isg:fact:EK} proves initiality. For $B=\{0\}$, the construction
is simply $J(B)=\Z$.

Conversely, if $H\subseteq\Oz$ is initial and nonzero, it contains
$\Z$ and each constant coefficient of $H$ is integral. Thus
$H=\Z\oplus\ker(\ct|_H)$. By \cref{isg:prop:quotient}, deleting the
initial plus sign on the positive exponents of that kernel gives an
initial group $B$. Reattaching the sign recovers the kernel exactly,
including its actual infinite-support membership conditions. Hence
$J(B)=H$, and the two constructions are inverse. Inclusions are
preserved and reflected because both operations are fixed reindexings
of the nonconstant normal forms. Finally,
$(b,m)\mapsto\sigma_*(b)+m$ is additive and has the stated
lexicographic order.
\end{proof}

\begin{corollary}[A relative criterion for discrete ordered groups]\label{isg:cor:relative}
Let $G$ be a nonzero discrete ordered abelian group with least positive
member $e$. Then $G$ admits an initial realization in $\No$ if and only
if $G/\Z e$ admits an initial realization and the exact sequence
\[
 0\longrightarrow\Z e\longrightarrow G
   \longrightarrow G/\Z e\longrightarrow0
\]
splits as a sequence of abelian groups.
\end{corollary}
\begin{proof}
If $G$ admits an initial realization, the least positive element is
carried to the least positive element there. Apply
\cref{isg:thm:C} and transport its splitting and quotient realization back.
For any discrete ordered group, $\Z e$ is convex: finite induction
shows that $[0,me]$ consists of $0,e,\ldots,me$ for $m\in\N$.
Conversely, a group-theoretic splitting identifies $G$ with
$(G/\Z e)\lex\Z$ as an ordered group: a lift of a nonzero positive
quotient element lies above every element of $\Z e$, by convexity.
Choose an initial realization $B$ of the quotient and apply
\cref{isg:thm:suspension}.
\end{proof}

This criterion is relative: it reduces initial realizability to the
quotient together with a splitting condition, rather than deciding
initial realizability of every possible quotient.

No result of this section states that every convex quotient of every
initial group has an initial realization; that was the extent of
source~01. The cyclic least-positive case is special here because it has
a normal-form coefficient projection and the exact root-deletion
operation above. Source~02 removes the restriction:
\cref{isg:cf:main:factor} realizes every convex factor, with the
coefficient projection replaced by the difference of two leading
truncations and root deletion replaced by sign-tree compression, and
\cref{isg:cf:prop:cyclic} shows that for $\Z\varepsilon$ its construction
is exactly the one of this section.

\section{Surcomplex and linear-algebraic consequences}\label{isg:sec:complex}

\subsection{Complex normal forms and rectangular modules}
Let $\No[i]=\No\oplus i\No$ with $i^2=-1$. Combining the real and
imaginary supports gives a normal form with complex coefficients and
surreal exponents; the union of the two reverse-well-ordered supports is
again reverse-well-ordered. Conjugation acts coefficientwise.
Define the Gaussian omnific module
\[
 \Oz[i]=\{a+ib:a,b\in\Oz\}.
\]
For an additive group $A\subseteq\No$, put
$A[i]=A\oplus iA$. It is a $\Z[i]$-module because multiplication by
$i$ sends $(a,b)$ to $(-b,a)$.

\begin{definition}[Rectangular initiality]\label{isg:def:rectangular}
A rectangular subclass $A+iA$ is coordinatewise initial if $A$ is
initial in $\No$. This is a deliberately coordinate-dependent notion;
no canonical single sign-tree order on $\No[i]$ is being postulated.
\end{definition}

\begin{theorem}[Surcomplex additive normalization]\label{isg:thm:complex}
Let $A$ be a nonzero discrete initial subgroup and $H=\Rect[A]$.
The map
\[
 \Rect^{\C}(a+ib)=\Rect(a)+i\Rect(b)
\]
is a $\Z[i]$-linear isomorphism
\[
 A[i]\longrightarrow H[i]\subseteq\Oz[i].
\]
It commutes with conjugation and has a complex-linear extension to the
full complex-coefficient Hahn spaces on $C_\alpha$ and
$\No_{\ge0}$. That extension preserves and reflects set-indexed strong
summability and commutes with its sums. The image is coordinatewise
initial. For fixed $\alpha,n$, its exact possible rectangular images
are the $H[i]$ with $H\subseteq\Oz$ initial and $K_\alpha\subseteq H$.
\end{theorem}
\begin{proof}
Additivity and multiplication by $i$ commute with the coordinatewise
map, so it is $\Z[i]$-linear. The inverse is coordinatewise
$\Rect^{-1}$. Conjugation negates the imaginary coordinate, and
additivity gives the asserted compatibility.

In complex normal forms, the map is exactly
\[
 \sum_{\gamma\ge p}z_\gamma\omega^\gamma
 \ \longmapsto\ c^{-1}z_p+
      \sum_{\gamma>p}z_\gamma\omega^{\tht(\gamma)}.
\]
The coefficient multiplier $c^{-1}$ is real. It commutes with every
complex scalar and with conjugation. The support argument in
\cref{isg:thm:transport} is unchanged; there is no additional infinite sum
of constants because strong summability remains point-finite.
Coordinatewise initiality follows from the initiality of $H$, and the
range assertion is exactly \cref{isg:thm:B} applied to the real factor.
\end{proof}

\subsection{What linear systems can be transported}
\begin{corollary}[Exact transport of integral linear equations]\label{isg:cor:linear}
Let $M\in\Z^{r\times s}$, let $b\in A^r$, and apply $\Rect$
coordinatewise. Then
\[
 Mx=b\quad\Longleftrightarrow\quad M\Rect(x)=\Rect(b).
\]
The corresponding statement holds for $M\in\Z[i]^{r\times s}$ and
$b\in A[i]^r$ under $\Rect^{\C}$. In the real case, coordinatewise
inequalities are preserved as well.
\end{corollary}
\begin{proof}
Every row is a finite integer or Gaussian-integer linear combination.
The relevant module isomorphism commutes with each such combination,
and its injectivity gives the reverse implication. The real order
assertion follows from the ordered-group isomorphism.
\end{proof}

The right-hand side is transported too. In general $\Rect(1)\ne1$,
so this corollary must not be rewritten as preservation of equations
with every ordinary constant fixed. Nor does it apply to nonlinear
polynomials, arbitrary surreal scalar matrices, norms, or inner products.

\section{Examples and failure modes}\label{isg:sec:examples}

\subsection{The motivating two-level example}
Ehrlich and Kaplan's example in the paragraph preceding their question
is
\[
 A=\D+\Z\omega^{-1},\qquad \varepsilon=\omega^{-1}.
\]
Our construction has $\alpha=1,n=0$, $\Theta_1(-1)=0$, and
$\Theta_1(0)=1$. It gives
\[
 d+m\omega^{-1}\longmapsto d\omega+m,
 \qquad H=\D\omega+\Z=K_1.
\]
Thus the proposed general construction specializes to their known map,
rather than treating that map as a new result \cite[Section~9]{EK}.
The example and the map stand in Section~9 of the journal version as
well, immediately before its Question~9.1.

\subsection{Rescale the bottom coefficient, not the whole series}
Consider
\[
 A=\tfrac12\Z+\Z\omega.
\]
The exponent class $\{0,1\}$ is initial, there are no right-ancestor
dyadic obligations, and its coefficient groups $\frac12\Z$ and $\Z$
are initial. Hence $A$ is initial by \cref{isg:fact:EK}. Its least positive
element is $1/2$.

Scalar multiplication by $2$ yields $\Z+2\Z\omega$, which omits
$\omega$ despite containing $2\omega$. Equivalently, its coefficient
group at exponent $1$ is not initial in $\R$. By contrast,
\[
 \mathcal N_{0,1}(a/2+b\omega)=a+b\omega
\]
has initial image $\Z+\Z\omega$. Only the bottom coefficient changes.

\subsection{An exponent translation fails at the first infinite branch}
Let
\[
 \Gamma=\{-1\}\cup\N\cup\{\omega\}
\]
be an exponent class, and let $A$ be the finite-support group with
coefficient $\Z$ at $-1$, coefficient $\D$ at $0$, and coefficient
$\Z$ at every positive member of $\Gamma$. This exponent class is
initial: $-1$ contributes the prefix $0$, and the ordinal $\omega$
contributes all finite ordinals. Its only right-ancestor obligation is
at $-1$, forcing the chosen dyadic group at zero. Thus $A$ is initial,
with least positive element $\omega^{-1}$.

Multiplication by $\omega$ replaces the support class by
\[
 \Gamma+1=\{0,1,2,\ldots\}\cup\{\omega+1\}.
\]
This class is not initial because it contains $\omega+1$ but omits
its prefix $\omega$. The resulting group is not initial. In contrast,
root relocation satisfies
\[
 \Theta_1(-1)=0,\qquad
 \Theta_1(k)=k+1\ (k<\omega),\qquad
 \Theta_1(\omega)=\omega.
\]
The last identity uses $\plus\cat\plus^\omega=\plus^\omega$;
it is precisely where ordinal concatenation differs from arithmetic
translation. The correct image has exponent class
$\N\cup\{\omega\}$, coefficient $\Z$ at zero, $\D$ at one,
and $\Z$ at the other exponents. It is initial.

\subsection{A genuine limit-support example}
Take $\alpha=\omega$, $n=0$, and consider the group of all series
\begin{equation}\label{isg:eq:limitexample}
 x=\sum_{k<\omega}a_k\omega^{-k}+m\omega^{-\omega},
 \qquad a_k\in\D,\quad m\in\Z.
\end{equation}
Zeros in this expression are ignored. The allowed support is a subset
of the reverse-well-ordered sequence
$0,-1,-2,\ldots,-\omega$. The exponent class is initial, its ancestor
coefficients are dyadic, and its bottom coefficient is integral, so the
full group \eqref{isg:eq:limitexample} is initial.

Since $\Theta_\omega(-k)=q_k=2^{-k}$ for finite $k$ and
$\Theta_\omega(-\omega)=0$, normalization gives the exact strong-sum
identity
\begin{equation}\label{isg:eq:limitimage}
 \mathcal N_{\omega,0}(x)
   =\sum_{k<\omega}a_k\omega^{2^{-k}}+m.
\end{equation}
For example, when $a_k=1$ for every $k$ and $m=1$, the source and image
have support order type $\omega+1$. The decreasing exponents
$1,1/2,1/4,\ldots,0$ are a legitimate reverse-well-ordered support;
no analytic limiting process is involved.

The finite-support minimal group $L_{\omega,0}$ is a proper subgroup of
the group in \eqref{isg:eq:limitexample}. Its image $K_\omega$ consists
only of finite sums of the positive terms appearing in
\eqref{isg:eq:limitimage}, together with an integer constant. The full
limit-support example therefore illustrates why the image cannot be
reconstructed from its coefficient groups alone.

\subsection{Inverse initiality needs the whole dyadic spine}
Let $H=\Z$ and $\alpha=1,n=0$. The inverse formula produces
\[
 \mathcal N_{1,0}^{-1}[\Z]=\Z\omega^{-1},
\]
which is not initial: $1$ is simpler than $\omega^{-1}$ but is missing.
Here $K_1\not\subseteq H$, so \cref{isg:thm:B} correctly rules out
initiality.

There is also a purely coefficient-level obstruction. Let
$H=\Z+\Z\omega$ and keep $\alpha=1,n=0$. Its exponent class contains
$\Lambda_1=\{0,1\}$, so the inverse exponent class is initial.
Nevertheless, the inverse group is
\[
 \Z+\Z\omega^{-1},
\]
which is not initial: its monomial $\omega^{-1}$ forces every dyadic
constant through \eqref{isg:eq:dyadic}, but its constant group is only
$\Z$. Equivalently, $1/2$ is a missing prefix of $\omega^{-1}$.
Thus the exponent condition in \cref{isg:prop:inversetree} alone does not
suffice for the group-level inverse.

\subsection{Multiplication and ordinary constants are not preserved}
For $A=\D+\Z\omega^{-1}$, one has $1\in A$ and
$\mathcal N_{1,0}(1)=\omega$. The product $1\cdot1$ also belongs to
$A$, but
\[
 \mathcal N_{1,0}(1\cdot1)=\omega
 \ne\omega^2=
 \mathcal N_{1,0}(1)\mathcal N_{1,0}(1).
\]
Even on products that happen to stay in the source, multiplicativity
fails. The surcomplex extension inherits this limitation. In particular,
no preservation of the complex modulus or of multiplication by general
surcomplex scalars follows from its Gaussian-linearity.

\section{Convex subgroups as exponent cuts}\label{isg:cf:sec:cuts}

This section and the next seven are source~02's
(\cref{isg:cf:sec:intro}). They extend \cref{isg:sec:quotient} from the
least positive cyclic subgroup to every convex subgroup.

\subsection{Conventions of the second source}\label{isg:cf:sec:conventions}
All statements about arbitrary set-sized ordered groups can be read in
ZFC. For the uniform class constructions we work, as in
\cref{isg:sec:background}, in NBG with global choice, allowing class
parameters. A quotient of a proper-class group is not treated as a set of
proper-class cosets. When such notation is used, the quotient is
represented by the explicitly constructed normal-form complement. The
class operations restrict to set operations on every individual support.
In particular, none of the proofs sums a proper class of terms.

Signs, prefixes, $\Pred$, $\Rs$, normal forms, leading truncations,
strong summability, $G_\gamma$ and cross-sectionality are as in
\cref{isg:sec:background}. For two sign sequences $x,y$, let $x\wedge y$
be their longest common prefix; if one is a prefix of the other, their
meet is the shorter one. Source~02 uses the Ehrlich--Kaplan criterion in
exactly the concrete form of \cref{isg:fact:EK}, including its proof that
the canonical normal-form image is initial, and the description of initial
real coefficient groups in \cref{isg:fact:real}; on an active exponent such
a group is nonzero and contains $1$. It imports one more result.

\begin{fact}[Ehrlich's divisible embedding]\label{isg:cf:fact:ehrlich}
Every divisible ordered abelian group has an initial realization in
$\No$ \cite{Ehrlich}.
\end{fact}
This is Ehrlich's earlier embedding theorem, also recalled in \cite{EK}.
It is used below only for set-sized groups. Neither source re-proves
\cref{isg:fact:EK} or \cref{isg:cf:fact:ehrlich}.

\begin{remark}[Actual membership is part of the data]\label{isg:cf:rem:membership}
A truncation-closed subgroup need not contain every series allowed by its
individual coefficient groups. In particular it need not be closed under
arbitrary infinite subseries or all ambient strong sums. All transported
groups below are actual images, not silently enlarged Hahn completions.
This is the caution of \cref{isg:rem:nocompletion}, stated by both
sources.
\end{remark}

Several symbols of source~02 are renamed here, because the letters are
already in use in \cref{isg:sec:intro,isg:sec:background,isg:sec:bottom,isg:sec:surgery,isg:sec:transport,isg:sec:normalization,isg:sec:classification,isg:sec:quotient,isg:sec:complex,isg:sec:examples}.
The renaming changes no statement.
\begin{center}\small
\begin{tabular}{@{}lll@{}}
\toprule
here & in source~02 & reason for the change\\
\midrule
$\res_G$ & $\rho_G$ & $\rho$ is root deletion (\cref{isg:lem:rho})\\
$G^{\dv}$ & $D(G)$ & $\D$ is $\Z[1/2]$\\
$\zeta,\ \zeta_*\in\Zhat$; $G_\zeta$ & $\alpha,\ \alpha_*$; $G_\alpha$ &
 $\alpha$ is the ordinal of the bottom exponent $-\alpha$\\
$\tau_*$ & $p_*$ & $p=-\alpha$ is the bottom exponent\\
$W=\Q\omega^2+\Z\omega+\Q$ & $K$ & $K_\alpha$ is the spine group\\
$\widetilde G_\zeta=G_\zeta\lex\Q$ & $H_\alpha$ & $H$ is the image in
 \cref{isg:thm:A}\\
$C\subseteq C'$ & $C_1\subseteq C_2$ & $C_\alpha$ is the cone
 $[-\alpha,\infty)$\\
$\Gamma_C$, $\Gamma^C=\Gamma\setminus\Gamma_C$ & $\Lambda_C$, $\Delta$ &
 $\Lambda_\alpha$ is the spine exponent class\\
$C^\natural=G[\Gamma^C]$ & $H_C$ & as for $H$\\
$\cmp E$, $\iota_x$ & $c_E$, $e_x$ & $c=2^{-n}$; $e$ is a least positive
 element\\
$\pi_1$ & $\pi$ & $\pi_p$ is the bottom coefficient projection\\
Theorems D, E, F & Theorems A, B, C & source~01's letters are kept\\
\bottomrule
\end{tabular}
\end{center}
Kept unchanged: $e$ for the least positive element of an abstract
discrete group (as in \cref{isg:cor:relative}); $\Pi$ for the purely
infinite normal forms, zero included (as in \cref{isg:sec:conventions});
$J(B)$ for the suspension of \cref{isg:thm:suspension}; ``coordinatewise
initial'' of \cref{isg:def:rectangular}, which source~02 calls
``rectangularly initial''. In
\cref{isg:cf:sec:presburger,isg:cf:sec:factorial,isg:cf:sec:modeltheory,isg:cf:sec:complex}
the letter $n$ is an ordinary modulus, not the dyadic exponent of
\cref{isg:prop:bottom}. The profinite residue $\res_G$ is the residue
$\res_{\mathcal I}$ of the trigonometry report
(\replabel{trigonometry:per:eq:residue}), and $G^{\dv}$ its
$\mathcal I^{\mathrm{div}}$ (\replabel{trigonometry:per:prop:divisible}).

\subsection{Exponent cuts and the canonical splitting}
Throughout \cref{isg:cf:sec:cuts,isg:cf:sec:compression,isg:cf:sec:factors}
let $G\subseteq\No$ be initial, with active exponent class
$\Gamma=\Gamma(G)$. For an exponent subclass $E$ define the ambient
coefficient restriction
\[
 P_E(x)=\sum_{\gamma\in\supp(x)\cap E}r_\gamma\omega^\gamma.
\]
It is always an ambient normal form. Its membership in $G$ needs proof.

\begin{lemma}[Archimedean detection]\label{isg:cf:lem:arch}
Let $C$ be a convex subgroup of $G$ and set
\[
 \Gamma_C=\{\gamma\in\Gamma:\omega^\gamma\in C\}.
\]
Then $\Gamma_C$ is a numerical lower segment of $\Gamma$. For nonzero
$x\in G$, with greatest exponent $\gamma$, the following are equivalent:
\[
 x\in C;\qquad \gamma\in\Gamma_C;\qquad
 \supp(x)\subseteq\Gamma_C.
\]
\end{lemma}
\begin{proof}
The numbers $|x|$ and $\omega^\gamma$ are Archimedean equivalent:
each is bounded by a positive ordinary integer multiple of the other.
Indeed their ratio has positive, nonzero real leading coefficient and
an infinitesimal remainder. Convexity and subgroup closure therefore
show that $x\in C$ exactly when $\omega^\gamma\in C$.
If $\delta<\gamma$ lies in $\Gamma$, then
$0<\omega^\delta<\omega^\gamma$, so membership of the latter in $C$
forces membership of the former. This proves the lower-segment assertion
and the support equivalence. Zero belongs to every subgroup separately.
\end{proof}

\begin{lemma}[Restriction to a convex interval]\label{isg:cf:lem:restriction}
If $E$ is a convex subclass of the ordered class $\Gamma$, then
$P_E(G)\subseteq G$. More precisely, $P_E$ is the difference of two
leading-truncation projections. Its image equals
\[
 G[E]:=\{x\in G:\supp(x)\subseteq E\}.
\]
This subgroup is truncation closed and cross-sectional on $E$.
\end{lemma}
\begin{proof}
There is nothing to prove when $E$ is empty. Otherwise set
\[
 U=\{\gamma\in\Gamma:(\exists\delta\in E)\ \delta\le\gamma\},
 \qquad
 V=\{\gamma\in\Gamma:(\forall\delta\in E)\ \delta<\gamma\}.
\]
Both are upper segments, $V\subseteq U$, and convexity gives $E=U\setminus V$.
For every reverse-well-ordered support, restriction to an upper segment
retains a leading segment. Truncation closure, with zero and the entire
form admitted, gives $P_U(x),P_V(x)\in G$. Hence
$P_E(x)=P_U(x)-P_V(x)\in G$. The image description follows from
$P_E^2=P_E$. Leading truncations of an element supported in $E$ remain
supported there, and every $\omega^\gamma$ with $\gamma\in E$ lies in
the image.
\end{proof}

\begin{theorem}[Canonical splitting in a chosen realization]\label{isg:cf:thm:split}
For a convex subgroup $C$ of $G$, put $\Gamma^C=\Gamma\setminus\Gamma_C$.
Then
\[
 C=G[\Gamma_C],\qquad C^\natural:=G[\Gamma^C],\qquad G=C^\natural\oplus C.
\]
The map $x\mapsto P_{\Gamma^C}(x)$ realizes $G/C$ as the ordered subgroup
$C^\natural$, and addition induces an ordered-group isomorphism
\[
 C^\natural\lex C\ \longrightarrow\ G,\qquad (h,c)\longmapsto h+c.
\]
\end{theorem}
\begin{proof}
The first equality is \cref{isg:cf:lem:arch}. Since $\Gamma^C$ is an
upper segment, $P_{\Gamma^C}(x)$ is a leading truncation in $G$. The
difference $x-P_{\Gamma^C}(x)$ is supported in $\Gamma_C$ and belongs to
$C$. The supports of the two summands are disjoint, proving uniqueness.
If $h\ne0$ and $c\in C$, the leading exponent of $h$ exceeds every
exponent of $c$. Consequently $h+c$ has the sign of $h$; when $h=0$
it has the sign of $c$. This is the lexicographic ordering.
The kernel of $P_{\Gamma^C}$ is $C$, so it gives the stated quotient model
and the inclusion of $C^\natural$ is an order-preserving section.
\end{proof}

For a discrete initial $A$ and $C=\Z\varepsilon$ one has
$\Gamma_C=\{p\}$ and $C^\natural=A^{>p}$, and \cref{isg:cf:thm:split} is
\cref{isg:lem:split} (see \cref{isg:cf:prop:cyclic}).

\begin{remark}[Which projection preserves order?]\label{isg:cf:rem:projection}
The upper projection $P_{\Gamma^C}$ is nondecreasing. The lower projection
$P_{\Gamma_C}$ usually is not: for example, a positive element
$\omega-1$ has negative constant projection. Splitting the abelian-group
sequence does not make both component projections positive maps.
\end{remark}

\begin{proposition}[Convex factors and commuting projections]\label{isg:cf:prop:factorcoordinates}
If $C\subseteq C'$ are convex in $G$, then
$\Gamma_C\subseteq\Gamma_{C'}$ and $E=\Gamma_{C'}\setminus\Gamma_C$
is convex in $\Gamma$. Restriction induces an order-isomorphism
\[
 C'/C\ \cong\ G[E].
\]
For arbitrary exponent subclasses $E,F$ the ambient identities are
\begin{equation}\label{isg:cf:eq:projectionlaw}
 P_E P_F=P_{E\cap F}=P_F P_E.
\end{equation}
If $E,F$ are convex intervals, all these maps preserve $G$.
\end{proposition}
\begin{proof}
The lower-segment descriptions give the interval assertion. Every
$x\in C'$ splits uniquely into its $E$-part and its $\Gamma_C$-part.
The first is a representative of the quotient, and its leading
coefficient determines the sign of a nonzero coset. The identities
\eqref{isg:cf:eq:projectionlaw} are coefficientwise. Intersections of
convex intervals are convex, so \cref{isg:cf:lem:restriction} handles
membership.
\end{proof}

These arguments already yield an intrinsic obstruction: an ordered
abelian group with a nonsplit convex subgroup sequence cannot be
initially realizable. The next step shows that the explicit complements
and interval subgroups can themselves be given initial realizations.

\section{Compression of a meet-closed sign tree}\label{isg:cf:sec:compression}

The central construction is elementary at the level of signs, but its
retained positions must be indexed by their \emph{order type}, not by
ordinary integer subtraction.

\begin{definition}[Retained-ancestor compression]\label{isg:cf:def:compression}
A nonempty subclass $E\subseteq\No$ is \emph{meet closed} if
$x\wedge y\in E$ for all $x,y\in E$. For $x\in E$ let
\[
 I_E(x)=\{\beta<\len(x):x\restriction\beta\in E\},
 \qquad \kappa_E(x)=\otp(I_E(x)).
\]
Let $\iota_x:\kappa_E(x)\to I_E(x)$ be the unique increasing enumeration.
Define $\cmp E(x)$ by
\begin{equation}\label{isg:cf:eq:compression}
 \len(\cmp E(x))=\kappa_E(x),\qquad
 \cmp E(x)(\xi)=x(\iota_x(\xi)).
\end{equation}
Thus the sign at position $\beta$ is retained exactly when its preceding
prefix belongs to $E$.
\end{definition}

The set $I_E(x)$ is a subset of the ordinal $\len(x)$, even when $E$ is
a proper class. The definition is therefore elementwise legitimate.

\begin{lemma}[Convex intervals are meet closed]\label{isg:cf:lem:meetclosed}
If $\Gamma$ is initial in $\No$ and $E\subseteq\Gamma$ is a nonempty
numerically convex subclass, then $E$ is meet closed.
\end{lemma}
\begin{proof}
If one of $x,y\in E$ is a prefix of the other, their meet is an endpoint.
Otherwise, their common prefix $m=x\wedge y$ belongs to $\Gamma$ by
initiality. At the first differing sign, one continuation is minus and
the other plus. Thus $m$ lies strictly between $x$ and $y$ numerically,
and convexity puts $m$ in $E$.
\end{proof}

\begin{lemma}[Retained prefix formula]\label{isg:cf:lem:retainedprefix}
If $y=x\restriction\beta\in E$ with $\beta<\len(x)$, then
\[
 I_E(y)=I_E(x)\cap\beta.
\]
Writing $\xi=\otp(I_E(x)\cap\beta)$, one has
\[
 \cmp E(y)=\cmp E(x)\restriction\xi,
 \qquad \xi<\kappa_E(x),\qquad \iota_x(\xi)=\beta.
\]
\end{lemma}
\begin{proof}
The prefixes of $y$ are exactly the prefixes of $x$ of length less than
$\beta$. Since $\beta\in I_E(x)$, it is the next retained position after
those earlier positions. The signs in the two compressed sequences
agree at all earlier retained positions by definition.
\end{proof}

\begin{theorem}[Compression theorem]\label{isg:cf:thm:compression}
For nonempty meet-closed $E\subseteq\No$, the map
$\cmp E:E\to\No$ is an order-isomorphism onto an initial subclass
$E^\flat:=\cmp E[E]$. For $x,y\in E$,
\begin{align}
 y\pr x&\quad\Longleftrightarrow\quad
       \cmp E(y)\pr\cmp E(x),\label{isg:cf:eq:prefixreflect}\\
 \cmp E(x\wedge y)&=\cmp E(x)\wedge\cmp E(y),\label{isg:cf:eq:meetreflect}\\
 \Rs(\cmp E(x))&=\cmp E[\Rs(x)\cap E].\label{isg:cf:eq:rightreflect}
\end{align}
It also preserves and reflects left-ancestor relations.
\end{theorem}
\begin{proof}
First suppose $y\pr x$. The retained prefix formula shows that
$\cmp E(y)$ is a proper prefix of $\cmp E(x)$. At its first new
position the compressed continuation has sign $x(\len(y))$.
Consequently the numerical comparison between $y$ and $x$ is unchanged.

Now suppose $x,y$ are incomparable by prefix and put $m=x\wedge y$.
Meet closure gives $m\in E$. The retained positions preceding
$\len(m)$ are identical for $x$ and $y$, and the position $\len(m)$
itself is retained in both. Their compressed sequences therefore have
common prefix $\cmp E(m)$ and then opposite signs. The first differing
sign has not changed. This proves order preservation, injectivity, and
\eqref{isg:cf:eq:meetreflect} in the incomparable case. The comparable
case of the meet identity follows from the retained prefix formula.
These alternatives also prove the reverse implication in
\eqref{isg:cf:eq:prefixreflect}: incomparable points remain incomparable.

Let $z$ be an arbitrary proper prefix of $\cmp E(x)$, of length
$\xi<\kappa_E(x)$. Put $\beta=\iota_x(\xi)$ and
$y=x\restriction\beta$. By definition $y\in E$, and the retained
prefix formula gives $\cmp E(y)=z$. Thus $E^\flat$ is initial, at
successor and limit prefix lengths alike. Combining this description
of every prefix with order preservation proves
\eqref{isg:cf:eq:rightreflect}; the left-ancestor assertion is identical
with the inequality reversed.
\end{proof}

\begin{remark}[No unjustified limit step]\label{isg:cf:rem:limit}
A limit ordinal $\xi<\kappa_E(x)$ still has the position $\iota_x(\xi)$.
The proof uses that position directly. It does not assume $E$ is closed
under limits of chains or that an infinite sequence of deletions
converges to a surreal. This is why the construction applies to cuts
without endpoints.
\end{remark}

\begin{proposition}[Uniqueness and nested coherence]\label{isg:cf:prop:coherence}
The map $\cmp E$ is the unique isomorphism of the numerical and prefix
orders from $E$ onto an initial subclass of $\No$. If $F\subseteq E$
is nonempty and meet closed, then
\begin{equation}\label{isg:cf:eq:coherence}
 \cmp{\cmp E[F]}\bigl(\cmp E(x)\bigr)=\cmp F(x)
 \qquad(x\in F).
\end{equation}
\end{proposition}
\begin{proof}
In an initial image, the proper prefixes of the image of $x$ correspond
exactly to the proper ancestors of $x$ in $E$, in their simplicity
order. Hence the image length must be $\kappa_E(x)$. The sign at the
position corresponding to $y\pr x$ is plus exactly when $y<x$.
Both the positions and signs are forced; they agree with
\eqref{isg:cf:eq:compression}.

By \eqref{isg:cf:eq:meetreflect}, the subclass $\cmp E[F]$ is meet closed.
In a second compression, a sign of $\cmp E(x)$ is retained exactly
when its preceding prefix is $\cmp E(y)$ for some $y\in F$.
By the retained prefix formula and prefix reflection, these are exactly
the original positions $\beta$ with $x\restriction\beta\in F$.
Thus the two-stage operation retains the same well-ordered positions,
in the same order, as direct compression by $F$. This proves
\eqref{isg:cf:eq:coherence}.
\end{proof}

\begin{example}[Compression is not exponent translation]\label{isg:cf:ex:compression}
Take the full finite sign tree of lengths at most two, in numerical order:
\[
 --\ <\ -\ <\ -+\ <\ 0\ <\ +-\ <\ +\ <\ ++.
\]
On the convex interval $E=\{-+,0,+-\}$, compression gives
\[
 -+\longmapsto -,
 \qquad 0\longmapsto0,
 \qquad +-\longmapsto+.
\]
The original exponents are $-1/2,0,1/2$, and their images are $-1,0,1$.
On $E=\{+,++\}$ the images are $0,+$. The definition is a tree operation,
not a uniform rule of subtracting a real or surreal endpoint.
\end{example}

\begin{example}[Meet closure cannot be omitted]\label{isg:cf:ex:badmeet}
For $E=\{-,+\}$ neither point has a proper prefix in $E$. Both therefore
compress to $0$, so injectivity fails. This set omits the meet $0$.
Convexity in an initial exponent class prevents precisely this defect.
\end{example}

The compression theorem is compatible with the general viewpoint of
surreal substructures and convex partitions in \cite{BagHoeven}.
It does not assert that a set-sized convex interval is a copy of the
whole proper-class structure $\No$; its target is an arbitrary initial
subclass, often small.

\section{Initial realizations of all convex factors}\label{isg:cf:sec:factors}

For the empty interval we set $E^\flat=\varnothing$,
$\mathcal F_E(G)=\{0\}$, and let $\Phi_E$ be the unique map between
the zero Hahn spaces. This convention also covers equal convex subgroups
in the factor statements.

\begin{theorem}[Hahn transport on an interval]\label{isg:cf:thm:transport}
Let $G\subseteq\No$ be initial with exponent class $\Gamma$, and let
$E\subseteq\Gamma$ be a nonempty convex interval. Define
\begin{equation}\label{isg:cf:eq:transport}
 \Phi_E\left(\sum_{\gamma\in E}r_\gamma\omega^\gamma\right)
  =\sum_{\gamma\in E}r_\gamma\omega^{\cmp E(\gamma)}.
\end{equation}
Then $\Phi_E$ is an ordered real-linear isomorphism
$\Hahn{\R}{E}\to\Hahn{\R}{E^\flat}$. It restricts to an ordered-group
isomorphism of $G[E]$ onto an initial subgroup $\mathcal F_E(G)$ of
$\No$. It preserves support order types, leading coefficients, leading
truncations, and strong summability in both directions.
\end{theorem}
\begin{proof}
By \cref{isg:cf:lem:meetclosed,isg:cf:thm:compression}, the exponent map
is an order-isomorphism with initial image. It therefore sends each
set-sized reverse-well-ordered support to another such support, with the
same order type. Its inverse has the same property. Coefficients are
unchanged, so addition and real scalar multiplication commute with the
map, and the leading coefficient retains its sign.

Set $B=\Phi_E[G[E]]$. By \cref{isg:cf:lem:restriction}, the source
subgroup is cross-sectional and truncation closed. These properties pass
to $B$, and its coefficient group at $\cmp E(\gamma)$ is exactly
$G_\gamma$. The latter is initial in $\R$. Every right ancestor of
$\cmp E(\gamma)$ has the form $\cmp E(\delta)$ with
$\delta\in\Rs(\gamma)\cap E$, by \eqref{isg:cf:eq:rightreflect}.
The original dyadic obligation \eqref{isg:eq:dyadic} supplies the required
coefficient group. The sufficiency direction of \cref{isg:fact:EK} now
shows that the actual image $B$ is initial.

Leading truncations retain the same initial sets of support indices
before and after reindexing. For a set-family $(x_i)$, the union of its
transported supports is the image under $\cmp E$ of the original union.
Reverse well-ordering is equivalent in both directions. At any exponent,
the set of indices contributing a nonzero coefficient is unchanged.
Point-finiteness is therefore also equivalent. In a strongly summable
family, every resulting coefficient is a finite real sum, and the
transport commutes with that sum.
\end{proof}

This is the analogue, for an interval and without a rescaled bottom
coefficient, of \cref{isg:thm:transport}; the proof has the same
structure.

\begin{proof}[Proof of \cref{isg:cf:main:factor}]
Transport $G$ to an actual initial realization. Apply
\cref{isg:cf:prop:factorcoordinates} to obtain
$C'/C\cong G[\Gamma_{C'}\setminus\Gamma_C]$.
If the exponent interval is empty, the factor is zero. Otherwise
\cref{isg:cf:thm:transport} gives the required initial image. Splitting is
\cref{isg:cf:thm:split}; the composition laws are
\eqref{isg:cf:eq:projectionlaw} and \eqref{isg:cf:eq:coherence}.
\end{proof}

\begin{corollary}[Subgroups, quotients, and transfinite cuts]\label{isg:cf:cor:closure}
The class $\IR$ is closed under convex subgroups, convex ordered
quotients, and their subquotients. For a chain of convex subgroups
$(C_\xi)_{\xi<\lambda}$ of an initially realizable group, any union or
intersection which is a convex subgroup has the same conclusion for
itself and for the corresponding quotient. No successor-stage
restriction is necessary.
\end{corollary}
\begin{proof}
The first assertions are special cases of \cref{isg:cf:main:factor}. A
union of an increasing chain and an intersection of convex subgroups are
convex. The theorem applies directly to their exponent cuts, rather
than asserting that images previously chosen at earlier stages have a
particular union or intersection.
\end{proof}

\begin{corollary}[Coherent compressed factors]\label{isg:cf:cor:coherentfactors}
For nonempty convex intervals $F\subseteq E\subseteq\Gamma$, the direct
map $\Phi_F$ agrees with first transporting by $\Phi_E$ and then
compressing the interval $\cmp E[F]$. Thus the two constructions give
the same normal forms, not just abstractly isomorphic groups.
\end{corollary}
\begin{proof}
The coefficient at every retained exponent is unchanged in both
constructions, and \eqref{isg:cf:eq:coherence} identifies the final
exponents. Actual membership follows from the restriction and image
definitions.
\end{proof}

\begin{remark}[Strong sums and factor projection]\label{isg:cf:rem:strongsums}
The reindexing $\Phi_E$ preserves and reflects strong summability on its
ambient space. The restriction $P_E$ preserves strong summability but
need not reflect it, because terms outside $E$ are discarded. Neither
assertion enlarges $G$ to a group closed under all ambient strong sums.
\end{remark}

\subsection{Examples with an infinite exponent cut}\label{isg:cf:sec:cutexamples}
Let $G=\Q[\omega]$ additively, meaning finite sums
$\sum_{j=0}^{m}q_j\omega^j$. The exponent class $\N$ consists of
all finite all-plus sequences and is initial. There are no right-ancestor
obligations within this class, so the criterion shows that $G$ is
initial. The convex subgroup of degree less than $k$ has quotient
represented by degrees at least $k$; compression sends the ordinal
exponent $k+j$ to $j$. This recovers an initial quotient isomorphic to
$G$.

For a genuinely transfinite example, take finite-support rational forms
on the initial exponent class of all ordinals below $\omega+\omega$.
The lower segment of finite exponents defines a convex subgroup. On the
upper interval $[\omega,\omega+\omega)$, compression deletes the first
$\omega$ ancestor positions and sends $\omega+n$ to $n$. The quotient
again has an initial rational-polynomial realization. This calculation
uses sign sequences, not an assumption that general surreal subtraction
preserves initiality.

\subsection{The cyclic case is root relocation followed by root deletion}
\label{isg:cf:sec:cyclic}
The following statement, added in the merge, identifies source~02's
construction for $C=\Z\varepsilon$ with the construction of
\cref{isg:sec:quotient}. Neither source states it.

\begin{proposition}[Compression of the bottom point; \textup{[merge]}]\label{isg:cf:prop:cyclic}
Let $A$ be a nonzero discrete initial subgroup, with $\alpha$, $n$,
$p=-\alpha$, $c$, $\varepsilon=c\omega^p$, $\tht$, $\Rect$, $\rho$,
$\pi_p$, $A^{>p}$, $H$ and $B$ as in
\cref{isg:sec:bottom,isg:sec:surgery,isg:sec:transport,isg:sec:normalization,isg:sec:quotient},
and let $\Gamma=\Gamma(A)$ and $C=\Z\varepsilon$. Then
$\Gamma_C=\{p\}$, $E:=\Gamma^C=\Gamma\setminus\{p\}$,
$C^\natural=A[E]=A^{>p}$, and
\begin{equation}\label{isg:cf:eq:cyclic}
 \cmp E(\gamma)=\rho\bigl(\tht(\gamma)\bigr)\quad(\gamma\in E),
 \qquad\text{that is,}\qquad
 \cmp E(p\cat\plus\cat u)=p\cat u,\quad
 \cmp E(x)=x\ \ (p\not\pre x).
\end{equation}
Consequently $\Phi_E(P_E(x))=\rho_*\bigl(\Rect(x)-\pi_p(x)\bigr)$ for every
$x\in A$, the realization $\mathcal F_E(A)$ of $A/\Z\varepsilon$ is the
group $B$ of \eqref{isg:eq:Bquotient}, and the splitting of
\cref{isg:cf:thm:split} for $C=\Z\varepsilon$ is that of
\cref{isg:lem:split,isg:thm:C}.
\end{proposition}
\begin{proof}
Since $\omega^p=2^n\varepsilon\in\Z\varepsilon$ and no other monomial
$\omega^\gamma$, $\gamma\in\Gamma$, lies in the convex subgroup
$\Z\varepsilon$ (its positive members are the multiples of
$c\omega^p$), $\Gamma_C=\{p\}$. By \cref{isg:cf:lem:arch}, an element of
$A$ lies in $A[E]$ exactly when its support avoids $p$, that is, when
$\pi_p(x)=0$; so $A[E]=A^{>p}$.

\emph{First route: the formula.} By \cref{isg:prop:bottom} every
$x\in E$ lies in the cone $C_\alpha$ and $x>p$, so, as in
\cref{isg:sec:surgery}, either $p\not\pre x$ or $x=p\cat\plus\cat u$.
Every proper prefix of $x$ lies in the initial class $\Gamma$, hence in
$E$ unless it equals $p$. If $p\not\pre x$, no prefix equals $p$, every
position is retained and $\cmp E(x)=x=\rho(\plus\cat x)=\rho(\tht(x))$.
If $x=p\cat\plus\cat u$, the only prefix equal to $p$ is
$x\restriction\alpha$, so exactly the plus sign at position $\alpha$ is
deleted and $\cmp E(x)=p\cat u=\rho(\plus\cat p\cat u)=\rho(\tht(x))$.
Positions are read as ordinal intervals, so limit lengths of $u$ need no
separate case.

\emph{Second route: uniqueness.} On $E$, the map $\rho\circ\tht$ is an
order-isomorphism (\cref{isg:lem:thetaorder}, and $\rho$ preserves
comparison), it preserves and reflects prefixes (\cref{isg:thm:tree} for
points above $p$, and \cref{isg:lem:rho}), and its image
$\rho[\tht[\Gamma]\setminus\{0\}]$ is initial
(\cref{isg:cor:initialimage,isg:lem:rho}). By the uniqueness in
\cref{isg:cf:prop:coherence} it equals $\cmp E$.

For $x=\sum_\gamma r_\gamma\omega^\gamma\in A$, \eqref{isg:eq:N} gives
$\Rect(x)-\pi_p(x)=\sum_{\gamma>p}r_\gamma\omega^{\tht(\gamma)}$, since the
constant term of $\Rect(x)$ is $c^{-1}r_p=\pi_p(x)$. Deleting the root of
each exponent gives $\sum_{\gamma\in E}r_\gamma\omega^{\rho\tht(\gamma)}$,
which is $\Phi_E(P_E(x))$ by \eqref{isg:cf:eq:cyclic}. An element of $H$
has zero constant coefficient exactly when it is the image of an element
of $A^{>p}$, so $H^+=\Rect[A^{>p}]$ and
$B=\rho_*[H^+]=\Phi_E[A[E]]=\mathcal F_E(A)$. The two splittings have the
same complement $A^{>p}=C^\natural$ and the same cyclic summand.
\end{proof}

Thus the construction of source~02, applied to $C=\Z\varepsilon$,
compresses away the bottom exponent and gives exactly source~01's
$\rho\circ\tht$ on exponents, with the same coefficients. The two
constructions also agree by the uniqueness of compression, and a finite
computation at placement confirmed \eqref{isg:cf:eq:cyclic} on words of
small length (\cref{isg:sec:provenance}).

\subsection{Limit stages of iterated bottom-layer removal}\label{isg:cf:sec:limit}
\Cref{isg:cor:iteration} (source~01) covers finite stages only. The
following corollary, formulated in the merge from source~02's
\cref{isg:cf:cor:closure,isg:cf:cor:coherentfactors} and
\cref{isg:cf:prop:cyclic}, covers every stage.

\begin{corollary}[Transfinite iteration; \textup{[merge]}]\label{isg:cf:cor:limit}
Let $A\subseteq\No$ be initial with exponent class $\Gamma$. Define
subgroups $C_\xi$ of $A$ by recursion: $C_0=\{0\}$; if $A/C_\xi$ is
nonzero and discrete, $C_{\xi+1}$ is the preimage in $A$ of the least
positive cyclic subgroup of $A/C_\xi$, and otherwise the recursion stops
at $\xi$; at a limit $\lambda$, $C_\lambda=\bigcup_{\xi<\lambda}C_\xi$.
Put $E_\xi=\Gamma\setminus\Gamma_{C_\xi}$. Then, at every stage $\xi$
reached, limit stages included:
\begin{enumerate}[label=\textup{(\roman*)}]
\item $C_\xi$ is convex, and $x+C_\xi\mapsto\Phi_{E_\xi}(P_{E_\xi}(x))$ is
      an ordered-group isomorphism of $A/C_\xi$ onto the initial group
      $\mathcal F_{E_\xi}(A)$;
\item for finite $\xi$ this is the realization produced by $\xi$
      applications of \cref{isg:thm:C}, as in \cref{isg:cor:iteration};
\item for $\xi<\eta$ with $E_\eta\ne\varnothing$, the realization at
      stage $\eta$ is obtained from that at stage $\xi$ by compressing
      the interval $\cmp{E_\xi}[E_\eta]$:
      $\Phi_{E_\eta}=\Phi_{\cmp{E_\xi}[E_\eta]}\circ\Phi_{E_\xi}$ on normal
      forms supported in $E_\eta$.
\end{enumerate}
No hypothesis beyond the initiality of $A$ is needed at a limit stage.
\end{corollary}
\begin{proof}
(i) The preimage of a convex subgroup under the order-preserving quotient
map by a convex subgroup is convex, and the union of a chain of convex
subgroups is a convex subgroup; so every $C_\xi$ is convex, and
\cref{isg:cf:prop:factorcoordinates,isg:cf:thm:split,isg:cf:thm:transport}
give the realization, with $\mathcal F_\varnothing(A)=\{0\}$ when $E_\xi$
is empty. (iii) is \cref{isg:cf:cor:coherentfactors} for
$E_\eta\subseteq E_\xi$.

(ii) By induction on finite $k$. For $k=0$, $E_0=\Gamma$ is initial, so
every position is retained and $\Phi_\Gamma$ is the identity. Suppose the
stage-$k$ realization is $B_k=\mathcal F_{E_k}(A)$, an initial group with
exponent class $\cmp{E_k}[E_k]$. The convex subgroup $C_{k+1}/C_k\cong\Z$
corresponds by \cref{isg:cf:prop:factorcoordinates} to the interval
$\Gamma_{C_{k+1}}\setminus\Gamma_{C_k}$, which is nonempty and cannot
contain two points (their monomials would be Archimedean inequivalent in a
cyclic group); so it is the single point $\min E_k$, and
$E_{k+1}=E_k\setminus\{\min E_k\}$. The least positive cyclic subgroup
of $B_k$ is the image of $C_{k+1}/C_k$, with bottom exponent
$\cmp{E_k}(\min E_k)$. By \cref{isg:cf:prop:cyclic}, applied to $B_k$,
\cref{isg:thm:C} realizes $B_k/(\text{that subgroup})$ by compressing
$\cmp{E_k}[E_k]$ minus its bottom point, which is $\cmp{E_k}[E_{k+1}]$.
By (iii) the result is $\mathcal F_{E_{k+1}}(A)$.
\end{proof}

This answers \cref{isg:q:limit}: the limit stage needs no hypothesis, its
quotient is realized directly from the cut $\Gamma_{C_\lambda}$ rather
than as a limit of earlier realizations, and it is coherent with every
earlier stage. The recursion itself uses no choice: each least positive
element is unique.

\section{The residue obstruction for Presburger groups}\label{isg:cf:sec:presburger}

\subsection{The bottom cyclic group and coherent residues}
Let $G$ be a $\Z$-group with least positive element $e$. The subgroup
$\Z e$ is convex: finite induction shows that the interval $[0,me]$
contains exactly $0,e,\ldots,me$ for every ordinary $m\in\N$ (as in the
proof of \cref{isg:cor:relative}). Its quotient is torsion free, as every
ordered abelian group is. In particular $\Z e$ is pure in $G$:
\[
 ny\in\Z e\quad\Longrightarrow\quad y\in\Z e.
\]

\begin{lemma}[Residues]\label{isg:cf:lem:residues}
For each ordinary $n\ge1$ and $x\in G$ there is a unique
$r_n(x)\in\{0,\ldots,n-1\}$ such that
$x-r_n(x)e\in nG$. The resulting homomorphisms
$G\to\Z/n\Z$ are compatible and define $\res_G$.
Moreover
\[
 \ker\res_G=\bigcap_{n\ge1}nG.
\]
\end{lemma}
\begin{proof}
Divisibility of $G/\Z e$ gives $x=ny+ke$ for some $y\in G$ and
$k\in\Z$. Reduce $k$ modulo $n$. If two remainders worked, then
$(r-r')e=ny$. Purity puts $y$ in $\Z e$, so $n$ divides $r-r'$ and
the chosen range forces equality. Addition and compatibility under
divisibility of moduli follow from the same uniqueness. The kernel
identity is immediate from the zero remainder condition.
\end{proof}

\begin{lemma}[The divisible kernel, classical]\label{isg:cf:lem:divkernel}
The subgroup $G^{\dv}:=\ker\res_G$ is divisible and is the largest
divisible subgroup of $G$.
\end{lemma}
\begin{proof}
Let $a\in G^{\dv}$ and fix $k\ge1$. Choose $b\in G$ with $kb=a$.
For every $n\ge1$, the membership $a\in knG$ gives $c\in G$ such that
$a=knc$. Torsion freeness yields $b=nc$, so $b\in nG$ for all $n$.
Thus $b\in G^{\dv}$. Any divisible subgroup is contained in every $nG$,
and therefore in $G^{\dv}$.
\end{proof}

This is the divisible-kernel part of
\cite[Proposition~4.1 and Remark~4.2]{Jerabek}; the elementary proof is
included to make the subsequent equivalences explicit.

\begin{proposition}[Split bottom versus ordinary residues]\label{isg:cf:prop:residuesplit}
For a $\Z$-group $G$, the following are equivalent:
\begin{enumerate}[label=\textup{(\roman*)}]
\item $\Z e$ is a direct summand of $G$;
\item $\res_G(G)=\Z\subseteq\Zhat$;
\item $G=G^{\dv}\oplus\Z e$.
\end{enumerate}
In case \textup{(iii)}, the order is $G^{\dv}\lex\Z$. Every group
complement to $\Z e$ equals $G^{\dv}$.
\end{proposition}
\begin{proof}
Assume $G=V\oplus\Z e$. The quotient map identifies $V$ with
$G/\Z e$, so $V$ is divisible and $V\subseteq G^{\dv}$. For $v+me$
its coherent residue is $m$. Hence $\res_G(G)=\Z$, and its kernel is
exactly $V$.

Conversely, if every coherent residue is an ordinary integer, then for
each $x$ there is $m\in\Z$ such that $\res_G(x)=m$.
It follows that $x-me\in G^{\dv}$. The intersection
$G^{\dv}\cap\Z e$ is zero, since no nonzero ordinary integer is divisible
by every positive ordinary integer. Thus (iii) holds, and implies (i).
Convexity of $\Z e$ gives the lexicographic ordering for any complement:
a nonzero positive quotient lift is larger than every member of
$\Z e$. The first argument also proves uniqueness of that complement.
\end{proof}

\subsection{Initial omnific realization in the split case}
Recall $\Oz=\Z\oplus\Pi$ from \eqref{isg:eq:Oz} and
\cref{isg:sec:conventions}. The symbol $\Pi$ denotes purely
\emph{infinite} forms, with zero included, not positive elements and not
infinitesimals. Source~02 uses the integer-bottom suspension
$J(B)=\Z\oplus\sigma_*[B]$ of \cref{isg:thm:suspension}: for initial
$B\subseteq\No$, $J(B)$ is an initial subgroup of $\Oz$ order-isomorphic
to $B\lex\Z$. It re-proves this from \cref{isg:fact:EK} by the argument of
the first paragraph of the proof of \cref{isg:thm:suspension}, so that
\cref{isg:cf:main:presburger} does not depend on source~01's separate
normalization proof. The statement and proof are printed once, there.

\begin{proof}[Proof of \cref{isg:cf:main:presburger}]
If $G\in\IR$, convex splitting (\cref{isg:cf:thm:split}) applied to
$\Z e$ in an initial realization makes the bottom cyclic group a direct
summand. \Cref{isg:cf:prop:residuesplit} proves the algebraic and residue
equivalences. Conversely, in the split case $G^{\dv}$ is divisible and
ordered. Choose an initial realization $B$ of it by
\cref{isg:cf:fact:ehrlich}. The suspension $J(B)$ is an initial omnific
realization of $G$. An initial subgroup of $\Oz$ is, by the meaning of
initial here, an initial subgroup of $\No$.
\end{proof}

Source~01's relative criterion gives the same answer by a second route.
For a $\Z$-group the quotient $G/\Z e$ is divisible, hence initially
realizable by \cref{isg:cf:fact:ehrlich}, so \cref{isg:cor:relative}
reduces initial realizability of $G$ to the splitting of $\Z e$: the
first and third conditions of \cref{isg:cf:main:presburger}.

\begin{corollary}[Rigidity on the initial side]\label{isg:cf:cor:rigidityinitial}
The only initially realizable rigid $\Z$-group is $\Z$ itself.
\end{corollary}
\begin{proof}
In $G^{\dv}\lex\Z$, the map $(d,m)\mapsto(2d,m)$ is an ordered-group
automorphism, because $G^{\dv}$ is divisible and torsion free. It is
nontrivial unless $G^{\dv}=0$. The standard group $\Z$ has no nonidentity
order-preserving group automorphism.
\end{proof}

\section{An explicit factorial family of nonsplit models}\label{isg:cf:sec:factorial}

\subsection{Presentation inside an ordinary rational plane}
Factorial moduli are cofinal for divisibility, so a profinite integer
$\zeta\in\Zhat$ is specified by its compatible residues modulo $n!$.
Let $S_n=S_n(\zeta)$ be the unique integer with
\[
 0\le S_n<n!,\qquad S_n\equiv\zeta\pmod{n!},\qquad n\ge1.
\]
Then $S_1=0$ and
\begin{equation}\label{isg:cf:eq:digits}
 S_{n+1}=S_n+a_n n!,\qquad a_n\in\{0,1,\ldots,n\}.
\end{equation}
In the lexicographically ordered rational plane $\Q\lex\Q$, put
\begin{equation}\label{isg:cf:eq:groupfamily}
 e=(0,1),\qquad
 x_n=\left(\frac1{n!},\frac{S_n}{n!}\right),\qquad
 G_\zeta=\bigcup_{n\ge1}(\Z x_n\oplus\Z e).
\end{equation}
The groups in the union are nested because
\begin{equation}\label{isg:cf:eq:transition}
 (n+1)x_{n+1}=x_n+a_n e.
\end{equation}

\begin{proposition}[Structure of the factorial family]\label{isg:cf:prop:family}
The group $G_\zeta$ is countable, has least positive element $e$, and
first-coordinate projection $\pi_1$ gives an exact sequence of ordered
groups
\begin{equation}\label{isg:cf:eq:extension}
 0\longrightarrow\Z e\longrightarrow G_\zeta
   \xrightarrow{\ \pi_1\ }\Q\longrightarrow0.
\end{equation}
Consequently it is a $\Z$-group. Its first generator has residue
$\res_{G_\zeta}(x_1)=-\zeta$.
\end{proposition}
\begin{proof}
Each stage is a free rank-two group, and the union is countable. Its
first-coordinate image is $\bigcup_n(n!)^{-1}\Z=\Q$.
If $kx_n+me$ has first coordinate zero, then $k=0$; hence the kernel is
exactly $\Z e$. A positive element of nonzero first coordinate exceeds
every integer multiple of $e$, and the positive elements of the kernel
are $e,2e,\ldots$. This proves discreteness and convexity.
The quotient is $\Q$, hence divisible. Finally,
$n!x_n=x_1+S_ne$, so $x_1\equiv-S_n e\pmod{n!G_\zeta}$.
Compatibility over factorial moduli gives the residue assertion.
\end{proof}

\begin{theorem}[Exact splitting test]\label{isg:cf:thm:familytest}
For $\zeta\in\Zhat$,
\[
 G_\zeta\in\IR\quad\Longleftrightarrow\quad
 \text{\eqref{isg:cf:eq:extension} splits}\quad\Longleftrightarrow\quad
 \zeta\in\Z.
\]
\end{theorem}
\begin{proof}
A retraction $r:G_\zeta\to\Z$ with $r(e)=1$ is determined by the
integers $b_n=r(x_n)$. The relations imply
\begin{equation}\label{isg:cf:eq:retraction}
 (n+1)b_{n+1}=b_n+a_n,
 \qquad n!b_n=b_1+S_n.
\end{equation}
Thus $b_1+\zeta$ is zero modulo every $n!$, so $\zeta=-b_1\in\Z$.
Conversely, when $\zeta=m\in\Z$, put
$b_n=(S_n-m)/n!$. These are integers satisfying all transition relations.
On each stage define $r(kx_n+le)=kb_n+l$; compatibility makes this a
well-defined retraction on the union. The equivalence with initial
realizability follows from \cref{isg:cf:main:presburger}.
\end{proof}

\subsection{A completely explicit noninteger parameter}
Define
\begin{equation}\label{isg:cf:eq:explicitparameter}
 \zeta_* =\sum_{k=1}^{\infty}k!\quad\text{in }\Zhat,
 \qquad S_n=\sum_{k=1}^{n-1}k!.
\end{equation}
For each fixed modulus, all sufficiently late factorials vanish, so
this is a well-defined profinite sum. It is not an ordinary convergent
real series and is not a Hahn sum of omnific integers.

\begin{lemma}\label{isg:cf:lem:noninteger}
One has $0\le S_n<n!$, $S_n\to+\infty$, and $S_n/n!\to0$ in the
ordinary real sense. In particular $\zeta_*\notin\Z$.
\end{lemma}
\begin{proof}
The first inequality is direct for $n=1,2$. If $S_n<n!$, then
$S_{n+1}=S_n+n!<2n!\le(n+1)!$ for $n\ge1$, with the small equality
case harmless because the first inequality is strict.
For $n\ge3$, induction gives
$S_n\le2(n-1)!$: it holds at $n=3$, and
\[
 S_{n+1}\le2(n-1)!+n!\le2n!\qquad(n\ge2).
\]
Hence $S_n/n!\le2/n$. Unboundedness follows from the summand
$(n-1)!$.
If $\zeta_*$ were a nonnegative ordinary integer $m$, its least
residues modulo all sufficiently large $n!$ would be the fixed value
$m$, contrary to unboundedness. If it were $-m$ with $m>0$, those
residues would be $n!-m$, whose ratio to $n!$ tends to one, contrary to
the proved ratio limit.
\end{proof}

For this parameter all digits $a_n$ equal $1$, and $(n+1)x_{n+1}=x_n+e$.
The first generators are shown below. Thus the abstract obstruction has a
presentation entirely in terms of pairs of ordinary rational numbers.
\begin{center}\small
\begin{tabular}{rrrr}
\toprule
$n$ & $n!$ & $S_n$ & $x_n$\\
\midrule
1&1&0&$(1,0)$\\
2&2&1&$(1/2,1/2)$\\
3&6&3&$(1/6,1/2)$\\
4&24&9&$(1/24,3/8)$\\
5&120&33&$(1/120,11/40)$\\
6&720&153&$(1/720,17/80)$\\
7&5040&873&$(1/5040,97/560)$\\
\bottomrule
\end{tabular}
\end{center}

\begin{remark}[Every finite retraction problem is solvable]\label{isg:cf:rem:finite}
Fix $N$. Taking $b_1=-S_N$ makes
$b_n=(S_n-S_N)/n!$ integral for every $n\le N$ and satisfies every
transition visible through stage $N$. There is nevertheless no single
integer $b_1$ satisfying all stages. The obstruction is precisely the
failure of this infinite compatibility problem to have a solution in
$\Z$, although it has the profinite solution $-\zeta_*$.
\end{remark}

\subsection{Extension classes and rigidity}
\begin{proposition}[Classical extension parameter, explicit form]\label{isg:cf:prop:extensionclasses}
There is an isomorphism $G_\zeta\to G_\eta$ inducing the identity on
both endpoints of \eqref{isg:cf:eq:extension} if and only if
$\zeta-\eta\in\Z$.
\end{proposition}
\begin{proof}
Such a map must send $e$ to $e$ and $x_n^\zeta$ to
$x_n^\eta+b_ne$ with $b_n\in\Z$. Comparing $n!$ times these elements
gives
\[
 n!b_n=b_1+S_n(\zeta)-S_n(\eta).
\]
Thus $b_1+\zeta-\eta=0$ in $\Zhat$. Conversely, if this equality
holds for an ordinary integer $b_1$, the displayed formula defines
integers $b_n$ satisfying the transition relations. They define compatible
isomorphisms on the stages and hence an isomorphism on the union;
replacing $b_n$ by $-b_n$ gives the inverse.
\end{proof}

This is the concrete presentation behind the familiar parameter
$\Zhat/\Z$. No new computation of the functor $\Ext$ is being claimed.
For completeness, arbitrary extensions of $\Q$ by $\Z$ admit the same
kind of presentation: choose lifts of $1/n!$, adjust their differences
by integers, and collect the resulting compatible residues. Only the
explicit family and its equivalence relation are needed here.

\begin{theorem}[Continuum many countable rigid countermodels]\label{isg:cf:thm:continuum}
There are $2^{\aleph_0}$ pairwise nonisomorphic countable rigid
$\Z$-groups not in $\IR$, among the $G_\zeta$.
\end{theorem}
\begin{proof}
The profinite integers have cardinality $2^{\aleph_0}$, and each coset
modulo $\Z$ is countable. \Cref{isg:cf:prop:extensionclasses} therefore
gives continuum many endpoint-marked extension classes. An ordered-group
isomorphism must preserve the least positive element and induces an
order-preserving automorphism of $\Q$. Every such automorphism is
multiplication by an element of $\Q_{>0}$, a countable set of choices.
Thus one ordered-group isomorphism class can account for at most
countably many of the marked extension classes. There remain continuum
many ordered-group types, with only the zero extension class split.

Let $\zeta\notin\Z$. If $G_\zeta^{\dv}$ were nonzero, its image in
$\Q$ would be a nonzero divisible subgroup, hence all of $\Q$.
The intersection $G_\zeta^{\dv}\cap\Z e$ is zero. Consequently the
restriction of $\pi_1$ would be an isomorphism $G_\zeta^{\dv}\to\Q$,
providing a section and contradicting \cref{isg:cf:thm:familytest}.
Hence $G_\zeta^{\dv}=0$ and $\res_{G_\zeta}$ is injective.
Every ordered-group automorphism preserves $e$ and all its residues,
so injectivity of the residue map forces that automorphism to be the
identity.
\end{proof}

The final argument is the residue-injectivity rigidity mechanism studied
in \cite{Jerabek}. Here it also supplies a particularly explicit family
on the noninitial side of \cref{isg:cf:main:presburger}.

\section{Elementary equivalence, dense models, and omitted types}\label{isg:cf:sec:modeltheory}

\subsection{Initial realizability is not elementary}
\begin{corollary}\label{isg:cf:cor:nonelementary}
Initial realizability is not preserved by elementary equivalence, even
among countable $\Z$-groups. It is not first-order axiomatizable as a
property of ordered abelian groups.
\end{corollary}
\begin{proof}
The standard group $\Z$ is initial. Every $G_\zeta$ is a $\Z$-group,
hence elementarily equivalent to $\Z$ by the characterization of
Presburger models \cite{Jerabek}. For $\zeta\notin\Z$ it is not
initially realizable. A first-order axiomatizable class is closed under
elementary equivalence.
\end{proof}

The same contrast already occurs with the explicit initial nonstandard
$\Z$-group $\Q\omega+\Z$, order-isomorphic to $\Q\lex\Z$.
It has the same Presburger theory as every nonsplit $G_\zeta$.

\begin{theorem}[Dense elementary-equivalence counterexample]\label{isg:cf:thm:dense}
For $\zeta\notin\Z$, let $\widetilde G_\zeta=G_\zeta\lex\Q$.
Then $\widetilde G_\zeta$ is countable and densely ordered, is not
initially realizable, and is elementarily equivalent as an ordered group
to the initial subgroup
\[
 W=\Q\omega^2+\Z\omega+\Q\subseteq\No.
\]
\end{theorem}
\begin{proof}
The final $\Q$ layer is convex and densely ordered. Between two elements
with distinct upper coordinates one can increase the lower coordinate
of the smaller element; with equal upper coordinates one uses density
of $\Q$. Thus the whole group is densely ordered.
The quotient by the bottom copy of $\Q$ is $G_\zeta$. If
$\widetilde G_\zeta$ were initially realizable, convex-factor closure
would make $G_\zeta$ initially realizable, a contradiction.

Since $G_\zeta\equiv\Q\lex\Z$, adjoining a disjoint second sort
interpreted as the fixed ordered group $\Q$ preserves elementary
equivalence. This can be verified by induction on formulas: formulas
in disjoint sorts reduce to Boolean combinations of formulas in the
individual sorts. The Cartesian product with coordinatewise addition
and lexicographic order is uniformly interpretable in that two-sorted
structure. It follows that
\[
 \widetilde G_\zeta\equiv(\Q\lex\Z)\lex\Q\cong W.
\]
Finally, the exponent class $\{0,1,2\}$ is initial and has no right
ancestors internally, and the coefficient groups $\Q,\Z,\Q$ are
initial real groups. Finite support gives truncation closure and
cross-sectionality, so the criterion proves that $W$ is initial.
\end{proof}

\subsection{A single recursive type witnessing the obstruction}
With the computable integers $S_n$ of \eqref{isg:cf:eq:explicitparameter},
consider the partial type in one variable
\begin{equation}\label{isg:cf:eq:type}
 \tau_*(X)=\{X\equiv S_ne\pmod{n!G}:n=1,2,\ldots\}.
\end{equation}
Each congruence abbreviates the first-order formula
$\exists Y\ (X=n!Y+S_ne)$ in the ordered-group language with $e$ named.
The least positive element is itself definable, so naming it does not
change the substantive distinction.

\begin{theorem}[An omitted computable type]\label{isg:cf:thm:type}
The type $\tau_*$ is finitely satisfiable in every $\Z$-group and omitted
by every initially realizable $\Z$-group. It is a recursive type under
a fixed standard coding of these canonical formulas. Consequently no
recursively saturated $\Z$-group is initially realizable.
\end{theorem}
\begin{proof}
A finite set of congruences is satisfied by $S_Ne$ for any sufficiently
large $N$, because $S_N\equiv S_n\pmod{n!}$ for $n\le N$.
A realization of all congruences has coherent residue $\zeta_*$.
That residue is not in $\Z$, by \cref{isg:cf:lem:noninteger}, whereas
\cref{isg:cf:main:presburger} requires the entire residue image to be
$\Z$. Thus the type is omitted in every initially realizable $\Z$-group.

The sequences $n!$ and $S_n$ are computable. Using canonical integer
numerals, the set of displayed formulas is decidable: read the modulus,
decide whether it is a factorial, and check the corresponding residue,
handling the duplicate modulus $1$ at the start by our convention
$n\ge1$. Standard definitions of recursive saturation therefore apply.
A recursively saturated model realizes every finitely satisfiable
recursive type over finitely many parameters, contradicting the omission.
\end{proof}

\begin{corollary}[Countable saturation]\label{isg:cf:cor:saturation}
If a $\Z$-group is countably saturated, its residue map is surjective
onto $\Zhat$. In particular it is not initially realizable.
\end{corollary}
\begin{proof}
For an arbitrary $\zeta\in\Zhat$, use its compatible residues
$S_n(\zeta)$ instead of $S_n$ in \eqref{isg:cf:eq:type}. The resulting
countable type is finitely satisfiable by ordinary multiples of $e$.
Countable saturation supplies a realization. Taking $\zeta_*$ yields
the final assertion without any assumption on the continuum hypothesis.
\end{proof}

\begin{theorem}[Every nonprincipal sequential ultrapower fails]\label{isg:cf:thm:ultrapower}
Let $G$ be any $\Z$-group and $U$ any nonprincipal ultrafilter on $\N$.
Then $G^\N/U$ is not initially realizable.
\end{theorem}
\begin{proof}
By \L o\'s's theorem it is a $\Z$-group. In the ultrapower take
$a=[(S_Ne)_{N\ge1}]_U$. For any fixed $n$, the congruence
$S_Ne\equiv S_ne\pmod{n!G}$ holds for all $N\ge n$, hence on a set
in $U$. Therefore $a$ realizes $\tau_*$. The previous theorem rules out
an initial realization. The argument uses only cofinite sets and does
not invoke an ultrapower saturation theorem.
\end{proof}

\subsection{An obstruction for nonstandard Peano models}\label{isg:cf:sec:PA}
\begin{theorem}[Nonstandard arithmetic cannot have an initial additive realization]
\label{isg:cf:thm:PA}
Let $M$ be a nonstandard model of Peano arithmetic, and let $G(M)$ be
the ordered abelian group completion of its additive monoid. Then
$G(M)$ is not initially realizable in $\No$.
\end{theorem}
\begin{proof}
The ordinary division algorithm, as proved in Peano arithmetic, gives
unique division with remainder by every positive standard integer.
Consequently $G(M)$ is a $\Z$-group with least positive element $1$.

Choose a nonstandard $h\in M$. Peano arithmetic defines the factorial
function and the primitive-recursive partial sums
\[
 F(h)=\sum_{k=1}^{h}k!.
\]
The sum here is the internally coded finite sum in $M$, not an external
infinite group sum. For each fixed ordinary $n$, arithmetic proves
that $n!$ divides $k!$ whenever $k\ge n$, and therefore proves
\[
 h\ge n\quad\Longrightarrow\quad F(h)\equiv S_n\pmod{n!}.
\]
This follows by induction on the finite-sum length inside the model;
all the displayed factorial and summation recursions are provably
total in Peano arithmetic. Since $h$ exceeds every standard $n$,
$F(h)$ realizes the external type $\tau_*$ in $G(M)$.
\Cref{isg:cf:thm:type} excludes initial realizability.
\end{proof}

The arithmetic infrastructure used here is standard; factorials and
division can in fact be handled in much weaker arithmetic, as noted in
\cite{EnayatHamkinsWcislo}. The fragment is not optimized here. The
conclusion is stronger than merely observing that the usual
omnific multiplication does not make the positive cone of $\Oz$ a
nonstandard Peano model: it excludes \emph{any} such model having an
initially realizable ordered additive group.

\begin{proof}[Proof of \cref{isg:cf:main:models}]
The discrete family and rigidity are \cref{isg:cf:thm:continuum} and
Presburger elementary equivalence. The dense examples are
\cref{isg:cf:thm:dense}. The remaining assertions are
\cref{isg:cf:thm:type,isg:cf:thm:ultrapower,isg:cf:thm:PA}.
\end{proof}

\section{Omnific and surcomplex consequences}\label{isg:cf:sec:complex}

\subsection{The precise residue image of omnific integers}
\begin{proposition}\label{isg:cf:prop:omnificresidue}
The additive group of $\Oz$ has least positive element $1$, its purely
infinite subgroup $\Pi$ is divisible, and for every ordinary $n\ge1$
\[
 \Oz/n\Oz\cong\Z/n\Z,\qquad
 \res_{\Oz}(z)=\ct(z)\in\Z\subseteq\Zhat.
\]
The type $\tau_*$ is finitely satisfiable in $\Oz$ and has no realization
there.
\end{proposition}
\begin{proof}
The normal form \eqref{isg:eq:Oz} gives $z=y+m$ uniquely with
$y\in\Pi$ and $m\in\Z$. Every nonzero purely infinite form dominates
all real constants in absolute value, proving the least-positive claim.
For a positive ordinary $n$, dividing the real coefficients of $y$ by
$n$ leaves its support strictly positive. Thus $\Pi$ is divisible and
$z\equiv m\pmod{n\Oz}$. The kernel of the residue map to $\Z/n\Z$
is exactly $n\Oz$, proving both formulas. The final assertion follows
by the same explicit residue argument as \cref{isg:cf:thm:type}.
\end{proof}

This is a class-function statement; it does not require a set of all
omnific integers or a set quotient by a proper-class subgroup. The
finite quotient is represented directly by its ordinary remainder.

\begin{corollary}[Unit-preserving omnific embeddings]\label{isg:cf:cor:unit}
For a set-sized $\Z$-group $G$ with least positive element $e$, the
following are equivalent: $G\in\IR$; there is an order-preserving
additive embedding $f:G\to\Oz$ with $f(e)=1$; there is any additive
homomorphism $f:G\to\Oz$ with $f(e)=1$.
\end{corollary}
\begin{proof}
An initial omnific realization supplies the stated embedding and takes
the least positive element to $1$. Conversely, for any such homomorphism,
$\ct\circ f:G\to\Z$ is a retraction on $\Z e$. Thus the bottom cyclic
group splits, and \cref{isg:cf:main:presburger} applies.
\end{proof}

In particular, the ordered additive completion of a nonstandard Peano
model cannot even admit a unit-preserving additive embedding into
$\Oz$, whether or not the image is initial. There is no contradiction
with ordinary ordered-group embeddings that move the distinguished
unit: the copy $(a,b)\mapsto a\omega^2+b\omega$ embeds each
$G_\zeta$ into $\Oz$ but sends $e$ to $\omega$, not to $1$.

\subsection{Period groups of surcomplex phases}\label{isg:cf:sec:periods}
The trigonometry report studies the normalized period group
$\mathcal I(\Psi)$ of a phase $\Psi$ on a field $F$
(\replabel{trigonometry:per:def:periodgroup}). It is an additive integer
part of $F$ with least positive element $1$, it sits in the exact
sequence $0\to\Z\to\mathcal I\to\Pi_F\to0$ given by the principal part
(\replabel{trigonometry:per:thm:dictionary}), and for set-sized $F$ it is
a $\Z$-group (after \replabel{trigonometry:per:prop:division}). Its
residue map $\res_{\mathcal I}$ is the map $\res_G$ of
\cref{isg:cf:lem:residues} for $G=\mathcal I$, $e=1$
(\replabel{trigonometry:per:eq:residue}), and its profinite defect
$\partial_\Psi:\Pi_F\to\Zhat/\Z$
(\replabel{trigonometry:per:thm:defect}) is that residue taken on a lift
and read modulo $\Z$. Source~02's criterion therefore decides which
period groups are initially realizable. The statement is the merge's.

\begin{corollary}[Initially realizable period groups; \textup{[merge]}]\label{isg:cf:cor:periods}
Let $F$ be $\No$ or one of the set-sized fields of
\replabel{trigonometry:per:sec:phases}, and let $\Psi$ be a phase on $F$
with normalized period group $\mathcal I=\mathcal I(\Psi)$, ordered as a
subgroup of $F$. The following are equivalent:
\begin{enumerate}[label=\textup{(\roman*)}]
\item $\mathcal I$ is initially realizable;
\item the profinite defect $\partial_\Psi$ vanishes;
\item the period sequence $0\to\Z\to\mathcal I\to\Pi_F\to0$ splits.
\end{enumerate}
In that case $\mathcal I\cong\Pi_F\lex\Z$ as ordered groups; for
$F=\No$ and for the initial fields $\No(\lambda)$ this is an
order-isomorphism onto the canonical period group $\Oz\cap F$, which is
initial.
\end{corollary}
\begin{proof}
(ii)$\Leftrightarrow$(iii) is \replabel{trigonometry:per:thm:split},
conditions (1) and (2).

(i)$\Rightarrow$(iii). The element $1$ is the least positive element of
$\mathcal I$ and $\mathcal I\cap\mathcal O_F=\Z$. An element of
$\mathcal I$ outside $\Z$ is infinite, so its sign is the sign of its
principal part; hence the principal part induces an ordered isomorphism
$\mathcal I/\Z\cong\Pi_F$. If $\mathcal I$ is initially realizable, the
convex subgroup $\Z$ is a direct summand by \cref{isg:cf:thm:split}
(or by \cref{isg:cor:relative}), and the period sequence splits.

(iii)$\Rightarrow$(i). A splitting gives $\mathcal I=V\oplus\Z$ with
$V$ mapped isomorphically onto $\Pi_F$ by the principal part. By
convexity of $\Z$ the order is lexicographic, and since the principal
part preserves signs on $V$, the map $v+m\mapsto\mathrm{pr}(v)+m$ is an
order-isomorphism $\mathcal I\to\Pi_F\oplus\Z$, the latter ordered as a
subgroup of $F$, which is $\Pi_F\lex\Z$. For $F=\No$ the target is
$\Pi\oplus\Z=\Oz$, and for $F=\No(\lambda)$ it is $\Oz\cap F$; both are
initial ($\Oz$ by \eqref{isg:eq:Oz}, and an intersection of initial
classes is initial). For a set-sized $F$, $\Pi_F$ is a divisible ordered
group, so it has an initial realization $B$ by
\cref{isg:cf:fact:ehrlich}, and $J(B)$ realizes $\Pi_F\lex\Z$
(\cref{isg:thm:suspension}).

Second route for set-sized $F$: $\mathcal I$ is a set-sized $\Z$-group
with $\res_{\mathcal I}(\mathcal I)=q^{-1}(\partial_\Psi(\Pi_F))$, $q$
the quotient map $\Zhat\to\Zhat/\Z$
(\replabel{trigonometry:per:thm:defect}); this is $\Z$ exactly when
$\partial_\Psi=0$, and \cref{isg:cf:main:presburger} gives
(i)$\Leftrightarrow$(ii).
\end{proof}

So a period group is initially realizable exactly when its profinite
defect is zero. The corollary is additive: it says nothing about
multiplicative closure of the periods, which the trigonometry report
shows to be independent of the defect
(\replabel{trigonometry:per:sec:tests}), and it does not make a
noncanonical split period group equal to $\Oz\cap F$ as a subclass. The
nonsplit period groups are further instances of the obstruction of
\cref{isg:cf:main:presburger}.

\subsection{Rectangular surcomplex factors}
Let $\No[i]$, $A[i]$ and coordinatewise initiality be as in
\cref{isg:sec:complex} (\cref{isg:def:rectangular}). This is an
explicitly coordinate-dependent notion; no single canonical sign-tree
order on $\No[i]$ is postulated.

\begin{theorem}[Rectangular convex-factor transport]\label{isg:cf:thm:complex}
Let $G\subseteq\No$ be initial and $C\subseteq C'$ convex in $G$.
For the interval $E=\Gamma_{C'}\setminus\Gamma_C$, the map
\[
 a+ib\longmapsto\Phi_E(P_Ea)+i\Phi_E(P_Eb),
 \qquad a,b\in C',
\]
induces a $\Z[i]$-linear isomorphism
\[
 C'[i]/C[i]\ \cong\ \mathcal F_E(G)[i].
\]
It commutes with conjugation. The target is coordinatewise initial, and
the ambient complex-coefficient Hahn transport preserves and reflects
set-indexed strong summability.
\end{theorem}
\begin{proof}
The real factor map has kernel $C$ in each coordinate. The induced
map is therefore bijective on the rectangular quotient. Additivity and
compatibility with $(a,b)\mapsto(-b,a)$ give $\Z[i]$-linearity, and
negating the imaginary coordinate gives conjugation compatibility.
Initiality of the real target is \cref{isg:cf:thm:transport}. A complex
normal form is obtained by the union of the two real supports, still
reverse well ordered. Reindexing changes neither that order nor the
finite set of contributors at any exponent, so the strong-sum proof is
the same coefficientwise proof as in the real case.
\end{proof}

\begin{proposition}[Gaussian omnific residues]\label{isg:cf:prop:gaussian}
For each ordinary $n\ge1$,
\[
 \Oz[i]/n\Oz[i]\cong\Z[i]/n\Z[i],\qquad
 \res_{\Oz[i]}(a+ib)=\ct(a)+i\ct(b)\in\Z[i]\subseteq\Zhat[i].
\]
\end{proposition}
\begin{proof}
Apply \cref{isg:cf:prop:omnificresidue} to both coordinates.
Multiplication by $i$ interchanges them with a sign, so the coordinate
residue map is $\Z[i]$-linear. Its kernel modulo $n$ is $n\Oz[i]$.
\end{proof}

These are additive and module-theoretic conclusions. The exponent
compression need not respect addition of exponents, so it does not in
general preserve products of normal forms. No field, ring, differential,
or exponential equivalence is inferred from \cref{isg:cf:thm:complex}.

\section{Boundaries of the convex-factor theory}\label{isg:cf:sec:boundaries}

\subsection{Why these conclusions do not classify every ordered group}
Convex splitting is necessary for initial realizability, but no
sufficiency theorem for that condition alone is proved here. Even a
truncation-closed cross-sectional Hahn presentation must meet the
coefficient and right-ancestor conditions. The proof of
\cref{isg:cf:main:presburger} succeeds because the remaining quotient is
divisible and hence has an initial realization by an established theorem.
For a general discrete group that quotient need not be divisible.

Arbitrary subgroup closure is false. In fact the rational-coordinate
map $(a,b)\mapsto a\omega+b$ embeds every $G_\zeta$ into the initial
group $\Q\omega+\Q$. For $\zeta\notin\Z$ this gives a concrete
noninitially-realizable subgroup of a two-level initial group.
The structural theorem requires convex factors or separately verified
coordinate restrictions.
Likewise, a convex subgroup may fail literal initiality in its original
position. In the initial group $\D+\Z\omega^{-1}$ of
\cref{isg:sec:examples}, the subgroup $\Z\omega^{-1}$ is convex and does
not contain the simpler number $1$. It nevertheless has the initial
realization $\Z$, as the theorem predicts.

\subsection{The full optimality question is still separate}\label{isg:cf:sec:optimality}
Ehrlich and Kaplan prove an optimality theorem allowing class models,
for theories between ordered abelian groups and divisible ordered
abelian groups, and ask for a set-model counterpart (Section~8 of
\cite{EK}: Theorem~6 and Question~1; in the journal version,
Theorem~8.1 and Question~8.1 \cite{EKjournal}). The explicit dense groups
of \cref{isg:cf:thm:dense} are set-sized noninitial examples, but are not
a construction uniform over every proper intermediate theory.

In particular, Presburger arithmetic asserts discreteness and is not a
subtheory of divisible ordered abelian groups. Thus
\cref{isg:cf:main:presburger,isg:cf:thm:continuum} alone cannot answer
their full question by simply substituting Presburger arithmetic for an
intermediate theory. The dense variant in \cref{isg:cf:thm:dense} removes
discreteness of the whole group, but still does not establish the
required uniformity over all intermediate theories. Source~02 makes no
claim to have resolved that published question in its entirety.

Source~02 answers none of Ehrlich and Kaplan's questions. Question~2
(journal 9.1) is the subject of source~01 (\cref{isg:cor:answer}), and
source~02 does not use that answer: from source~01 it takes only the
suspension, which it re-proves. Question~3 (journal 9.2), the
initial-monoid problem, is not touched. Question~1 (journal 8.1) is the
one discussed here.

\begin{remark}[The two formulations of the optimality theorem; \textup{[merge]}]\label{isg:cf:rem:formulations}
The arXiv version quantifies over theories $T$ with
$T_{\mathrm{OAG}}\subseteq T\subseteq T_{\mathrm{DOAG}}$; since
$\mathrm{Th}(W)$ proves non-divisibility, it is not such a theory. The
journal version, read at placement in the authors' text hosted on the
second author's web page, states Theorem~8.1(i) for theories in the
language $\{<,+,0\}$ \emph{containing} the theory $T_D$ of nontrivial
densely ordered abelian groups, and Question~8.1 asks whether its ``only
if'' part can be established in NBG with set models alone. In that
formulation $\mathrm{Th}(W)$ is in scope, and \cref{isg:cf:thm:dense}
gives, for every theory $T\supseteq T_D$ in that language that holds in
$W$, a countable model of $T$ that is not initially realizable. This
settles the set-model ``only if'' part for those theories only, as
Ehrlich and Kaplan's own example of a $p$-divisible group does for
others; it gives no construction for every theory lacking full
divisibility, and the question stays open in both versions.
\end{remark}

\subsection{Product and completion pitfalls}
The group $G_\zeta$ for $\zeta\notin\Z$ has both convex kernel
$\Z$ and quotient $\Q$ initially realizable, but is not itself initially
realizable. Thus closure under convex factors has no unrestricted
converse for extensions. The missing datum is the extension class.

By contrast, \cref{isg:cf:thm:split} gives an actual complement whenever
the middle group is initial. This complement depends on the normal-form
realization; a group automorphism can change a splitting in general.
Only the Presburger split case has the uniqueness proved in
\cref{isg:cf:prop:residuesplit}.

The maps $P_E$ for infinitely many disjoint intervals may form an
ambient strongly summable family when applied to one fixed normal form,
but this does not identify an arbitrary initial group with the full
Hahn product of its factors. Such an identification would reintroduce
the unsupported completion assumption excluded in
\cref{isg:cf:rem:membership}.

\section{Verification, provenance, and a formalization route}\label{isg:sec:audit}

Each subsection below first gives source~01's record, as delivered, and
then source~02's; ``this manuscript'' in source~01's text means source~01.

\subsection{The exact dependency chain}
The proof separates into three kinds of input. Standard surreal
normal forms identify the ambient additive spaces. The established
Ehrlich--Kaplan criterion decides when a specified canonical Hahn image
is initial. Everything connecting the original discrete group to the
normalized image is then given by the explicit formulas and proofs in
this manuscript.

The essential chain is
\begin{gather*}
 \text{bottom coefficient structure}
 \ \longrightarrow\ \text{root relocation}\\
 \ \longrightarrow\ \text{strong Hahn transport}
 \ \longrightarrow\ \text{initiality of the image}.
\end{gather*}
For the inverse classification, the additional chain is
\begin{gather*}
 \text{missing ancestors of }p
 \ \longrightarrow\ \text{spine exponents}\\
 \ \longrightarrow\ \text{forced dyadic coefficients}
 \ \longrightarrow\ K_\alpha\subseteq H.
\end{gather*}
The sign-sequence proof is independent of the arithmetic of exponents.
In particular, no unproved assertion that $\tht$ preserves addition of
exponents occurs in the chain.

\emph{Source~02.} The convex-factor theorem depends on normal forms and
the initial Hahn criterion, then only on the proved projection and sign
lemmas. The Presburger criterion adds the divisible initial-embedding
theorem (\cref{isg:cf:fact:ehrlich}). The explicit nonsplit examples add
no further structural input. The elementary-equivalence consequences use
the classical characterization of Presburger models; the ultrapower
statement additionally uses \L o\'s's theorem. The Peano statement uses
the usual provably total factorial and summation recursions. The
surcomplex module assertions are coordinatewise consequences of the real
factor theorem. The merge statements
\cref{isg:cf:prop:cyclic,isg:cf:cor:limit,isg:cf:cor:periods} add
source~01's \cref{isg:sec:surgery,isg:sec:quotient} and, for the last,
statements of the trigonometry report cited by label.

\subsection{Finite regression checks supplied with the article}
The accompanying Python program \texttt{verify.py} (shipped in this
report as \fn{code/verify.py}) uses only the standard
library, exact rational arithmetic, and a fixed random seed. Its recorded
run performs \textbf{450,862 assertions}, all passing (shipped as
\fn{data/verification_report.json}). The tests cover
finite sign words through length seven for all relevant finite cone
minima; every prefix-closed subset of the binary tree through depth
three, including the empty tree; inverse-spine tests on a further family
of nonnegative initial trees; and 1,500 random pairs of finite rational
normal forms.

The count of prefix-closed trees is independently specified by
\[
 T_{-1}=1,\qquad T_d=1+T_{d-1}^2,
\]
so $T_0=2$, $T_1=5$, $T_2=26$, and $T_3=677$.
The recurrence counts the empty tree, or a root with independently
chosen initial left and right subtrees. This makes the finite-tree
coverage exhaustive at the stated depth, rather than a random sample.

Tests check inverse identities, order preservation, the exact predecessor
identity, the right-ancestor
identity away from the minimum, loss of the minimum's right ancestors,
initiality of transported trees, the exact inverse exponent criterion,
and the source dyadic obligations recovered by adjoining the spine.
The series tests check inverse transport, additivity, order, truncations,
and integrality of the bottom coefficient. Separate negative tests reject
the false global right-ancestor identity and unrestricted inversion.

The program does not represent arbitrary ordinals, arbitrary real
coefficients, or infinite Hahn supports. It does not check the imported
characterization theorem or establish absence of a transfinite proof gap.
Its role is regression testing of the finite combinatorics and coefficient
formulas. The written proofs, not these experiments, carry the general
claims.

When this report was placed, the program was rerun on a copy of the
directory with an explicit output file. It passed with the same 450,862
assertions, and its JSON output was identical in content to the shipped
record. The program's default output file is
\fn{verification_report.json} in the current directory, and the
delivered build script calls it with that name, so a rerun should always
be made on a copy (the collection README gives the commands).

\emph{Source~02.} Its program, shipped as
\fn{code/02-convex-factors-verify.py} (delivered as
\fn{code/verify.py}), uses only the Python standard library and exact
integers and rational fractions. Its recorded run, shipped as
\fn{data/02-convex-factors-verification.json} (delivered as
\fn{data/verification.json}), returns \texttt{PASS}. The main coverage
is:
\begin{center}\small
\begin{tabular}{@{}p{0.72\textwidth}r@{}}
\toprule
Check family & Count\\
\midrule
All nonempty convex intervals in the sign tree of lengths $\le5$ & 2,016\\
All nonempty meet-closed subsets of the sign tree of lengths $\le2$ & 61\\
Pairs checked for order, prefix reflection, and meet preservation & 1,398,609\\
Nested convex interval pairs in the sign tree of lengths $\le4$ & 46,376\\
Points checked for two-stage compression coherence & 324,632\\
Factorial transition identities & 100\\
Compatibility checks between factorial residues & 5,050\\
Integer-parameter splitting identities & 1,769\\
Finite projection composition identities & 24,300\\
\bottomrule
\end{tabular}
\end{center}
No row certifies a transfinite theorem. The program additionally checks
initiality of every finite compressed image (145,740 prefix checks), the
non-meet-closed collision example, the solvability of all finite
retraction systems through sixty stages, the factorial bound used in
\cref{isg:cf:lem:noninteger}, and additivity of interval projections
(2,970 checks); the full counts are in the JSON record. The successful
finite retractions are a deliberate regression case: the code must not
mistake finite solvability for a global splitting. The tests do not
verify the Ehrlich--Kaplan theorem, arbitrary ordinals, class
comprehension, completeness of Presburger arithmetic, recursive
saturation, or the Peano-model argument. Those depend on the written
proofs and the explicitly imported results, not on the test counts.

The program prints its JSON report and writes a file only when
\texttt{--output} is given; the delivered README's rerun command names the
delivered record as that output and would overwrite it, so it too should
be rerun on a copy. At this write it was rerun on a copy with a scratch
output file: it passed in about ten seconds, and the JSON content was
identical to the shipped record.

\subsection{Relationship to the Surreal repository}
The repository review examined the public README, the introductory
coverage and dependency ledger, and targeted searches for overlap
\cite{Repo}. Those materials describe extensive sign-sequence, canonical
cut, normal-form, Hahn-summation, and algebraic infrastructure. They also
explicitly distinguish source manuscripts from checked Lean declarations.
The present argument is designed to fit that distinction.

The review took place while the repository was changing. The initially
inspected README and ledger had blob hashes
\begin{center}\small\ttfamily
391afd7c64cebd48b47fbe43dae26c27e8ee8047\\
6a89497cb1de3e5c7a346201f7fc006d85082426.
\end{center}
A later metadata read reported commit
\begin{center}\small\ttfamily
173eb522fdf5d4bf201c097f76b38273eb4f95e1.
\end{center}
The initially returned recursive-tree hash is not used as a commit
identifier. This manuscript does not claim to have reviewed every
repository article or every companion archive, and no repository result
is silently promoted from ``pending'' to ``proved.''
The manuscript has no pinned commit. The commit \texttt{173eb52} above
is an ancestor of the commit \texttt{c6359e4} that placed the manuscript
in the collection. The manuscript makes no claim about the content of
any report, so no statement of it had to be corrected against the
collection as it now stands.

A sound implementation would first formalize the following mathematical
interfaces; the names are proposed work units, not declarations claimed
to exist:
\begin{enumerate}
\item \emph{Cone root relocation.} Define the three-way sign split,
      the inverse, the order equivalence, and
      \eqref{isg:eq:predidentity} for ordinal-length sequences.
\item \emph{Hahn reindexing with one coefficient multiplier.} Prove
      support transport, real-linearity, order preservation, truncation
      compatibility, and strong-sum interchange.
\item \emph{Concrete initial-subgroup criterion.} Supply a checked
      version of \cref{isg:fact:EK}, including both its coefficient and
      right-ancestor conditions, before claiming the final theorem.
\item \emph{Normalization and exact inverse.} Instantiate the criterion
      for the transformed groups; then derive the spine classification,
      quotient realization, and complex-coordinate extension.
\end{enumerate}

The largest imported proof obligation is the third item. A checked
conditional theorem using that criterion would be valuable, but would
not by itself constitute a checked theorem about every actual initial
subgroup of $\No$. This package supplies no Lean file with that
obligation hidden as an axiom or an unfinished proof.

\emph{Source~02.} Its repository review was made at commit
\begin{center}\small\ttfamily
9693b28c24e6fcb317ce47969a40185c5dfef402,
\end{center}
and included the main README, the formalization ledger, the article
catalogue exposed there, and the guide of the omnific groups and lattices
report; these distinguish mathematical drafts from Lean results. It did
not constitute a claim-by-claim audit of the repository or a local Lean
build. Source~02 read source~01 from the project library, including its
quotient, suspension and relative splitting arguments, and recorded that
its location in the inspected Git tree was not established, citing it as
a preceding project draft without assigning it a repository path. That
was correct at its pin: \texttt{9693b28} is an ancestor of the commit
\texttt{c6359e4} that placed source~01, and the preceding draft is now
\cref{isg:sec:intro,isg:sec:background,isg:sec:bottom,isg:sec:surgery,isg:sec:transport,isg:sec:normalization,isg:sec:classification,isg:sec:quotient,isg:sec:complex,isg:sec:examples}
of this report (its ``Section~8'' is \cref{isg:sec:quotient}). No theorem
of source~02 is called Lean-verified. A sensible implementation order
would formalize the finite and transfinite sign-tree compression first,
followed by the support projection lemmas and the transfer of the initial
Hahn criterion. The residue and factorial-presentation arguments can be
formalized in ordinary abelian-group and ordered-group libraries
independently of a full implementation of omnific integers. This is a
dependency plan, not delivered formal code.

\subsection{Novelty and literature status}
The new contributions claimed provisionally here are the explicit
root-relocation map together with its exact predecessor identity, its
use in omnific normalization, and the sharp dyadic-spine image
classification. The extremal, quotient, and surcomplex statements are
derived refinements of that construction. The known two-level example,
the general initial-Hahn-subgroup characterization, and the classification
of initial real coefficient groups are not claimed as new.

The literature check recovered the stated question in the original
preprint, thesis, and workshop account \cite{EK,Kaplan,Workshop}, and
examined adjacent sign-concatenation work \cite{BagHoeven}. Targeted
searches did not locate a prior resolution or the specific root-relocation
argument. That negative search result is limited evidence: it does not
exclude an unindexed paper, unpublished proof, or a reformulation already
known to a specialist. Neither mathematical correctness nor priority is
certified by the absence of a search result.

\emph{Source~02.} The convex-factor theorem and its surreal applications
are proposed contributions with complete proofs; neither priority nor
independent referee certification is claimed. The Ehrlich--Kaplan
criterion, Ehrlich's divisible embedding, Je\v r\'abek's residue theory
and residue-injectivity rigidity mechanism \cite{Jerabek}, and
$\Ext(\Q,\Z)\cong\Zhat/\Z$ \cite{Strickland} are classical and credited;
no new $\Ext$ computation is claimed, and the existence of rigid
Presburger models is not claimed as new. Bagayoko and van der Hoeven's
surreal substructures \cite{BagHoeven} are a related viewpoint, not a
dependency, and \cite{EnayatHamkinsWcislo} is background on internal
arithmetic only. A targeted source review did not find a prior
formulation of the main convex-factor theorem; the search was not
exhaustive.

\subsection{Provenance, placement review, and place in the collection}
\label{isg:sec:provenance}

\paragraph{Two sources.} This report is built from two manuscripts.

\emph{Source~01} is \emph{Discrete Initial Groups and Omnific
Normalization} (24 pages, 23~September 2026), manuscript~09 of the
collection's batch~28. It was delivered with a README, a source audit,
the finite checks \texttt{verify.py}, their recorded result and a build
script. It was placed in commit \texttt{c6359e4} as this directory and
written in \texttt{3d9dbe9}. Its LaTeX source is the basis of
\fn{article.tex}; its source audit is shipped unchanged as
\fn{source_audit.md}, the program and build script as \fn{code/verify.py}
and \fn{code/build.sh}, and the record as
\fn{data/verification_report.json}. Its delivered PDF and README are not
shipped. It has no pinned repository commit; its read identifiers are
those recorded above. It contributes
\cref{isg:sec:intro,isg:sec:background,isg:sec:bottom,isg:sec:surgery,isg:sec:transport,isg:sec:normalization,isg:sec:classification,isg:sec:quotient,isg:sec:complex,isg:sec:examples},
Appendix~\ref{isg:app:compact}, Questions
\ref{isg:q:convex}--\ref{isg:q:formal}, non-claims
\textup{(N1)}--\textup{(N15)} and most of the checklist of
Appendix~\ref{isg:app:checklist}.

\emph{Source~02} is \emph{Convex Factors and Profinite Obstructions for
Initial Surreal Groups} (23 pages, 23~September 2026), manuscript~09 of
the collection's batch~29. It was delivered with a README, a provenance
record, the program \texttt{code/verify.py} and its record
\texttt{data/verification.json}, and placed in commit \texttt{66d7e55}
under the prefix \fn{02-convex-factors-}: the provenance record is shipped
unchanged as \fn{02-convex-factors-PROVENANCE.md}, the program as
\fn{code/02-convex-factors-verify.py} and the record as
\fn{data/02-convex-factors-verification.json}. Its LaTeX source, PDF and
README are not shipped; the provenance record names the delivered file
names (\fn{code/verify.py}, \fn{data/verification.json}), which in this
directory carry the prefix. Its repository pin is \texttt{9693b28}. It
contributes \cref{isg:cf:sec:intro},
\cref{isg:cf:sec:cuts,isg:cf:sec:compression,isg:cf:sec:factors,isg:cf:sec:presburger,isg:cf:sec:factorial,isg:cf:sec:modeltheory,isg:cf:sec:complex,isg:cf:sec:boundaries}
except the statements marked \textup{[merge]}, Appendix~\ref{isg:cf:app:core},
Questions \ref{isg:cf:q:intrinsic}--\ref{isg:cf:q:multiplicative},
non-claims \textup{(N16)}--\textup{(N29)} and the corresponding rows of
the checklist.

\paragraph{How the merge was made (source~02).} Source~02 was placed as an
addition because it answers two questions this report named,
\cref{isg:q:convex,isg:q:limit}. Its material is printed as it stands,
with these choices. Results that both sources contain are printed once:
the Ehrlich--Kaplan criterion (\cref{isg:fact:EK}), the classification of
initial real groups (\cref{isg:fact:real}), the caution against Hahn
completion (\cref{isg:rem:nocompletion,isg:cf:rem:membership}), the
integer-bottom suspension (\cref{isg:thm:suspension}; source~02 re-proves
it from the criterion, by the same argument), the description of $\Oz$
(\eqref{isg:eq:Oz}) and the rectangular convention
(\cref{isg:def:rectangular}). Source~02's foundations section is replaced
by \cref{isg:sec:background} and \cref{isg:cf:sec:conventions}; its
symbols that collide with source~01's are renamed
(\cref{isg:cf:sec:conventions}); its Theorems A--C are Theorems D--F. Its
verification, repository, novelty, question and non-claim material is
merged into the corresponding parts of
\cref{isg:sec:audit,isg:sec:further}; its notation table is replaced by
the renaming table of \cref{isg:cf:sec:conventions} and the definitions
at their places, its dependency summary is in \cref{isg:sec:audit}, and
its independent-review checks are in
Appendix~\ref{isg:app:checklist}. Three statements are added in the merge, each
with a complete proof and marked \textup{[merge]}:
\cref{isg:cf:prop:cyclic} (source~02's construction for $\Z\varepsilon$
is source~01's $\rho\circ\tht$), \cref{isg:cf:cor:limit} (limit stages of
the iteration of \cref{isg:cor:iteration}) and \cref{isg:cf:cor:periods}
(period groups of the trigonometry report); \cref{isg:cf:rem:formulations}
records how the two versions of Ehrlich and Kaplan's optimality theorem
differ. The Ehrlich--Kaplan numbering table of \cref{isg:sec:intro} gained
the two rows for Section~8. No mathematical statement, proof or example
of source~01 was changed. Where source~01 said that convex quotients and
limit stages were not covered (after
\cref{isg:cor:iteration,isg:cor:relative}, in
\cref{isg:q:convex,isg:q:limit} and in \textup{(N10)}), the text now says
which source proves what; the title page, abstract, status box,
conventions, this section and the conclusion mention the second source.

\paragraph{What the write changed (source~01).} Every label now carries the prefix
\texttt{isg:}. The journal numbering of the Ehrlich--Kaplan results
(\cref{isg:sec:intro}) and the wording of the journal version of their
criterion (after \cref{isg:fact:EK}) were added, as were
\cref{isg:sec:conventions}, this subsection, the shipped file names, a
definition of $\ct$ (which the manuscript used without defining it),
labels on formerly unlabelled statements, and the collected list of
non-claims (\cref{isg:sec:nonclaims}). Cross-references now name each kind
of statement correctly; in the delivered source every reference to a
lemma, proposition, corollary, remark or imported fact printed as
``theorem''. No mathematical statement, proof or example was changed.

\paragraph{Placement review (source~01).} When source~01 was placed, every step
of the argument was checked by hand: the bottom-coefficient structure,
both cases of the predecessor identity \eqref{isg:eq:predidentity} and its
failure at the old minimum, the ambient transport, every dyadic obligation
in the proof of \cref{isg:thm:A}, both directions of \cref{isg:thm:B}
with its two counterexamples, the extremal groups, the size and depth
corollaries, the splitting and suspension theorems, and every example of
\cref{isg:sec:examples}. No gap was found. The Ehrlich--Kaplan criterion
was checked as a statement, against the arXiv text and the author-hosted
text of the journal version: \cref{isg:fact:EK} is a faithful reading of
\cite[Theorem~1 and its proof]{EK}, equivalently of Theorem~5.1 of
\cite{EKjournal}, and \cref{isg:fact:real} is \cite[Lemma~3]{EK}. The
criterion was not re-proved; the whole argument rests on it. This review
is not a refereeing, and the solution proposed here has not been refereed.
Its scope is source~01's material; it does not cover
\cref{isg:cf:sec:cuts,isg:cf:sec:compression,isg:cf:sec:factors,isg:cf:sec:presburger,isg:cf:sec:factorial,isg:cf:sec:modeltheory,isg:cf:sec:complex,isg:cf:sec:boundaries}.

\paragraph{Placement review (source~02).} When source~02 was placed, its
proofs were checked by hand and no error was found: the exponent-cut
lemmas and the canonical splitting, the compression theorem with its
uniqueness and nested coherence, the transfer of every dyadic obligation
in \cref{isg:cf:thm:transport} (coefficient groups are unchanged and
every target right ancestor comes from a source one by
\eqref{isg:cf:eq:rightreflect}), the residue lemmas and
\cref{isg:cf:prop:residuesplit}, the factorial family (its table, the
transitions, $n!x_n=x_1+S_ne$, $S_n\le2(n-1)!$ and the finite
retractions), the extension classes and rigidity, the dense example
(initiality of $W$: its exponents $0,1,2$ have no right ancestors in the
class, and $\Q$ and $\Z$ are initial in $\R$), the omitted type, the
ultrapower and Peano arguments, the omnific residues and the rectangular
transport. The compression was re-implemented independently and tested
on all 496 nonempty convex intervals of the sign tree of lengths at most
four (order, injectivity, initial image, prefix and right-ancestor
reflection) and for nested coherence on the tree of lengths at most
three. The elementary equivalence in \cref{isg:cf:thm:dense} rests on a
standard two-sorted interpretation, checked as an argument. The
Ehrlich--Kaplan criterion and Ehrlich's embedding theorem were not
re-proved. For the merge statements, \eqref{isg:cf:eq:cyclic} was
confirmed by a finite computation for $\alpha\le3$ and words of length at
most six, and the iteration in \cref{isg:cf:cor:limit} for three removal
steps on those trees; that computation is not shipped. This review is not
a refereeing, and it does not extend to source~01's material beyond the
consistency statements of \cref{isg:cf:sec:cyclic,isg:cf:sec:limit}.

\paragraph{Literature.} The placement review of source~01 ran targeted web
searches (``discrete initial subgroup'', ``initial subgroup'' together
with ``omnific integers'', and the title of \cite{EK}). They found only
the Ehrlich--Kaplan paper, Kaplan's thesis and unrelated pages, and no
resolution of the question. The search was targeted, not exhaustive. The
citations of Kaplan's thesis (its Sections~6.2 and~7.1) and of the
Oberwolfach report (printed pages 3355--3356) are the manuscript's; they
were not re-checked at placement. For source~02 no separate literature
search is recorded at placement; its own review was targeted and found no
matching formulation of the convex-factor theorem. Its citations of
Je\v r\'abek, Strickland, Enayat--Hamkins--Wcis\l o and Bagayoko--van der
Hoeven are the manuscript's and were not re-checked. Its account of
Ehrlich and Kaplan's Section~8 was compared with the arXiv text and the
authors' text of the journal version at this write
(\cref{isg:cf:rem:formulations}).

\paragraph{Place in the collection.} At the placement of source~01, no
other report named Question~2 (Question~9.1) of Ehrlich and Kaplan, and
none treated initial subgroups: a search of the reports and of the Lean sources for
``initial subgroup'', ``discrete initial'' and ``discretely ordered
initial'' found nothing outside this directory. The paper \cite{EK} is
cited as background in two reports: the foundations report cites it among
the initial-embedding results next to its sign-tree categoricity
(\replabel{found:sub:signtree}), and the definable-surreals report for
the simplicity theory of simplest separators
(\replabel{dsn:sec:foundations}). The latter's maximal initial elementary
exponential core (\replabel{dsn:thm:core}) concerns initial
\emph{subfields} of definable surreals; it is related in spirit and
neither overlaps nor conflicts with the present groups. The omnific
integers used here are those of the Diophantine and set-sized-quotient
reports (\replabel{odg:def:rings}, \replabel{osq:eq:Ozdef}); their
discreteness is \replabel{odg:prop:units}, and the constant coefficient
is the additive map of \replabel{osq:lem:ct}. The Diophantine report and
the foundations report also cite a different Ehrlich--Kaplan paper,
\emph{Surreal ordered exponential fields}, for the initial integer part
$\Oz\cap K$ of an initial subfield $K$ (the status note after
\replabel{odg:thm:floor}, and \replabel{found:sub:canonicalexp}). That
statement concerns subrings and integer parts. The groups $H\subseteq\Oz$
here are additive subgroups and $\Rect$ is not multiplicative, so there is
no conflict. No \texttt{isg:} statement has a Lean implementation.

Source~02 adds one neighbour. The trigonometry report studies the
normalized period groups $\mathcal I(\Psi)$ of surcomplex phases, which
are $\Z$-groups for set-sized fields, with the same Je\v r\'abek residue
theory: its residue $\res_{\mathcal I}$ and maximal divisible subgroup
$\mathcal I^{\mathrm{div}}$ (\replabel{trigonometry:per:sec:defect}) are
$\res_G$ and $G^{\dv}$ here, and its profinite defect is that residue
modulo $\Z$. By \cref{isg:cf:cor:periods} a period group is initially
realizable exactly when its profinite defect vanishes. There is no
conflict: that report's initiality statements concern the kernel-multiplier
ring (\replabel{trigonometry:per:thm:robustness}), a subring, and the
initial integer part $\Oz\cap F$, not the realizability of period groups.
The divisible-kernel lemma is classical in both places. The set-sized
quotient report (\replabel{osq:eq:Ozdef}) and the Diophantine report use
the same $\Pi$ and $\ct$ as \cref{isg:cf:prop:omnificresidue}.

Related (batch 33). The exponential-relations report
(\nolinkurl{docs/foundations-and-computation/exponential-relations-over-omnific-integers})
reads the class-sized
$\Oz$ formula by formula as a $\Z$-group, and so as a model of Presburger
arithmetic in the sense of \cref{isg:cf:sec:intro}. Its division with
remainder (\replabel{exr:prop:division}) is \cref{isg:cf:lem:residues} for
$G=\Oz$, $e=1$, re-derived independently, and its lexicographic splitting
$\Oz=\Pi\oplus\Z$ is the split form of \cref{isg:cf:main:presburger} with
$G^{\dv}=\Pi$. It expands this group by the zero predicates of finite
exponential sums with rational slopes and proves the result a definitional
expansion of Presburger arithmetic, with $\Z$ an elementary substructure of
$\Oz$ (\replabel{exr:thm:rational}). Nothing there concerns initiality, and
no statement here concerns exponential predicates.

\section{Further questions and exact boundaries}\label{isg:sec:further}

The construction establishes an additive normal form while deliberately
relinquishing several stronger kinds of compatibility. The following
questions isolate possible further work without asserting that each is
a previously published open problem.

\begin{question}[Other convex quotients; answered by source~02]\label{isg:q:convex}
Which convex subgroups $C$ of an initial group $A\subseteq\No$ admit
an explicit initial realization of $A/C$ by a normal-form transport?
\end{question}
Source~01 handles $C=\Z\varepsilon$. For a larger convex subgroup,
one would need to identify the removed exponent region and determine
whether its boundary can be relocated without creating unmet dyadic
obligations. A numerical quotient of exponent chains alone would not
settle the simplicity condition.

\emph{Status: answered, by source~02.} All of them.
\Cref{isg:cf:main:factor,isg:cf:thm:transport} realize $A/C$ explicitly:
the removed exponent region is the lower segment $\Gamma_C$
(\cref{isg:cf:lem:arch}), the quotient is represented by the upper
projection $P_{\Gamma^C}$, a difference of leading truncations
(\cref{isg:cf:lem:restriction,isg:cf:thm:split}), and compression of
$\Gamma^C$ creates no unmet dyadic obligation because it reflects right
ancestors \eqref{isg:cf:eq:rightreflect}; it is not a numerical quotient
of exponent chains (\cref{isg:cf:ex:compression}). For
$C=\Z\varepsilon$ the transport is source~01's
(\cref{isg:cf:prop:cyclic}).

\begin{question}[Limit iteration; answered by source~02]\label{isg:q:limit}
When an iterated removal of discrete bottom layers continues through
infinitely many stages, which hypotheses give a coherent initial
realization at a limit stage?
\end{question}
In source~01, finite iteration follows from \cref{isg:cor:iteration};
limit iteration requires a separate control of supports and of the
changing sign-tree presentations. No commutation with such an unproved
transfinite limit is used in source~01.

\emph{Status: answered, by source~02 with a corollary of the merge.} No
hypothesis is needed. The union of the removed layers is a convex
subgroup, whose quotient is realized directly from its exponent cut
(\cref{isg:cf:cor:closure}), coherently with every earlier stage
(\cref{isg:cf:cor:coherentfactors}); \cref{isg:cf:cor:limit} states this
for the iteration of \cref{isg:cor:iteration}. The required control is
supplied by compression, which is defined position by position and never
as a limit of successive deletions (\cref{isg:cf:rem:limit}).

\begin{question}[Multiplicative compatibility on restricted classes]\label{isg:q:multiplicative}
For which special initial groups does some normalization to \Oz also
preserve products wherever those products are part of a specified
algebraic structure?
\end{question}
The counterexample with $1\cdot1$ rules out a general claim for our
normalization. Any useful positive result must impose additional
conditions on exponent addition and coefficient rescaling, not merely
on the order of supports.

\begin{question}[Formal certification and uniqueness principles]\label{isg:q:formal}
Can the root-relocation proof and the concrete initiality criterion be
fully verified in the repository's foundational model, and what natural
extra conditions single out this normalization among the possible
ordered-group isomorphisms to initial omnific groups?
\end{question}
Our formulas are canonical relative to the given sign model, normal
forms, and least positive scale. They are not asserted to be invariant
under arbitrary abstract group isomorphisms. The uniqueness in
\cref{isg:thm:B} concerns the inverse under the fixed map, not uniqueness of
all possible initial embeddings.

The next four questions are source~02's. They mark extensions of the
proved statements rather than prerequisites left unproved in their
proofs, and they are not asserted to be previously published open
problems (the second is related to the published Question~1 / 8.1 of
Ehrlich and Kaplan, \cref{isg:cf:sec:optimality}).

\begin{question}[Intrinsic characterization]\label{isg:cf:q:intrinsic}
Which initially realizable ordered groups are characterized by intrinsic
conditions on all convex factors and the splitting of convex extensions,
without first choosing a Hahn presentation?
\end{question}

\begin{question}[Intermediate theories]\label{isg:cf:q:intermediate}
Can the nonsplit convex-extension obstruction be arranged inside a
set-sized model of every proper intermediate theory appearing in the
Ehrlich--Kaplan optimality question?
\end{question}

\begin{question}[Arithmetic fragments]\label{isg:cf:q:fragment}
What is the weakest arithmetic fragment whose nonstandard models have
ordered additive group completions excluded by the factorial residue
type? Which analogous computable residue types yield sharper bounds?
\end{question}

\begin{question}[Compression and multiplication]\label{isg:cf:q:multiplicative}
Under what additional compatibility conditions does exponent-tree
compression preserve a specified multiplication or a specified action
of a coefficient domain? The unrestricted additive result provides no
such compatibility automatically.
\end{question}
This is the counterpart for compression of \cref{isg:q:multiplicative},
which asks the same for the normalization $\Rect$.

\subsection{Non-claims, collected}\label{isg:sec:nonclaims}
The limitations stated at their places in the text, in the shipped
source audit and in the shipped provenance record of source~02 are
collected here, so that none is lost. Each stays in force where it was
stated. Items \textup{(N1)}--\textup{(N15)} are source~01's; in them
``the manuscript'' means source~01. Items \textup{(N16)}--\textup{(N29)}
are source~02's.
\begin{enumerate}[label=\textup{(N\arabic*)},leftmargin=*]
\item \emph{Status.} The manuscript is unrefereed. It is a proposed
      solution with full proposed proofs, not an independently certified
      resolution; neither correctness nor novelty is certified, and no
      priority is claimed. The placement review
      (\cref{isg:sec:provenance}) is not a refereeing.
\item \emph{Literature.} The literature and repository searches were
      targeted, not exhaustive. A negative search result is limited
      evidence and does not exclude an unindexed paper, an unpublished
      proof or a reformulation known to specialists. That a question was
      posed in 2015 and 2018 does not establish that it was still
      unresolved on the review date.
\item \emph{Imported criterion.} \Cref{isg:fact:EK} is imported, not
      re-proved, by the manuscript or at placement. No proof here
      replaces it by the weaker assertion that initiality of the exponent
      class alone suffices.
\item \emph{Not claimed as new.} The two-level example and its map,
      the initial-Hahn-subgroup characterization, the classification of
      initial real groups (\cref{isg:fact:real}) and Conway normal forms
      are imported and cited. The extremal, quotient and surcomplex
      statements are derived refinements of the construction.
\item \emph{Not a simplicity isomorphism.} $\Rect$ is an ordered-group
      isomorphism onto an initial image. It need not preserve the
      simplicity relation (\cref{isg:rem:notsimplicity}); it is not
      multiplicative, does not preserve norms, and does not fix ordinary
      constants: $\Rect(1)\ne1$ in general (\cref{isg:sec:examples}).
\item \emph{No classification.} No classification of all ordered abelian
      groups admitting initial embeddings is claimed, and the
      initial-monoid problem (Question~3 of \cite{EK}, Question~9.2 of
      the journal version) is not addressed.
\item \emph{Exponent maps.} $\Theta_\alpha$ is not a homomorphism of
      exponent groups, and \eqref{isg:eq:leadingtransport} is not
      preservation of a multiplicative valuation
      (\cref{isg:rem:preserved}). \Cref{isg:prop:length} bounds exponent
      lengths, not the birthday of a transformed normal form.
\item \emph{No internal summation.} The strong-sum statements concern the
      ambient Hahn spaces; they give no internal summation on a subgroup
      $A$ (\cref{isg:rem:nocompletion}), and \cref{isg:lem:isolate} does
      not license arbitrary infinite subseries.
\item \emph{Invariance and uniqueness.} The ordinal $\alpha$ is not an
      invariant of the abstract ordered group. The uniqueness in
      \cref{isg:thm:B} is relative to the fixed map, not uniqueness of all
      initial embeddings, and the formulas are not claimed invariant under
      abstract isomorphisms.
\item \emph{Quotients.} It is not claimed that every convex quotient of
      an initial group has an initial realization. Iteration is covered
      only at finite stages (\cref{isg:cor:iteration}), with no claim at a
      limit stage, and \cref{isg:cor:relative} is a relative criterion
      that does not decide realizability of every quotient. (This is a
      statement about source~01. Source~02 proves both missing
      statements: \cref{isg:cf:main:factor} and
      \cref{isg:cf:cor:closure,isg:cf:cor:limit}. The relative criterion
      remains relative; \cref{isg:cf:main:presburger} decides the case of
      $\Z$-groups.)
\item \emph{Surcomplex scope.} \Cref{isg:thm:complex} is Gaussian-linear
      and coordinatewise initial. It gives no field isomorphism, no
      canonical sign-tree or simplicity relation on $\No[i]$, no
      preservation of the modulus, and no compatibility with
      multiplication by general surcomplex scalars. \Cref{isg:cor:linear}
      transports the right-hand side as well and does not apply to
      equations with ordinary constants held fixed, to nonlinear
      polynomials, to arbitrary surreal scalar matrices, or to norms and
      inner products.
\item \emph{Finite checks.} The finite checks do not verify transfinite
      concatenation or limit cases, arbitrary set-supported Hahn sums, the
      imported criterion, or the absence of a transfinite proof gap; the
      written proofs carry the general claims. Compiling and rendering the
      document establish its integrity, not its correctness.
\item \emph{No Lean.} No new Lean formalization and no Lean build are
      claimed. The proposed formalization units of \cref{isg:sec:audit} are
      not existing declarations, and a checked conditional theorem using
      the imported criterion would not by itself be a checked theorem
      about every initial subgroup of $\No$.
\item \emph{Repository review.} The manuscript's repository review was
      targeted: it did not review every report or archive, it treated
      truncated metadata as neither a complete inventory nor a proof of
      absence, it took no repository description as a certificate, and it
      promoted no repository result from pending to proved. The argument
      relies on no repository manuscript, and no repository Lean build was
      run for it.
\item \emph{Sources used only for context.} Bagayoko and van der Hoeven's
      sign-sequence conventions are context, not a source of the
      normalization or of its range theorem. The questions of this section
      are not asserted to be previously published open problems, and
      $d(H)=\On$ is a formal marker, not a surreal number or a set
      ordinal.
\item \emph{Status of source~02.} The convex-factor theorem and its
      applications are proposed contributions with complete proofs.
      Neither priority nor independent referee certification is claimed,
      no Lean formalization or Lean build is claimed, and the placement
      review (\cref{isg:sec:provenance}) is not a refereeing.
\item \emph{Credited, not new.} The Ehrlich--Kaplan criterion,
      Ehrlich's divisible embedding (\cref{isg:cf:fact:ehrlich}),
      Je\v r\'abek's residue theory and rigidity mechanism, and
      $\Ext(\Q,\Z)\cong\Zhat/\Z$ are classical and credited; no new
      computation of $\Ext$ is claimed, and the existence of rigid
      Presburger models is not claimed as a new phenomenon. The results
      of source~01 that source~02 uses (the suspension, the relative
      splitting criterion) are not claimed again as new.
\item \emph{No classification.} There is no classification of initially
      realizable groups. Convex splitting is necessary, and its
      sufficiency is not proved (\cref{isg:cf:sec:boundaries}).
\item \emph{Optimality question.} Source~02 does not resolve the
      Ehrlich--Kaplan set-model optimality question (Question~1 of
      \cite{EK}, Question~8.1 of the journal version) and gives no
      construction uniform over the intermediate theories
      (\cref{isg:cf:sec:optimality,isg:cf:rem:formulations}). It answers
      none of Ehrlich and Kaplan's questions.
\item \emph{Realization-dependent splitting.} The splitting of
      \cref{isg:cf:thm:split} depends on the chosen realization and is not
      invariant under abstract automorphisms; only the Presburger split
      case has the uniqueness of \cref{isg:cf:prop:residuesplit}.
\item \emph{Realizability, not literal initiality.} The results concern
      initial \emph{realizability}, not literal initiality of $C$ or of a
      complement in its original position.
\item \emph{No Hahn product, no completion.} An initial group is not
      identified with the full Hahn product of its factors, and no closure
      under ambient strong sums is asserted
      (\cref{isg:cf:rem:membership,isg:cf:rem:strongsums}).
\item \emph{Additive only.} The results are additive and
      module-theoretic. There is no multiplicative, field, ring,
      differential or exponential equivalence; the compression need not
      respect addition of exponents.
\item \emph{Surcomplex scope of source~02.} The rectangular notion is
      coordinate-dependent; no canonical sign tree on $\No[i]$ is
      postulated.
\item \emph{Finite checks of source~02.} The finite regression does not
      verify the Ehrlich--Kaplan theorem, arbitrary ordinals, class
      comprehension, completeness of Presburger arithmetic, recursive
      saturation or the Peano-model argument; the successful finite
      retractions are not a global splitting (\cref{isg:cf:rem:finite}).
      The delivered build and rendering checks concern presentation, not
      correctness.
\item \emph{Arithmetic fragment.} The fragment of arithmetic used for
      \cref{isg:cf:thm:PA} is not optimized.
\item \emph{Repository review of source~02.} Its review, at pin
      \texttt{9693b28}, was targeted: it was not a claim-by-claim audit
      of the repository, and no local Lean build was run.
\item \emph{Predecessor citation.} Source~02 cited source~01 without
      inventing a repository path for it; the current location
      (\cref{isg:sec:audit}) was supplied at this write.
\item \emph{Finite solvability.} Solvability of every finite retraction
      system is not a global splitting (\cref{isg:cf:rem:finite}).
\end{enumerate}
The statements marked \textup{[merge]} add no claim of priority; they are
proved from the two sources and, for \cref{isg:cf:cor:periods}, from
statements of the trigonometry report cited by label, whose own
non-claims remain in force there.

\section{Conclusion}\label{isg:sec:conclusion}
The obstruction in the Ehrlich--Kaplan question is not the numerical
size of the least positive element. Scalar multiplication removes that
size immediately. The obstruction is the interaction of the exponent
tree with coefficient groups forced by right ancestors.

The explicit map $\Theta_\alpha$ moves the minimum exponent to zero
and preserves all right-ancestor relations away from it. The coefficient
map then normalizes the discrete bottom group without disturbing the
higher groups. This gives the proposed affirmative proof. The relations
lost at the minimum are exactly those recorded by the dyadic spine,
which yields a necessary-and-sufficient image criterion rather than
only a sufficient construction.

The resulting theory is additive but strong: it transports arbitrary
set-supported normal forms without changing their support order types,
commutes with strong sums, identifies extremal initial groups at each
least positive scale, and produces initial models of the bottom cyclic
quotients and rectangular surcomplex modules. The remaining obligations
are independent mathematical review, a more exhaustive priority check,
and a complete formal verification of the identified dependency chain.

Source~02's compression construction makes convex-factor closure a
concrete normal-form theorem: retain an exponent interval, delete
precisely the ancestor positions outside it, and invoke the initial Hahn
criterion. It is coherent for nested intervals and works at arbitrary
transfinite cuts; for the bottom cyclic layer it is source~01's
$\rho\circ\tht$. At the same time, normal-form projection forces all
convex extensions inside an initial realization to split. In additive
arithmetic this yields a sharp boundary: a Presburger group has an
initial surreal realization exactly when each element's coherent
standard residues are ordinary integers. The explicit factorial family
places continuum many countable rigid models on the other side of that
boundary. Dense examples and omitted-type arguments show that the
distinction survives beyond discreteness and cannot be captured by
first-order theory alone. None of this is a certified resolution of the
full classification problem or of Ehrlich and Kaplan's optimality
question.

\appendix
\section{A compact proof for independent checking}\label{isg:app:compact}

This appendix compresses the main existence argument without changing
its hypotheses. The full case analysis and support arguments are in the
body of the paper.

Let $A$ be a discrete initial subgroup. Its least positive member has a
single normal-form term: otherwise a nonzero tail in $A$ would be a
smaller positive element after taking absolute value. Write that term
as $c\omega^p$. The exponent class has minimum $p$, which, being the
minimum of an initial tree, is all-minus, say $p=-\alpha$.
The bottom real coefficient group is discrete and initial, hence
$c\Z$ with $c=2^{-n}$.

On $[p,\infty)$ define
\[
 p\mapsto0,\qquad
 p\cat\plus\cat u\mapsto\plus\cat p\cat u,\qquad
 x\mapsto\plus\cat x\quad(p\not\pre x).
\]
The inverse removes the initial plus and reinserts a plus immediately
after the block $p$ whenever the remaining word starts with $p$.
The two maps are increasing inverse bijections. For $x>p$,
\[
 \Pred(\Theta_\alpha x)
   =\Theta_\alpha[\Pred(x)\cup\{p\}],
 \qquad
 \Rs(\Theta_\alpha x)=\Theta_\alpha[\Rs(x)].
\]
Therefore an initial exponent class with minimum $p$ maps to an initial
one with minimum zero, and every new right-ancestor obligation comes
from an old one above $p$.

Transport each nonbottom term $r\omega^\gamma$ to
$r\omega^{\Theta_\alpha\gamma}$, and transport the bottom term
$r\omega^p$ to $r/c$. This is a real-linear order-isomorphism on the
ambient Hahn spaces. Its image of $A$ is truncation closed and
cross-sectional, has coefficient group $\Z$ at zero, and has unchanged
initial real coefficient groups at positive exponents. The right-ancestor
identity preserves all required dyadic containments. By the concrete
Ehrlich--Kaplan criterion, the image is initial. It lies in \Oz because
all exponents are nonnegative and the constant coefficient is integral.

For the inverse classification, only the ancestors of the old minimum
were discarded. They are $-\beta$ for $\beta<\alpha$, and their images
are $q_\beta=\plus\cat\minus^\beta$. Initial inversion requires
those exponents, while the right-ancestor condition at $p$ requires all
their dyadic coefficients. Together these requirements are exactly
$K_\alpha\subseteq H$.

\clearpage
\section{Proof-review checklist}\label{isg:app:checklist}
An independent review can focus on the following specific potential
failure points rather than checking experimental output as though it
were a proof.

\begin{longtable}{@{}p{0.30\textwidth}p{0.64\textwidth}@{}}
\toprule
Issue & Location and required check\\
\midrule
\endhead
Concrete versus existential initiality &
\Cref{isg:fact:EK}: the imported sufficiency statement must apply to the
specified canonical image, not only to an unspecified isomorphic copy.\\[4pt]
Minimum exponent &
\Cref{isg:prop:bottom}: truncation and subtraction isolate a nonzero tail;
its absolute value is smaller than the positive leading monomial.\\[4pt]
Limit sign lengths &
\Cref{isg:lem:inverse,isg:thm:tree}: concatenated tails use ordinal interval
order types. No subtraction of ordinal lengths or removal of a last
limit sign is permitted.\\[4pt]
Exceptional old minimum &
\Cref{isg:thm:tree,isg:rem:exception}: the right-ancestor equality excludes
$x=p$. Its failure is used, not ignored, in the inverse classification.\\[4pt]
Coefficient group scaling &
\Cref{isg:lem:coefficients}: only the bottom coefficient is rescaled;
all right-ancestor obligations in the target are at unchanged groups.\\[4pt]
Infinite subgroup membership &
\Cref{isg:sec:transport}: transport the actual subgroup. Never replace it
with all normal forms allowed by its coefficient groups.\\[4pt]
Sharp inverse criterion &
\Cref{isg:thm:B}: inverse exponent initiality is not enough; the missing
minimum's ancestors must carry all dyadic coefficients.\\[4pt]
Quotient order &
\Cref{isg:lem:split,isg:prop:quotient}: higher nonzero terms dominate the cyclic
bottom group, and positive-root deletion supplies an initial image.\\[4pt]
Class-sized scope &
All supports and summation indices are sets; correspondences of
class-sized groups are expressed with class parameters.\\[4pt]
Surcomplex scope &
\Cref{isg:thm:complex}: the conclusion is Gaussian-linear and coordinatewise
initial, not a field isomorphism or a new canonical complex simplicity
relation.\\[4pt]
\multicolumn{2}{@{}l}{\emph{Source~02 (convex factors).}}\\[2pt]
Actual normal-form image &
\Cref{isg:fact:EK}: its sufficiency part must apply to the actual
transported image $\mathcal F_E(G)$, not just to an abstract copy.\\[4pt]
Restriction stays in $G$ &
\Cref{isg:cf:lem:restriction}: $P_E$ is the difference of two leading
truncations, so no infinite subseries is used.\\[4pt]
Meet retention &
\Cref{isg:cf:thm:compression}: compression must retain the sign at the
meet of an incomparable pair; convexity puts the meet in $E$
(\cref{isg:cf:lem:meetclosed,isg:cf:ex:badmeet}).\\[4pt]
Every target right ancestor &
\Cref{isg:cf:thm:compression,isg:cf:thm:transport}: every target right
ancestor must be accounted for, not merely every source ancestor; this is
the equality \eqref{isg:cf:eq:rightreflect}.\\[4pt]
Residue sign &
\Cref{isg:cf:prop:family}: the residue of $x_1$ is $-\zeta$, not
$\zeta$; factorial moduli are cofinal for divisibility.\\[4pt]
Divisible complement &
\Cref{isg:cf:prop:residuesplit}: a complement to $\Z e$ is divisible
because it is isomorphic to the quotient, not because every torsion-free
subgroup is divisible.\\[4pt]
Noninteger parameter &
\Cref{isg:cf:lem:noninteger}: $\zeta_*$ is shown to be neither a
nonnegative nor a negative ordinary integer.\\[4pt]
Internal sums &
\Cref{isg:cf:thm:PA}: the sum in a nonstandard Peano model is an
internally coded finite sum, not an external infinite sum in $G$.\\[4pt]
Merge statements &
\Cref{isg:cf:prop:cyclic,isg:cf:cor:limit,isg:cf:cor:periods}: the
identification with $\rho\circ\tht$, the one-point layers of the
iteration, and the order of $\mathcal I/\Z\cong\Pi_F$.\\
\bottomrule
\end{longtable}

\section{The core convex-factor proof in one place}\label{isg:cf:app:core}
This appendix is source~02's. For independent checking, the logical chain
can be compressed without changing its hypotheses. Let $G\subseteq\No$ be
initial and let $C\subseteq C'$ be convex. The active exponents form an
initial class $\Gamma$. Convexity says that membership of a nonzero
element in $C$ or $C'$ is determined exactly by whether its leading
monomial belongs to that subgroup. Thus each has a lower exponent segment,
$\Gamma_C$ and $\Gamma_{C'}$.

The interval $E=\Gamma_{C'}\setminus\Gamma_C$ represents $C'/C$:
coefficient restriction to $E$ stays in $G$, because it is a difference
of two leading truncations. The kernel on $C'$ is $C$, and the
leading sign gives the quotient order.

Any two incomparable elements of $E$ have their common prefix between
them numerically. It is in $\Gamma$ by initiality and in $E$ by
convexity. Thus $E$ is meet closed. In each $\gamma\in E$, retain the
signs at precisely those positions $\beta$ for which
$\gamma\restriction\beta\in E$. Index them by their ordinal order
type. A source prefix in $E$ becomes a target prefix; incomparable
points keep opposite signs at their retained meet. Every target prefix
comes from one of the retained source prefixes. Hence the map is an
order-isomorphism onto an initial exponent class and preserves and
reflects right ancestry.

Reindex the normal forms supported in $E$ by this map. The coefficient
groups do not change, truncation closure and cross-sectionality persist,
and every target dyadic right-ancestor requirement is one of the
original requirements. The Ehrlich--Kaplan criterion makes this actual
image an initial subgroup of $\No$. This proves the convex-factor
assertion. Upper-exponent restriction simultaneously supplies the split
section for each convex subgroup sequence.

\clearpage
\phantomsection
\addcontentsline{toc}{section}{References}
\begin{thebibliography}{99}
\bibitem{EK}
P.~Ehrlich and E.~Kaplan,
\emph{Number systems with simplicity hierarchies: a generalization of
Conway's theory of surreal numbers II},
arXiv:1512.04001v1, 13 December 2015.
\url{https://arxiv.org/abs/1512.04001v1}.
The version explicitly used for theorem numbering: Lemma~3,
Theorem~1 and its concrete sufficiency proof, Proposition~9,
and Section~9, Question~2; for source~02 also Section~8, Theorem~6 and
Question~1.

\bibitem{EKjournal}
P.~Ehrlich and E.~Kaplan,
\emph{Number systems with simplicity hierarchies: a generalization of
Conway's theory of surreal numbers II},
The Journal of Symbolic Logic \textbf{83} (2018), no.~2, 617--633.
\url{https://doi.org/10.1017/jsl.2017.9}.
Publisher metadata checked; the arXiv version above supplies the cited
internal numbering. The journal numbering given in \cref{isg:sec:intro}
(Lemma~5.1, Theorem~5.1, Proposition~5.2, Theorem~8.1, Questions~8.1,
9.1 and~9.2) was
read, at placement, in the authors' text of this version hosted on the
second author's web page; the publisher's typeset pages were not
compared.

\bibitem{Kaplan}
E.~Kaplan,
\emph{Initial Embeddings in the Surreal Number Tree},
Honors thesis, Ohio University, 2015.
OhioLINK accession \texttt{ouhonors1429615758}.
See Sections~6.2 and~7.1 for omnific integers and the group question.
\url{https://etd.ohiolink.edu/acprod/odb_etd/ws/send_file/send?accession=ouhonors1429615758&disposition=inline}.

\bibitem{Workshop}
\emph{Mini-Workshop: Surreal Numbers, Surreal Analysis, Hahn Fields
and Derivations}, Oberwolfach Report 60/2016,
workshop account containing the initial-subgroup question and the
concrete Hahn criterion, especially printed pages 3355--3356.
\url{https://ems.press/content/serial-article-files/46663}.

\bibitem{BagHoeven}
V.~Bagayoko and J.~van der Hoeven,
\emph{Surreal substructures}, author-hosted manuscript,
first version 8 June 2019, corrected version 16 April 2022.
\url{https://www.texmacs.org/joris/sss/sss.html}.
See the sign-sequence conventions and concatenation operations.
Source~02 cites the arXiv version, arXiv:2305.02001 (2023),
\url{https://arxiv.org/abs/2305.02001}, for surreal substructures and
convex partitions, as a related viewpoint only.

\bibitem{Repo}
V.~Reshetnikov,
\emph{Surreal}, public research and Lean-formalization repository,
reviewed 23 September 2026.
\url{https://github.com/VladimirReshetnikov/Surreal}.
Reviewed entry points: \texttt{README.md} and
\texttt{docs/FORMALIZATION.md}; precise read identifiers and the scope
of the review are recorded in \cref{isg:sec:audit} and the accompanying
\texttt{source\_audit.md}. Source~02 inspected the repository at commit
\texttt{9693b28c24e6fcb317ce47969a40185c5dfef402}, 23 September 2026
(\cref{isg:sec:audit}); its review does not certify that every file was
read.

\bibitem{Ehrlich}
P.~Ehrlich, \emph{Number systems with simplicity hierarchies:
a generalization of Conway's theory of surreal numbers},
The Journal of Symbolic Logic \textbf{66} (2001), 1231--1258;
corrigendum, \textbf{70} (2005), 1022.
The divisible initial-embedding theorem (\cref{isg:cf:fact:ehrlich}) is
also recalled in \cite{EK}. Cited by source~02.

\bibitem{Conway}
J.~H.~Conway, \emph{On Numbers and Games}, 2nd ed., A K Peters, 2001.
Cited by source~02 for normal forms.

\bibitem{Gonshor}
H.~Gonshor, \emph{An Introduction to the Theory of Surreal Numbers},
London Mathematical Society Lecture Note Series 110,
Cambridge University Press, 1986. Cited by source~02 for normal forms.

\bibitem{Jerabek}
E.~Je\v r\'abek, \emph{Rigid models of Presburger arithmetic},
Mathematical Logic Quarterly \textbf{65} (2019), no.~1, 108--115.
\url{https://doi.org/10.1002/malq.201800019};
\url{https://arxiv.org/abs/1803.05797v2} (2018).
Cited by source~02 for $\Z$-groups, their residue map, divisible kernel
and rigidity mechanism.

\bibitem{Strickland}
N.~Strickland, \emph{Algebraic theory of abelian groups},
\url{https://arxiv.org/abs/2001.10469} (2020),
especially Section~8 and Example~8.27. Cited by source~02 for
$\Ext(\Q,\Z)\cong\Zhat/\Z$.

\bibitem{EnayatHamkinsWcislo}
A.~Enayat, J.~D.~Hamkins, and B.~Wcis\l o,
\emph{Topological models of arithmetic},
\url{https://arxiv.org/abs/1808.01270v3} (2020).
Used by source~02 only for background on internal arithmetic and
factorials; the residue obstruction of \cref{isg:cf:thm:PA} is proved
here.
\end{thebibliography}
\end{document}
